\documentclass[12pt,a4paper,openany]{book}
\usepackage[utf8]{inputenc}
\usepackage[T1]{fontenc}
\usepackage[french]{babel}
\usepackage[margin=2.5cm]{geometry}
\usepackage{amsmath,amssymb,amsthm,mathrsfs,mathtools}
\usepackage{tikz}
\usetikzlibrary{arrows.meta,calc,positioning,patterns,decorations.markings,decorations.pathmorphing}
\usepackage{pgfplots}
\pgfplotsset{compat=1.18}
\usepackage[most]{tcolorbox}
\tcbuselibrary{theorems,skins,breakable}
\usepackage[colorlinks=true,linkcolor=blue!60!black,citecolor=green!50!black]{hyperref}
\usepackage{makeidx}
\makeindex
\usepackage{enumitem}
\usepackage{fancyhdr}
\usepackage{caption}
\usepackage{booktabs}
\usepackage{array}
\usepackage{microtype}
\usepackage{xcolor}

\pagestyle{fancy}
\fancyhf{}
\fancyhead[LE]{\leftmark}
\fancyhead[RO]{\rightmark}
\fancyfoot[C]{\thepage}

% Theorem environments — MUST use \newtheorem (amsthm), NOT \newtcbtheorem
\theoremstyle{definition}
\newtheorem{definition}{Définition}[chapter]
\newtheorem{exemple}[definition]{Exemple}
\newtheorem{exercice}{Exercice}[chapter]

\theoremstyle{plain}
\newtheorem{theoreme}[definition]{Théorème}
\newtheorem{proposition}[definition]{Proposition}
\newtheorem{lemme}[definition]{Lemme}
\newtheorem{corollaire}[definition]{Corollaire}

\theoremstyle{remark}
\newtheorem{remarque}[definition]{Remarque}
\newtheorem{notation}[definition]{Notation}

% tcolorbox wrapping (visual only — do NOT use \newtcbtheorem)
\tcolorboxenvironment{definition}{enhanced,breakable,colback=blue!5,colframe=blue!70!black,fonttitle=\bfseries,before skip=10pt,after skip=10pt,left=8pt,right=8pt}
\tcolorboxenvironment{theoreme}{enhanced,breakable,colback=red!5,colframe=red!70!black,fonttitle=\bfseries,before skip=10pt,after skip=10pt,left=8pt,right=8pt}
\tcolorboxenvironment{proposition}{enhanced,breakable,colback=orange!5,colframe=orange!70!black,fonttitle=\bfseries,before skip=10pt,after skip=10pt,left=8pt,right=8pt}
\tcolorboxenvironment{lemme}{enhanced,breakable,colback=yellow!5,colframe=yellow!50!black,fonttitle=\bfseries,before skip=10pt,after skip=10pt,left=8pt,right=8pt}
\tcolorboxenvironment{corollaire}{enhanced,breakable,colback=green!5,colframe=green!50!black,fonttitle=\bfseries,before skip=10pt,after skip=10pt,left=8pt,right=8pt}
\tcolorboxenvironment{remarque}{enhanced,breakable,colback=gray!5,colframe=gray!50!black,before skip=8pt,after skip=8pt,left=6pt,right=6pt}
\tcolorboxenvironment{exemple}{enhanced,breakable,colback=purple!5,colframe=purple!50!black,before skip=8pt,after skip=8pt,left=6pt,right=6pt}

% Custom commands
\newcommand{\N}{\mathbb{N}}
\newcommand{\Z}{\mathbb{Z}}
\newcommand{\Q}{\mathbb{Q}}
\newcommand{\R}{\mathbb{R}}
\newcommand{\C}{\mathbb{C}}
\newcommand{\F}{\mathbb{F}}
\newcommand{\abs}[1]{\left|#1\right|}
\newcommand{\norm}[1]{\left\|#1\right\|}
\newcommand{\floor}[1]{\left\lfloor #1 \right\rfloor}
\newcommand{\ceil}[1]{\left\lceil #1 \right\rceil}
\newcommand{\legendre}[2]{\left(\frac{#1}{#2}\right)}
\newcommand{\jacobi}[2]{\left(\frac{#1}{#2}\right)}
\DeclareMathOperator{\pgcd}{pgcd}
\DeclareMathOperator{\ppcm}{ppcm}
\DeclareMathOperator{\ord}{ord}
\DeclareMathOperator{\Ker}{Ker}
\DeclareMathOperator{\im}{Im}

% ============================================================
\begin{document}

% ============================================================
%                       PAGE DE TITRE
% ============================================================
\begin{titlepage}
\begin{tikzpicture}[remember picture, overlay]
  \fill[left color=blue!15!white, right color=blue!5!white]
    (current page.south west) rectangle (current page.north east);
  \fill[color=orange!70!yellow]
    ([yshift=-2cm]current page.north west) rectangle ([yshift=-2.3cm]current page.north east);
  \fill[color=blue!80!black, rounded corners=8pt]
    ([xshift=3cm, yshift=-4cm]current page.north west) rectangle
    ([xshift=-3cm, yshift=-10cm]current page.north east);
  \node[anchor=center, text=white, font=\Huge\bfseries]
    at ([yshift=-6cm]current page.north) {Th\'eorie des Nombres};
  \node[anchor=center, text=orange!80!yellow, font=\Large]
    at ([yshift=-7.5cm]current page.north) {Notes de Cours};
  \node[anchor=center, text=white!80, font=\large]
    at ([yshift=-8.8cm]current page.north) {Master M1 --- 2025--2026};
  \node[anchor=center, text=blue!80!black, font=\Large\itshape]
    at ([yshift=-12cm]current page.north) {Ya\'e Ulrich Gaba};
  \draw[orange!70!yellow, line width=2pt]
    ([xshift=5cm, yshift=-13.5cm]current page.north west) --
    ([xshift=-5cm, yshift=-13.5cm]current page.north east);
  \node[anchor=center, text width=12cm, align=center, text=blue!60!black, font=\itshape]
    at ([yshift=-16cm]current page.north) {<<\,Les math\'ematiques sont la reine des sciences et la th\'eorie des nombres est la reine des math\'ematiques.\,>>\\[4pt]--- Carl Friedrich Gauss};
  \node[anchor=south, text=gray, font=\small]
    at ([yshift=1.5cm]current page.south) {\today};
\end{tikzpicture}
\end{titlepage}

% ============================================================
%                        PRÉFACE
% ============================================================
\chapter*{Préface}
\addcontentsline{toc}{chapter}{Préface}
\label{preface}

\index{preface@Préface}

La théorie des nombres, que Carl Friedrich Gauss\index{Gauss, Carl Friedrich} qualifiait
de \emph{reine des mathématiques} (\textit{die Königin der Mathematik}), est l'une
des branches les plus anciennes et les plus profondes des mathématiques. Ses origines
remontent à l'Antiquité, aux travaux d'Euclide\index{Euclide} dans ses
\textit{Éléments}, et ses questions fondamentales --- la nature des nombres premiers,
la résolution des équations diophantiennes, la distribution des entiers ayant des
propriétés arithmétiques particulières --- continuent de fasciner et de défier les
mathématiciens de notre époque.

Ce qui rend la théorie des nombres si singulière, c'est le contraste saisissant entre
la simplicité apparente de ses énoncés et la profondeur souvent abyssale de ses
démonstrations. La conjecture de Goldbach\index{conjecture!de Goldbach}, formulée en 1742,
peut être comprise par un collégien~: \emph{tout entier pair supérieur à $2$ est somme de
deux nombres premiers}. Pourtant, près de trois siècles plus tard, elle résiste
toujours aux efforts des plus brillants esprits.

L'histoire de la théorie des nombres est jalonnée de noms illustres~:
Euclide\index{Euclide}, Diophante\index{Diophante}, Fermat\index{Fermat, Pierre de},
Euler\index{Euler, Leonhard}, Gauss\index{Gauss, Carl Friedrich},
Riemann\index{Riemann, Bernhard}, Ramanujan\index{Ramanujan, Srinivasa}, et plus
récemment Andrew Wiles\index{Wiles, Andrew}, dont la démonstration du dernier théorème
de Fermat en 1995 a couronné trois siècles et demi de recherches. Chacun de ces
mathématiciens a apporté des outils nouveaux --- l'analyse complexe, la géométrie
algébrique, la théorie des représentations --- montrant que la théorie des nombres
se nourrit de toutes les branches des mathématiques et les enrichit en retour.

De surcroît, la théorie des nombres a acquis au cours des dernières décennies une
importance pratique considérable. La cryptographie à clé publique\index{cryptographie},
en particulier le système RSA\index{RSA}, repose de manière essentielle sur
l'arithmétique des grands nombres premiers et sur la difficulté du problème de la
factorisation. La théorie des codes correcteurs d'erreurs, les générateurs de nombres
pseudo-aléatoires et de nombreux algorithmes en informatique théorique puisent
directement dans l'arsenal de la théorie des nombres.

\medskip

\noindent\textbf{Organisation de ces notes.}
Ce cours est conçu pour des étudiants de niveau licence avancée ou master. Il suppose
une bonne maîtrise de l'algèbre élémentaire (groupes, anneaux) et de l'analyse de
base. Les chapitres sont organisés de manière progressive~:

\begin{itemize}[leftmargin=2cm]
  \item Les \textbf{Chapitres 1--2} posent les fondations~: divisibilité, algorithme
    d'Euclide, nombres premiers, théorème fondamental de l'arithmétique.
  \item Les \textbf{Chapitres 3--5} développent l'arithmétique modulaire, les
    résidus quadratiques et les fonctions arithmétiques.
  \item Les \textbf{Chapitres 6--8} abordent des sujets plus avancés~: fractions
    continues, approximation diophantienne, formes quadratiques.
  \item Les \textbf{Chapitres 9--10} offrent une introduction à la théorie analytique
    des nombres et aux méthodes du crible.
\end{itemize}

Chaque chapitre contient des démonstrations complètes, des exemples détaillés et
une série d'exercices classés par difficulté~:
\begin{itemize}[leftmargin=2cm]
  \item[$\star$] Exercices d'application directe.
  \item[$\star\star$] Exercices demandant une réflexion plus poussée.
  \item[$\star\star\star$] Problèmes de niveau olympiade ou concours.
\end{itemize}

\medskip

\noindent\textbf{Conventions.}
Sauf mention contraire, les lettres $a, b, c, d, m, n$ désignent des entiers. Le
mot \og nombre \fg{} sans qualificatif signifie \og entier \fg{}. Toutes les
démonstrations sont rédigées de manière complète~; le symbole $\qedsymbol$ marque
la fin d'une preuve.

\newpage

% ============================================================
%                   TABLE DES NOTATIONS
% ============================================================
\chapter*{Notations}
\addcontentsline{toc}{chapter}{Notations}
\label{sec:notations}

\index{notations}

Le tableau ci-dessous rassemble les notations utilisées tout au long de ce cours.

\bigskip

\begin{center}
\renewcommand{\arraystretch}{1.4}
\begin{tabular}{>{\raggedright}p{3.5cm} p{9cm}}
\toprule
\textbf{Notation} & \textbf{Signification} \\
\midrule
$\N$ & Ensemble des entiers naturels $\{0, 1, 2, \ldots\}$ \\
$\N^*$ & Ensemble des entiers naturels non nuls $\{1, 2, 3, \ldots\}$ \\
$\Z$ & Anneau des entiers relatifs \\
$\Q$ & Corps des nombres rationnels \\
$\R$ & Corps des nombres réels \\
$\C$ & Corps des nombres complexes \\
$\Z/n\Z$ & Anneau des entiers modulo $n$ \\
$(\Z/n\Z)^\times$ & Groupe des éléments inversibles de $\Z/n\Z$ \\
$\F_p$ & Corps fini à $p$ éléments ($p$ premier) \\
$a \mid b$ & $a$ divise $b$ \\
$a \nmid b$ & $a$ ne divise pas $b$ \\
$\pgcd(a,b)$ & Plus grand commun diviseur de $a$ et $b$ \\
$\ppcm(a,b)$ & Plus petit commun multiple de $a$ et $b$ \\
$a \equiv b \pmod{n}$ & $a$ est congru à $b$ modulo $n$ \\
$\floor{x}$ & Partie entière (plancher) de $x$ \\
$\ceil{x}$ & Partie entière supérieure (plafond) de $x$ \\
$\{x\}$ & Partie fractionnaire de $x$~: $\{x\} = x - \floor{x}$ \\
$v_p(n)$ & Valuation $p$-adique de $n$ \\
$\varphi(n)$ & Indicatrice d'Euler\index{indicatrice d'Euler} \\
$\mu(n)$ & Fonction de Möbius\index{fonction de Möbius} \\
$\tau(n)$, $d(n)$ & Nombre de diviseurs de $n$ \\
$\sigma(n)$ & Somme des diviseurs de $n$ \\
$\sigma_k(n)$ & Somme des puissances $k$-ièmes des diviseurs de $n$ \\
$\Lambda(n)$ & Fonction de von Mangoldt\index{fonction de von Mangoldt} \\
$\pi(x)$ & Nombre de premiers $\leq x$ \\
$\legendre{a}{p}$ & Symbole de Legendre\index{symbole de Legendre} \\
$\jacobi{a}{n}$ & Symbole de Jacobi\index{symbole de Jacobi} \\
$\ord_n(a)$ & Ordre de $a$ modulo $n$ \\
$f \ast g$ & Convolution de Dirichlet\index{convolution de Dirichlet} de $f$ et $g$ \\
$f = O(g)$ & $\abs{f(x)} \leq C\abs{g(x)}$ pour $x$ assez grand \\
$f \sim g$ & $f(x)/g(x) \to 1$ quand $x \to +\infty$ \\
\bottomrule
\end{tabular}
\end{center}

\newpage

% ============================================================
%                  TABLE DES MATIÈRES
% ============================================================
\tableofcontents

\newpage

% ============================================================
% ============================================================
%           CHAPITRE 1 : DIVISIBILITÉ ET ALGORITHME D'EUCLIDE
% ============================================================
% ============================================================

\chapter{Divisibilité et Algorithme d'Euclide}
\label{ch:divisibilite}

\index{divisibilite@divisibilité}

\begin{quote}
\itshape
\og Les mathématiques sont la reine des sciences, et la théorie des nombres est la
reine des mathématiques. \fg{}

\hfill --- Carl Friedrich Gauss
\end{quote}

\bigskip

\section{Introduction historique}
\label{sec:ch1:intro}

\index{Euclide!Éléments}

L'étude de la divisibilité est aussi ancienne que les mathématiques elles-mêmes. Dans
le Livre~VII de ses \textit{Éléments}\index{Elements@\textit{Éléments} (Euclide)},
rédigés vers 300 av.~J.-C., Euclide pose les fondations de l'arithmétique en
définissant la notion de diviseur, de nombre premier, et en établissant un algorithme
pour calculer la \og plus grande commune mesure \fg{} de deux entiers --- ce que nous
appelons aujourd'hui le \emph{plus grand commun diviseur}\index{pgcd}.

L'algorithme d'Euclide\index{algorithme d'Euclide} est remarquable à plus d'un titre~:
c'est l'un des plus anciens algorithmes connus, et il reste l'un des plus efficaces
pour le calcul du pgcd, même à l'ère des ordinateurs. Sa variante
\emph{étendue}\index{algorithme d'Euclide!étendu}, qui fournit les coefficients de
Bézout\index{Bezout@Bézout, identité de}, est au c\oe ur de la cryptographie
moderne, notamment du système RSA\index{RSA}.

Ce chapitre développe de manière rigoureuse les notions fondamentales de divisibilité
dans $\Z$, l'algorithme d'Euclide et ses applications.

% ============================================================
\section{Divisibilité dans $\Z$}
\label{sec:ch1:divisibilite}

\index{divisibilite@divisibilité!définition}

\begin{definition}[Divisibilité]
\label{def:divisibilite}
Soient $a, b \in \Z$. On dit que $a$ \emph{divise} $b$, et l'on note $a \mid b$,
s'il existe $k \in \Z$ tel que $b = ka$. On dit aussi que $a$ est un
\emph{diviseur}\index{diviseur} de $b$, ou que $b$ est un
\emph{multiple}\index{multiple} de $a$.

Si $a$ ne divise pas $b$, on note $a \nmid b$.
\end{definition}

\begin{exemple}[Exemples fondamentaux]
\label{ex:div:fondamentaux}
\begin{enumerate}[label=(\roman*)]
  \item $3 \mid 12$ car $12 = 4 \times 3$.
  \item $7 \mid (-21)$ car $-21 = (-3) \times 7$.
  \item $1 \mid n$ pour tout $n \in \Z$ car $n = n \times 1$.
  \item $n \mid 0$ pour tout $n \in \Z$ car $0 = 0 \times n$.
  \item $0 \mid n$ si et seulement si $n = 0$.
  \item $5 \nmid 13$ car il n'existe aucun $k \in \Z$ tel que $13 = 5k$.
\end{enumerate}
\end{exemple}

\begin{proposition}[Propriétés élémentaires de la divisibilité]
\label{prop:div:elem}
\index{divisibilite@divisibilité!propriétés}
Soient $a, b, c \in \Z$. Les propriétés suivantes sont vérifiées~:
\begin{enumerate}[label=(\roman*)]
  \item \textbf{Réflexivité}~: $a \mid a$.
  \item \textbf{Transitivité}~: si $a \mid b$ et $b \mid c$, alors $a \mid c$.
  \item \textbf{Combinaison linéaire}~: si $a \mid b$ et $a \mid c$, alors
    $a \mid (ub + vc)$ pour tous $u, v \in \Z$.
  \item \textbf{Antisymétrie (à signe près)}~: si $a \mid b$ et $b \mid a$, alors
    $a = \pm b$.
  \item \textbf{Compatibilité avec le produit}~: si $a \mid b$, alors $a \mid bc$
    pour tout $c \in \Z$.
  \item \textbf{Simplification}~: si $c \neq 0$ et $ca \mid cb$, alors $a \mid b$.
\end{enumerate}
\end{proposition}

\begin{proof}
\begin{enumerate}[label=(\roman*)]
  \item On a $a = 1 \cdot a$, donc $a \mid a$.

  \item Supposons $a \mid b$ et $b \mid c$. Il existe $k_1, k_2 \in \Z$ tels que
    $b = k_1 a$ et $c = k_2 b$. Alors $c = k_2 (k_1 a) = (k_1 k_2) a$, d'où
    $a \mid c$.

  \item Si $a \mid b$ et $a \mid c$, il existe $k_1, k_2 \in \Z$ tels que
    $b = k_1 a$ et $c = k_2 a$. Pour tous $u, v \in \Z$~:
    \[
      ub + vc = u(k_1 a) + v(k_2 a) = (uk_1 + vk_2)\, a,
    \]
    donc $a \mid (ub + vc)$.

  \item Si $a \mid b$ et $b \mid a$, il existe $k_1, k_2 \in \Z$ tels que
    $b = k_1 a$ et $a = k_2 b$. Si $a = 0$, alors $b = 0$ et $a = b$. Sinon,
    $a = k_2 k_1 a$, d'où $k_1 k_2 = 1$. Comme $k_1, k_2 \in \Z$, on a
    $k_1 = k_2 = 1$ ou $k_1 = k_2 = -1$, c'est-à-dire $a = \pm b$.

  \item Si $b = ka$, alors $bc = (kc)a$, donc $a \mid bc$.

  \item Si $ca \mid cb$, il existe $k \in \Z$ tel que $cb = k \cdot ca$. Comme
    $c \neq 0$, on simplifie par $c$ pour obtenir $b = ka$, soit $a \mid b$.
    \qedhere
\end{enumerate}
\end{proof}

\begin{remarque}
\label{rem:div:ordre}
La relation de divisibilité est une relation d'ordre partiel sur $\N^*$~: elle est
réflexive, antisymétrique et transitive. Elle n'est pas totale~: par exemple,
$3 \nmid 5$ et $5 \nmid 3$.
\end{remarque}

% ============================================================
\section{Division euclidienne}
\label{sec:ch1:div_eucl}

\index{division euclidienne}

\begin{theoreme}[Division euclidienne]
\label{thm:div_eucl}
Soient $a \in \Z$ et $b \in \Z \setminus \{0\}$. Il existe un unique couple
$(q, r) \in \Z \times \N$ tel que
\[
  a = bq + r \quad \text{et} \quad 0 \leq r < \abs{b}.
\]
L'entier $q$ est appelé le \emph{quotient}\index{quotient} et $r$ le
\emph{reste}\index{reste} de la division euclidienne de $a$ par $b$.
\end{theoreme}

\begin{proof}
\textbf{Existence.}
Considérons l'ensemble
\[
  S = \{a - bk : k \in \Z\} \cap \N = \{a - bk : k \in \Z,\ a - bk \geq 0\}.
\]
Montrons que $S$ est non vide. Si $b > 0$, prenons $k = -\abs{a}$~: alors
$a - bk = a + b\abs{a} \geq a + \abs{a} \geq 0$. Si $b < 0$, prenons
$k = \abs{a}$~: alors $a - bk = a + \abs{b}\abs{a} \geq 0$. Dans les deux cas,
$S \neq \emptyset$.

Comme $S \subseteq \N$, $S$ admet un plus petit élément par le principe du bon
ordre\index{principe du bon ordre}. Notons $r = \min S$ et $q$ l'entier correspondant,
de sorte que $r = a - bq \geq 0$.

Montrons que $r < \abs{b}$. Par l'absurde, supposons $r \geq \abs{b}$. Alors
\[
  a - b(q + \operatorname{sgn}(b)) = a - bq - \abs{b} = r - \abs{b} \geq 0,
\]
où $\operatorname{sgn}(b) = b/\abs{b}$. Donc $r - \abs{b} \in S$ avec
$r - \abs{b} < r$, ce qui contredit la minimalité de $r$. Par conséquent
$0 \leq r < \abs{b}$.

\medskip
\textbf{Unicité.}
Supposons qu'il existe deux couples $(q_1, r_1)$ et $(q_2, r_2)$ vérifiant les
conditions. Alors
\[
  bq_1 + r_1 = bq_2 + r_2 \implies b(q_1 - q_2) = r_2 - r_1.
\]
Or $0 \leq r_1, r_2 < \abs{b}$, d'où $\abs{r_2 - r_1} < \abs{b}$.
Comme $b(q_1 - q_2) = r_2 - r_1$ et $\abs{b(q_1-q_2)} = \abs{b}\abs{q_1-q_2}$,
on doit avoir $\abs{b}\abs{q_1-q_2} < \abs{b}$, ce qui impose
$\abs{q_1-q_2} < 1$, soit $q_1 = q_2$. Il s'ensuit que $r_1 = r_2$.
\end{proof}

\begin{exemple}[Calculs de division euclidienne]
\label{ex:div_eucl}
\begin{enumerate}[label=(\roman*)]
  \item Division de $a = 157$ par $b = 13$~:
    $157 = 13 \times 12 + 1$, donc $q = 12$ et $r = 1$.
  \item Division de $a = -157$ par $b = 13$~:
    $-157 = 13 \times (-13) + 12$, donc $q = -13$ et $r = 12$.
  \item Division de $a = 100$ par $b = -7$~:
    $100 = (-7)\times(-14) + 2$, donc $q = -14$ et $r = 2$.
\end{enumerate}
\end{exemple}

\begin{remarque}
\label{rem:div_eucl:divisible}
Le reste $r$ est nul si et seulement si $b \mid a$. La division euclidienne
généralise donc la notion de divisibilité~: $b \mid a$ si et seulement si
le reste de la division de $a$ par $b$ est $0$.
\end{remarque}

% ============================================================
\section{Plus grand commun diviseur}
\label{sec:ch1:pgcd}

\index{pgcd!définition}

\begin{definition}[Plus grand commun diviseur]
\label{def:pgcd}
Soient $a, b \in \Z$, non tous les deux nuls. Le \emph{plus grand commun
diviseur}\index{pgcd} de $a$ et $b$, noté $\pgcd(a,b)$, est le plus grand entier
$d \in \N^*$ tel que $d \mid a$ et $d \mid b$.
\end{definition}

\begin{proposition}[Existence et caractérisation du pgcd]
\label{prop:pgcd:caract}
Soient $a, b \in \Z$, non tous les deux nuls. Le $\pgcd(a,b)$ existe et est l'unique
entier $d \in \N^*$ vérifiant~:
\begin{enumerate}[label=(\roman*)]
  \item $d \mid a$ et $d \mid b$~;
  \item pour tout $c \in \Z$, si $c \mid a$ et $c \mid b$, alors $c \mid d$.
\end{enumerate}
\end{proposition}

\begin{proof}
L'ensemble des diviseurs communs de $a$ et $b$ est non vide (il contient $1$) et
majoré par $\max(\abs{a}, \abs{b})$ (tout diviseur de $a$ est $\leq \abs{a}$ en
valeur absolue, et de même pour $b$, sauf si $a = 0$ ou $b = 0$, auquel cas le
pgcd est $\abs{b}$ ou $\abs{a}$ respectivement). Il admet donc un plus grand
élément $d$.

Pour la propriété (ii), nous la démontrerons comme conséquence de l'identité de
Bézout (cf.\ 
ef{thm:bezout}).
\end{proof}

\begin{proposition}[Propriétés du pgcd]
\label{prop:pgcd:props}
\index{pgcd!propriétés}
Pour tous $a, b, c \in \Z$~:
\begin{enumerate}[label=(\roman*)]
  \item $\pgcd(a,b) = \pgcd(b,a)$ \textup{(commutativité)}.
  \item $\pgcd(a,b) = \pgcd(\abs{a}, \abs{b})$.
  \item $\pgcd(a, 0) = \abs{a}$ pour $a \neq 0$.
  \item $\pgcd(a, 1) = 1$.
  \item $\pgcd(a,b) = \pgcd(a, b - qa)$ pour tout $q \in \Z$.
  \item $\pgcd(ca, cb) = \abs{c}\, \pgcd(a,b)$ pour $c \neq 0$.
\end{enumerate}
\end{proposition}

\begin{proof}
Les propriétés (i)--(iv) découlent immédiatement de la définition.

(v) Il suffit de montrer que $a$ et $b$ ont les mêmes diviseurs communs que $a$ et
$b - qa$. Si $d \mid a$ et $d \mid b$, alors $d \mid (b - qa)$ par linéarité.
Réciproquement, si $d \mid a$ et $d \mid (b-qa)$, alors $d \mid ((b-qa) + qa) = b$.

(vi) Si $d = \pgcd(a,b)$, alors $d \mid a$ et $d \mid b$, donc $\abs{c}\,d \mid ca$
et $\abs{c}\,d \mid cb$. Réciproquement, tout diviseur commun de $ca$ et $cb$ est de
la forme $\abs{c}\,e$ où $e$ est un diviseur commun de $a$ et $b$, d'où le résultat.
\end{proof}

La propriété~(v) est la clé de l'algorithme d'Euclide.

% ============================================================
\section{Algorithme d'Euclide}
\label{sec:ch1:algo_euclide}

\index{algorithme d'Euclide}

L'algorithme d'Euclide calcule le pgcd de deux entiers par une suite de divisions
euclidiennes. Il repose sur la propriété fondamentale
$\pgcd(a,b) = \pgcd(b, a \bmod b)$.

\begin{theoreme}[Algorithme d'Euclide --- Correction]
\label{thm:algo_euclide}
Soient $a, b \in \N$ avec $b > 0$. La suite de divisions euclidiennes~:
\begin{align*}
  a   &= bq_1 + r_1, & 0 &\leq r_1 < b, \\
  b   &= r_1 q_2 + r_2, & 0 &\leq r_2 < r_1, \\
  r_1 &= r_2 q_3 + r_3, & 0 &\leq r_3 < r_2, \\
      &\;\;\vdots \\
  r_{n-2} &= r_{n-1} q_n + r_n, & 0 &\leq r_n < r_{n-1}, \\
  r_{n-1} &= r_n q_{n+1} + 0,
\end{align*}
se termine en un nombre fini d'étapes, et le dernier reste non nul $r_n$ est
$\pgcd(a,b)$.
\end{theoreme}

\begin{proof}
\textbf{Terminaison.} La suite $b > r_1 > r_2 > \cdots \geq 0$ est une suite
strictement décroissante d'entiers naturels. Par le principe du bon
ordre\index{principe du bon ordre}, elle atteint $0$ en un nombre fini d'étapes.

\textbf{Correction.} Par la propriété~(v) de la 
ef{prop:pgcd:props},
$\pgcd(a,b) = \pgcd(b, r_1) = \pgcd(r_1, r_2) = \cdots = \pgcd(r_{n-1}, r_n)
= \pgcd(r_n, 0) = r_n$.
\end{proof}

\begin{exemple}[Calcul de $\pgcd(252, 198)$]
\label{ex:euclide:252_198}
\index{algorithme d'Euclide!exemple}
Appliquons l'algorithme d'Euclide~:
\begin{align*}
  252 &= 198 \times 1 + 54, \\
  198 &= 54 \times 3 + 36, \\
  54  &= 36 \times 1 + 18, \\
  36  &= 18 \times 2 + 0.
\end{align*}
Le dernier reste non nul est $18$, donc $\pgcd(252, 198) = 18$.
\end{exemple}

\begin{exemple}[Calcul de $\pgcd(1071, 462)$]
\label{ex:euclide:1071_462}
\begin{align*}
  1071 &= 462 \times 2 + 147, \\
  462  &= 147 \times 3 + 21, \\
  147  &= 21 \times 7 + 0.
\end{align*}
Le dernier reste non nul est $21$, donc $\pgcd(1071, 462) = 21$.
\end{exemple}

\begin{proposition}[Complexité de l'algorithme d'Euclide]
\label{prop:euclide:complexite}
\index{algorithme d'Euclide!complexité}
L'algorithme d'Euclide appliqué à deux entiers $a \geq b > 0$ effectue au plus
$\floor{2\log_2(\max(a,b))} + 1$ divisions. Plus précisément, si l'algorithme
effectue $n$ divisions, alors $b \geq F_{n}$ et $a \geq F_{n+1}$, où $(F_k)$ est
la suite de Fibonacci\index{Fibonacci, suite de}.
\end{proposition}

\begin{proof}
Montrons que les restes décroissent au moins d'un facteur $2$ tous les deux pas.
Si $r_{i+1} \leq r_i / 2$, c'est immédiat. Si $r_{i+1} > r_i / 2$, alors
$q_{i+2} = 1$ et $r_{i+2} = r_i - r_{i+1} < r_i / 2$. Donc en deux étapes, le
reste est au moins divisé par $2$, ce qui donne au plus $2\log_2(b)$ étapes.

Pour la borne de Fibonacci, on montre par récurrence descendante que
$r_{n-k} \geq F_{k+2}$ pour $0 \leq k \leq n$, en utilisant $r_{n-k-2} \geq
r_{n-k-1} + r_{n-k} \geq F_{k+2} + F_{k+3} = F_{k+4}$, ce qui est la récurrence
de Fibonacci. En remontant, on obtient $b \geq F_{n}$ et $a \geq F_{n+1}$.
\end{proof}

\begin{remarque}
\label{rem:euclide:lamé}
Ce résultat est essentiellement le \textbf{théorème de Lamé}\index{Lamé, théorème de}
(1844)~: le nombre de divisions dans l'algorithme d'Euclide appliqué à deux entiers
inférieurs à $N$ est au plus $\floor{\log_\phi(N\sqrt{5})}$, où $\phi =
\frac{1+\sqrt{5}}{2}$ est le nombre d'or. Les pires cas sont atteints pour des
nombres de Fibonacci consécutifs.
\end{remarque}

% ============================================================
\section{Identité de Bézout}
\label{sec:ch1:bezout}

\index{Bezout@Bézout, identité de}

\begin{theoreme}[Identité de Bézout]
\label{thm:bezout}
Soient $a, b \in \Z$, non tous les deux nuls. Il existe $u, v \in \Z$ tels que
\[
  au + bv = \pgcd(a,b).
\]
De plus, l'ensemble $\{au + bv : u, v \in \Z\}$ est exactement l'ensemble des
multiples de $\pgcd(a,b)$, c'est-à-dire $a\Z + b\Z = \pgcd(a,b)\,\Z$.
\end{theoreme}

\begin{proof}
Considérons l'ensemble
\[
  I = \{au + bv : u, v \in \Z\} \cap \N^*.
\]
Cet ensemble est non vide car $\abs{a} = a \cdot \operatorname{sgn}(a) + b \cdot 0
\in I$ (en supposant $a \neq 0$~; sinon $\abs{b} \in I$). Soit $d = \min I$ le
plus petit élément de $I$ (par le principe du bon ordre). Par construction,
$d = au_0 + bv_0$ pour certains $u_0, v_0 \in \Z$.

\textbf{Montrons que $d \mid a$.} Effectuons la division euclidienne de $a$ par $d$~:
$a = dq + r$ avec $0 \leq r < d$. Alors
\[
  r = a - dq = a - (au_0 + bv_0)q = a(1 - u_0 q) + b(-v_0 q).
\]
Donc $r \in a\Z + b\Z$. Si $r > 0$, alors $r \in I$ et $r < d$, ce qui contredit
la minimalité de $d$. Donc $r = 0$ et $d \mid a$.

Par le même argument, $d \mid b$. Donc $d$ est un diviseur commun de $a$ et $b$.

\textbf{Montrons que $d = \pgcd(a,b)$.} Si $c$ est un diviseur commun quelconque de
$a$ et $b$, alors $c \mid (au_0 + bv_0) = d$. Donc $c \leq d$ (en valeur absolue),
ce qui montre que $d$ est bien le plus grand diviseur commun.

\textbf{Caractérisation de $I \cup \{0\}$.} Si $m \in a\Z + b\Z$, effectuons la
division euclidienne de $m$ par $d$~: $m = dq + r$ avec $0 \leq r < d$. Le même
argument montre $r = 0$, donc $d \mid m$, i.e.\ $m \in d\Z$. Réciproquement,
$d = au_0 + bv_0$, donc $dk = a(ku_0) + b(kv_0) \in a\Z + b\Z$. Ainsi
$a\Z + b\Z = d\Z$.
\end{proof}

\begin{corollaire}
\label{cor:bezout:caract_pgcd}
\index{pgcd!caractérisation par Bézout}
Le $\pgcd(a,b)$ est le plus petit élément strictement positif de l'ensemble
$\{au + bv : u, v \in \Z\}$.
\end{corollaire}

\subsection{Algorithme d'Euclide étendu}
\label{subsec:ch1:euclide_etendu}

\index{algorithme d'Euclide!étendu}

L'algorithme d'Euclide étendu permet de calculer simultanément le $\pgcd(a,b)$ et les
coefficients de Bézout $u, v$ tels que $au + bv = \pgcd(a,b)$. L'idée est de
remonter les étapes de l'algorithme d'Euclide.

\begin{exemple}[Coefficients de Bézout pour $\pgcd(252, 198)$]
\label{ex:bezout:252_198}
\index{Bezout@Bézout!exemple}
Reprenons l'exemple~\ref{ex:euclide:252_198}. On a trouvé~:
\begin{align*}
  252 &= 198 \times 1 + 54, \\
  198 &= 54 \times 3 + 36, \\
  54  &= 36 \times 1 + 18.
\end{align*}
Remontons les équations~:
\begin{align*}
  18 &= 54 - 36 \times 1 \\
     &= 54 - (198 - 54 \times 3) \times 1 \\
     &= 54 \times 4 - 198 \times 1 \\
     &= (252 - 198 \times 1) \times 4 - 198 \times 1 \\
     &= 252 \times 4 - 198 \times 5.
\end{align*}
Donc $u = 4$, $v = -5$ et $252 \times 4 + 198 \times (-5) = 1008 - 990 = 18$.
\end{exemple}

\begin{exemple}[Coefficients de Bézout pour $\pgcd(1071, 462)$]
\label{ex:bezout:1071_462}
Reprenons l'exemple~\ref{ex:euclide:1071_462}~:
\begin{align*}
  1071 &= 462 \times 2 + 147, \\
  462  &= 147 \times 3 + 21.
\end{align*}
Remontée~:
\begin{align*}
  21 &= 462 - 147 \times 3 \\
     &= 462 - (1071 - 462 \times 2) \times 3 \\
     &= 462 \times 7 - 1071 \times 3.
\end{align*}
Donc $1071 \times (-3) + 462 \times 7 = -3213 + 3234 = 21$. On a $u = -3$, $v = 7$.
\end{exemple}

\begin{remarque}[Implémentation tabulaire]
\label{rem:bezout:tableau}
En pratique, on peut calculer les coefficients de Bézout par un tableau. On maintient
à chaque étape les triplets $(r_i, u_i, v_i)$ tels que $r_i = au_i + bv_i$~:
\begin{center}
\begin{tabular}{cccc}
\toprule
Étape & $r_i$ & $u_i$ & $v_i$ \\
\midrule
$0$ & $a$ & $1$ & $0$ \\
$1$ & $b$ & $0$ & $1$ \\
$i+1$ & $r_{i-1} - q_i r_i$ & $u_{i-1} - q_i u_i$ & $v_{i-1} - q_i v_i$ \\
\bottomrule
\end{tabular}
\end{center}
Lorsque $r_{i+1} = 0$, on a $\pgcd(a,b) = r_i$ et les coefficients de Bézout sont
$(u_i, v_i)$.
\end{remarque}

% ============================================================
\section{Plus petit commun multiple}
\label{sec:ch1:ppcm}

\index{ppcm}

\begin{definition}[Plus petit commun multiple]
\label{def:ppcm}
Soient $a, b \in \Z \setminus \{0\}$. Le \emph{plus petit commun multiple} de $a$
et $b$, noté $\ppcm(a,b)$, est le plus petit entier $m \in \N^*$ tel que $a \mid m$
et $b \mid m$.
\end{definition}

\begin{theoreme}[Relation pgcd--ppcm]
\label{thm:pgcd_ppcm}
\index{pgcd!relation avec le ppcm}
\index{ppcm!relation avec le pgcd}
Pour tous $a, b \in \Z \setminus \{0\}$,
\[
  \pgcd(a,b) \cdot \ppcm(a,b) = \abs{ab}.
\]
\end{theoreme}

\begin{proof}
Posons $d = \pgcd(a,b)$ et écrivons $a = da'$, $b = db'$ avec $\pgcd(a',b') = 1$.
Montrons que $\ppcm(a,b) = da'b' = \abs{ab}/d$.

D'abord, $a \mid da'b'$ car $da'b' = ab'$, et $b \mid da'b'$ car $da'b' = a'b$.
Donc $da'b'$ est un multiple commun.

Soit $m$ un multiple commun quelconque~: $m = ak = bl$ pour certains $k, l \in \Z$.
Alors $da'k = db'l$, soit $a'k = b'l$. Comme $\pgcd(a',b') = 1$ et $a' \mid b'l$,
le lemme de Gauss (
ef{lemme:gauss}) donne $a' \mid l$. Donc
$m = bl = db'l \geq db'a' = da'b'$ (en valeur absolue). Ainsi $da'b'$ est le plus
petit multiple commun positif.

Enfin, $d \cdot da'b' = d^2 a'b' = (da')(db') = \abs{a}\abs{b} = \abs{ab}$.
\end{proof}

\begin{exemple}
\label{ex:pgcd_ppcm}
$\pgcd(12, 18) = 6$ et $\ppcm(12, 18) = 36$. Vérifions~:
$6 \times 36 = 216 = 12 \times 18 = \abs{12 \times 18}$. \checkmark
\end{exemple}

% ============================================================
\section{Entiers premiers entre eux et lemme de Gauss}
\label{sec:ch1:copremiers}

\index{entiers premiers entre eux}

\begin{definition}[Entiers premiers entre eux]
\label{def:copremiers}
Deux entiers $a$ et $b$ sont dits \emph{premiers entre eux}\index{premiers entre eux}
(ou \emph{copremiers}) si $\pgcd(a,b) = 1$.
\end{definition}

\begin{proposition}[Caractérisation par Bézout]
\label{prop:copremiers:bezout}
Deux entiers $a$ et $b$ sont premiers entre eux si et seulement s'il existe
$u, v \in \Z$ tels que $au + bv = 1$.
\end{proposition}

\begin{proof}
Si $\pgcd(a,b) = 1$, l'identité de Bézout fournit directement de tels $u, v$.
Réciproquement, si $au + bv = 1$ et si $d \mid a$ et $d \mid b$, alors
$d \mid (au + bv) = 1$, donc $d = \pm 1$, et $\pgcd(a,b) = 1$.
\end{proof}

\begin{lemme}[Lemme de Gauss]
\label{lemme:gauss}
\index{Gauss!lemme de}
Soient $a, b, c \in \Z$. Si $a \mid bc$ et $\pgcd(a,b) = 1$, alors $a \mid c$.
\end{lemme}

\begin{proof}
Puisque $\pgcd(a,b) = 1$, l'identité de Bézout fournit $u, v \in \Z$ tels que
$au + bv = 1$. En multipliant par $c$~:
\[
  acu + bcv = c.
\]
Or $a \mid acu$ (évidemment) et $a \mid bcv$ (car $a \mid bc$ par hypothèse). Donc
$a \mid (acu + bcv) = c$.
\end{proof}

\begin{corollaire}
\label{cor:gauss:produit_copremiers}
Si $a \mid c$, $b \mid c$ et $\pgcd(a,b) = 1$, alors $ab \mid c$.
\end{corollaire}

\begin{proof}
De $a \mid c$, on tire $c = ak$ pour un certain $k \in \Z$. Comme $b \mid c = ak$
et $\pgcd(a,b) = 1$, le lemme de Gauss donne $b \mid k$, soit $k = bl$. Donc
$c = abl$ et $ab \mid c$.
\end{proof}

\begin{corollaire}
\label{cor:gauss:premier_divise_produit}
Si $p$ est premier et $p \mid ab$, alors $p \mid a$ ou $p \mid b$.
\end{corollaire}

\begin{proof}
Si $p \nmid a$, alors $\pgcd(p,a) = 1$ (car les seuls diviseurs positifs de $p$ sont
$1$ et $p$, et $p \nmid a$). Par le lemme de Gauss, $p \mid b$.
\end{proof}

\begin{remarque}
\label{rem:gauss:generalisation}
Le 
ef{cor:gauss:premier_divise_produit} se généralise par récurrence~: si $p$ est
premier et $p \mid a_1 a_2 \cdots a_n$, alors $p \mid a_i$ pour au moins un indice
$i$. Cette propriété est cruciale pour la démonstration de l'unicité dans le théorème
fondamental de l'arithmétique (
ef{thm:fta}).
\end{remarque}

% ============================================================
\section{Diagramme du treillis de divisibilité}
\label{sec:ch1:tikz}

Le diagramme ci-dessous illustre le treillis des diviseurs\index{treillis de
divisibilité} de $30$~: deux n\oe uds sont reliés si l'un divise l'autre (avec
le plus petit en bas).

\begin{center}
\begin{tikzpicture}[
    node distance=1.5cm and 1.5cm,
    every node/.style={draw, circle, minimum size=8mm, inner sep=0pt,
      fill=blue!10, font=\small},
    >=Stealth
  ]
  % Niveau 0
  \node (1) {$1$};
  % Niveau 1
  \node (2) [above left=of 1, xshift=-0.5cm] {$2$};
  \node (3) [above=of 1] {$3$};
  \node (5) [above right=of 1, xshift=0.5cm] {$5$};
  % Niveau 2
  \node (6) [above right=of 2, xshift=-0.5cm] {$6$};
  \node (10) [above=of 5, xshift=-0.5cm] {$10$};
  \node (15) [above right=of 5, xshift=-0.5cm] {$15$};
  % Niveau 3
  \node (30) [above right=of 6, xshift=0.5cm] {$30$};

  % Arêtes
  \draw (1) -- (2);
  \draw (1) -- (3);
  \draw (1) -- (5);
  \draw (2) -- (6);
  \draw (2) -- (10);
  \draw (3) -- (6);
  \draw (3) -- (15);
  \draw (5) -- (10);
  \draw (5) -- (15);
  \draw (6) -- (30);
  \draw (10) -- (30);
  \draw (15) -- (30);
\end{tikzpicture}
\captionof{figure}{Treillis des diviseurs de $30$.}
\label{fig:treillis_30}
\end{center}

\index{treillis de divisibilité!de 30}

Dans ce treillis, le pgcd de deux éléments est leur borne inférieure et le ppcm est
leur borne supérieure. Par exemple, $\pgcd(6, 10) = 2$ (la borne inférieure de $6$
et $10$) et $\ppcm(6, 10) = 30$ (la borne supérieure).

% ============================================================
\section{Connexion avec la cryptographie~: le système RSA}
\label{sec:ch1:rsa}

\index{RSA!rôle du pgcd}
\index{cryptographie}

Le calcul efficace du pgcd, via l'algorithme d'Euclide, est une brique fondamentale
de la cryptographie à clé publique\index{cryptographie!clé publique}. Dans le système
RSA\index{RSA}, inventé par Rivest, Shamir et Adleman\index{Rivest, Ron}%
\index{Shamir, Adi}\index{Adleman, Leonard} en 1977~:

\begin{enumerate}[label=\arabic*.]
  \item On choisit deux grands nombres premiers $p$ et $q$ et on pose $n = pq$.
  \item On choisit un exposant de chiffrement $e$ tel que
    $\pgcd(e, \varphi(n)) = 1$, où $\varphi(n) = (p-1)(q-1)$ est l'indicatrice
    d'Euler.
  \item L'exposant de déchiffrement $d$ est l'inverse de $e$ modulo $\varphi(n)$,
    calculé par l'algorithme d'Euclide étendu~: on trouve $u, v$ tels que
    $eu + \varphi(n)v = 1$, et on prend $d = u \bmod \varphi(n)$.
\end{enumerate}

L'algorithme d'Euclide étendu intervient donc de manière cruciale dans la génération
des clés RSA. Sa rapidité (complexité en $O(\log^2 n)$ opérations sur les bits) est
essentielle pour la praticabilité du système.

% ============================================================
\section{Exercices}
\label{sec:ch1:exercices}

\begin{exercice}[$\star$]
\label{exo:ch1:pgcd_calculs}
Calculer, à l'aide de l'algorithme d'Euclide~:
\begin{enumerate}[label=(\alph*)]
  \item $\pgcd(420, 378)$,
  \item $\pgcd(1001, 770)$,
  \item $\pgcd(10403, 3731)$.
\end{enumerate}
Pour chaque cas, trouver les coefficients de Bézout.
\end{exercice}

\begin{exercice}[$\star$]
\label{exo:ch1:ppcm_calculs}
Calculer $\ppcm(120, 252)$ de deux manières~: directement et via la relation
$\pgcd \times \ppcm = \abs{ab}$.
\end{exercice}

\begin{exercice}[$\star$]
\label{exo:ch1:copremiers}
Montrer que pour tout $n \in \N^*$, les entiers $n$ et $n+1$ sont premiers entre
eux.
\end{exercice}

\begin{exercice}[$\star\star$]
\label{exo:ch1:bezout_equation}
Résoudre dans $\Z^2$ l'équation $252x + 198y = 36$.
\end{exercice}

\begin{exercice}[$\star\star$]
\label{exo:ch1:div_combinaison}
Soient $a, b \in \Z$ tels que $\pgcd(a,b) = 1$. Montrer que
$\pgcd(a+b, a-b) \in \{1, 2\}$.
\end{exercice}

\begin{exercice}[$\star\star$]
\label{exo:ch1:pgcd_fibo}
Montrer que pour tout $n \geq 1$, $\pgcd(F_n, F_{n+1}) = 1$, où $(F_n)$ est la
suite de Fibonacci\index{Fibonacci, suite de} définie par $F_1 = F_2 = 1$ et
$F_{n+2} = F_{n+1} + F_n$.
\end{exercice}

\begin{exercice}[$\star\star$]
\label{exo:ch1:pgcd_consecutifs}
Soient $a, b \in \N^*$. Montrer que $\pgcd(a^n - 1, a^m - 1) = a^{\pgcd(n,m)} - 1$.
\end{exercice}

\begin{exercice}[$\star\star$]
\label{exo:ch1:gauss_application}
Soient $a, b, c \in \Z$ avec $\pgcd(a,b) = 1$. Montrer que si $a \mid c$ et
$b \mid c$, alors $ab \mid c$.
\end{exercice}

\begin{exercice}[$\star\star\star$]
\label{exo:ch1:olympiade1}
Soient $a, b \in \N^*$ tels que $\pgcd(a,b) = 1$. Montrer que tout entier
$n \geq (a-1)(b-1)$ peut s'écrire sous la forme $n = xa + yb$ avec
$x, y \in \N$. En déduire que le plus grand entier qui \emph{ne peut pas} être
représenté est $ab - a - b$.
\end{exercice}

\begin{exercice}[$\star\star\star$]
\label{exo:ch1:olympiade2}
Soit $n \geq 2$ un entier. Montrer que
\[
  \sum_{k=1}^{n} \pgcd(k, n) = \sum_{d \mid n} d\, \varphi(n/d),
\]
où $\varphi$ désigne l'indicatrice d'Euler. \textit{(On pourra regrouper les termes
de la somme de gauche selon la valeur de $\pgcd(k,n)$.)}
\end{exercice}

\begin{exercice}[$\star\star\star$]
\label{exo:ch1:olympiade3}
\index{nombres de Fibonacci!et pgcd}
Montrer que $\pgcd(F_m, F_n) = F_{\pgcd(m,n)}$ pour tous $m, n \geq 1$, où $(F_n)$
est la suite de Fibonacci.
\end{exercice}

% ============================================================
\section{Résumé du chapitre}
\label{sec:ch1:resume}

\begin{tcolorbox}[enhanced,colback=blue!3,colframe=blue!50!black,
  title={\textbf{Chapitre 1 --- Résumé}},fonttitle=\large,
  breakable]
\begin{itemize}[leftmargin=1.5cm]
  \item \textbf{Divisibilité}~: $a \mid b \iff \exists k \in \Z,\; b = ka$.
    Relation d'ordre partiel sur $\N^*$.
  \item \textbf{Division euclidienne}~: $a = bq + r$ avec $0 \leq r < \abs{b}$~;
    existence et unicité.
  \item \textbf{Pgcd}~: plus grand diviseur commun~; calculé efficacement par
    l'algorithme d'Euclide en $O(\log(\min(a,b)))$ divisions.
  \item \textbf{Identité de Bézout}~: $\pgcd(a,b) = au + bv$ pour certains
    $u, v \in \Z$~; les coefficients sont calculés par l'algorithme étendu.
  \item \textbf{Ppcm}~: $\pgcd(a,b) \cdot \ppcm(a,b) = \abs{ab}$.
  \item \textbf{Lemme de Gauss}~: si $a \mid bc$ et $\pgcd(a,b) = 1$, alors
    $a \mid c$.
  \item \textbf{Corollaire fondamental}~: si $p$ premier divise $ab$, alors
    $p \mid a$ ou $p \mid b$.
\end{itemize}
\end{tcolorbox}


% ============================================================
% ============================================================
%    CHAPITRE 2 : NOMBRES PREMIERS — INFINITÉ, DISTRIBUTION,
%                 CRIBLE D'ÉRATOSTHÈNE
% ============================================================
% ============================================================

\chapter{Nombres Premiers --- Infinité, Distribution, Crible d'Ératosthène}
\label{ch:premiers}

\index{nombre premier}

\begin{quote}
\itshape
\og Les nombres premiers sont à la théorie des nombres ce que les atomes sont à la
chimie~: les constituants fondamentaux à partir desquels tout le reste est
construit. \fg{}
\end{quote}

\bigskip

\section{Introduction historique}
\label{sec:ch2:intro}

\index{nombre premier!histoire}

Les nombres premiers ont fasciné les mathématiciens depuis l'Antiquité.
Euclide\index{Euclide} a démontré leur infinitude dans la Proposition~20 du
Livre~IX de ses \textit{Éléments} --- une démonstration dont l'élégance reste
inégalée après plus de deux millénaires. Euler\index{Euler, Leonhard}, au
\textsc{xviii}\textsuperscript{e} siècle, a donné une preuve analytique de
l'infinitude des premiers en montrant que la série $\sum_{p} 1/p$, étendue aux
nombres premiers, diverge. Cette approche a ouvert la voie à la théorie analytique
des nombres.

Le \textsc{xix}\textsuperscript{e} siècle a vu l'émergence de la question centrale~:
\emph{comment les nombres premiers sont-ils distribués parmi les entiers ?} Les
travaux de Legendre\index{Legendre, Adrien-Marie},
Gauss\index{Gauss, Carl Friedrich},
Chebyshev\index{Chebyshev, Pafnouti},
Riemann\index{Riemann, Bernhard},
Hadamard\index{Hadamard, Jacques} et
de la Vallée-Poussin\index{Vallée-Poussin, Charles de la} ont culminé dans le
\emph{théorème des nombres premiers}\index{théorème des nombres premiers}, démontré
indépendamment par Hadamard et de la Vallée-Poussin en 1896.

Au \textsc{xx}\textsuperscript{e} siècle, Paul Erdős\index{Erdos@Erdős, Paul} et
Atle Selberg\index{Selberg, Atle} ont donné en 1949 une preuve \og élémentaire \fg{}
(n'utilisant pas l'analyse complexe) du théorème des nombres premiers, un résultat qui
a surpris la communauté mathématique.

% ============================================================
\section{Définitions fondamentales}
\label{sec:ch2:definitions}

\begin{definition}[Nombre premier, nombre composé]
\label{def:premier}
\index{nombre premier!définition}
\index{nombre composé}
Un entier $p \geq 2$ est dit \emph{premier} si ses seuls diviseurs positifs sont $1$
et $p$. Un entier $n \geq 2$ qui n'est pas premier est dit \emph{composé}~: il
possède un diviseur $d$ vérifiant $1 < d < n$.

Par convention, $1$ n'est ni premier ni composé.
\end{definition}

\begin{remarque}
\label{rem:premier:2}
Le nombre $2$ est le seul nombre premier pair\index{nombre premier!pair}. Tout nombre
premier $p \geq 3$ est impair.
\end{remarque}

\begin{proposition}[Existence d'un diviseur premier]
\label{prop:diviseur_premier}
\index{nombre premier!diviseur premier}
Tout entier $n \geq 2$ possède un diviseur premier. De plus, le plus petit diviseur
$d > 1$ de $n$ est premier.
\end{proposition}

\begin{proof}
Soit $d$ le plus petit diviseur de $n$ supérieur à $1$. Un tel $d$ existe car $n$
se divise lui-même. Montrons que $d$ est premier. Par l'absurde, si $d$ est composé,
il possède un diviseur $d'$ avec $1 < d' < d$. Par transitivité de la divisibilité,
$d' \mid n$, ce qui contredit la minimalité de $d$.
\end{proof}

\begin{corollaire}
\label{cor:diviseur_premier_sqrt}
Si $n \geq 2$ est composé, alors $n$ possède un diviseur premier $p \leq \sqrt{n}$.
\end{corollaire}

\begin{proof}
Si $n$ est composé, $n = ab$ avec $1 < a \leq b < n$. Si $a > \sqrt{n}$, alors
$n = ab > \sqrt{n} \cdot \sqrt{n} = n$, contradiction. Donc $a \leq \sqrt{n}$, et
$a$ possède un diviseur premier $p \leq a \leq \sqrt{n}$.
\end{proof}

% ============================================================
\section{Le théorème fondamental de l'arithmétique}
\label{sec:ch2:fta}

\index{théorème fondamental de l'arithmétique}

\begin{theoreme}[Théorème fondamental de l'arithmétique]
\label{thm:fta}
Tout entier $n \geq 2$ s'écrit de manière unique (à l'ordre des facteurs près)
comme produit de nombres premiers~:
\[
  n = p_1^{\alpha_1} p_2^{\alpha_2} \cdots p_k^{\alpha_k},
\]
où $p_1 < p_2 < \cdots < p_k$ sont des nombres premiers et
$\alpha_1, \alpha_2, \ldots, \alpha_k \geq 1$.
\end{theoreme}

\begin{proof}
\textbf{Existence} (par récurrence forte sur $n$).

\emph{Base}~: $n = 2$ est premier, donc $n = 2$ est un produit de premiers (un seul
facteur).

\emph{Hérédité}~: soit $n \geq 3$. Supposons que tout entier $m$ vérifiant
$2 \leq m < n$ s'écrive comme produit de premiers.
\begin{itemize}
  \item Si $n$ est premier, alors $n$ est un produit de premiers.
  \item Si $n$ est composé, alors $n = ab$ avec $2 \leq a, b < n$. Par hypothèse de
    récurrence, $a$ et $b$ s'écrivent comme produits de premiers, donc $n = ab$ aussi.
\end{itemize}

\medskip
\textbf{Unicité} (par récurrence forte sur $n$).

Supposons que $n$ admette deux factorisations~:
\[
  n = p_1 p_2 \cdots p_r = q_1 q_2 \cdots q_s,
\]
où les $p_i$ et $q_j$ sont premiers (non nécessairement distincts, non nécessairement
ordonnés).

Comme $p_1 \mid n = q_1 q_2 \cdots q_s$, le

ef{cor:gauss:premier_divise_produit} (appliqué par récurrence) donne $p_1 \mid q_j$
pour un certain $j$. Comme $q_j$ est premier et $p_1 \geq 2$, on a $p_1 = q_j$.

En simplifiant par $p_1$~:
\[
  p_2 \cdots p_r = q_1 \cdots q_{j-1}\, q_{j+1} \cdots q_s.
\]
Si $r = 1$, alors le membre gauche vaut $1$, donc $s = 1$ et les deux factorisations
sont identiques. Si $r \geq 2$, le produit $p_2 \cdots p_r < n$, et par hypothèse de
récurrence, les deux factorisations du membre simplifié sont identiques (à l'ordre
près). En réintroduisant $p_1 = q_j$, les deux factorisations de $n$ sont identiques
à l'ordre près.
\end{proof}

\begin{notation}[Factorisation canonique]
\label{not:factorisation}
\index{factorisation canonique}
On écrira la factorisation canonique de $n$ sous la forme
\[
  n = \prod_{p \mid n} p^{v_p(n)},
\]
où $v_p(n)$\index{valuation p-adique@valuation $p$-adique} désigne la
\emph{valuation $p$-adique} de $n$, c'est-à-dire le plus grand exposant $\alpha$
tel que $p^\alpha \mid n$.
\end{notation}

\begin{proposition}[Pgcd et ppcm via les factorisations]
\label{prop:pgcd_ppcm_fact}
\index{pgcd!via factorisation}
\index{ppcm!via factorisation}
Si $a = \prod_p p^{\alpha_p}$ et $b = \prod_p p^{\beta_p}$, alors~:
\[
  \pgcd(a,b) = \prod_p p^{\min(\alpha_p, \beta_p)}, \qquad
  \ppcm(a,b) = \prod_p p^{\max(\alpha_p, \beta_p)}.
\]
\end{proposition}

\begin{proof}
Posons $d = \prod_p p^{\min(\alpha_p, \beta_p)}$. Pour tout premier $p$, on a
$v_p(d) = \min(\alpha_p, \beta_p) \leq \alpha_p = v_p(a)$, donc $d \mid a$. De
même $d \mid b$. Si $c \mid a$ et $c \mid b$, alors $v_p(c) \leq \min(\alpha_p,
\beta_p) = v_p(d)$ pour tout $p$, donc $c \mid d$. Ainsi $d = \pgcd(a,b)$.

La preuve pour le ppcm est analogue avec $\max$ au lieu de $\min$.
\end{proof}

\begin{exemple}
\label{ex:pgcd_ppcm_fact}
$360 = 2^3 \cdot 3^2 \cdot 5$ et $756 = 2^2 \cdot 3^3 \cdot 7$. Donc~:
\begin{align*}
  \pgcd(360, 756) &= 2^{\min(3,2)} \cdot 3^{\min(2,3)} \cdot 5^{\min(1,0)} \cdot
    7^{\min(0,1)} = 2^2 \cdot 3^2 = 36, \\
  \ppcm(360, 756) &= 2^{\max(3,2)} \cdot 3^{\max(2,3)} \cdot 5^{\max(1,0)} \cdot
    7^{\max(0,1)} = 2^3 \cdot 3^3 \cdot 5 \cdot 7 = 7560.
\end{align*}
Vérifions~: $36 \times 7560 = 272160 = 360 \times 756$. \checkmark
\end{exemple}

% ============================================================
\section{Infinitude des nombres premiers}
\label{sec:ch2:infinite}

\index{nombre premier!infinitude}

Nous présentons trois démonstrations classiques de ce résultat fondamental.

\begin{theoreme}[Infinitude des nombres premiers --- Euclide]
\label{thm:euclide_infinite}
\index{Euclide!théorème de l'infinitude des premiers}
L'ensemble des nombres premiers est infini.
\end{theoreme}

\subsection{Première preuve~: par l'absurde (Euclide)}
\label{subsec:ch2:preuve1}

\begin{proof}[Preuve d'Euclide]
Supposons par l'absurde que l'ensemble des nombres premiers soit fini~:
$\{p_1, p_2, \ldots, p_n\}$. Considérons l'entier
\[
  N = p_1 p_2 \cdots p_n + 1.
\]
Par la 
ef{prop:diviseur_premier}, $N$ possède un diviseur premier $p$. Or pour
tout $i \in \{1, \ldots, n\}$, $p_i \mid p_1 p_2 \cdots p_n$, donc
$p_i \mid (N - p_1 \cdots p_n) = 1$ si $p = p_i$, ce qui est impossible. Donc
$p \neq p_i$ pour tout $i$, ce qui contredit l'hypothèse que $\{p_1, \ldots, p_n\}$
est l'ensemble de \emph{tous} les premiers.
\end{proof}

\subsection{Deuxième preuve~: par la divergence de $\sum 1/p$ (Euler)}
\label{subsec:ch2:preuve2}

\index{Euler, Leonhard!preuve de l'infinitude des premiers}

\begin{theoreme}[Euler, 1737]
\label{thm:euler_divergence}
La série $\displaystyle\sum_{p \text{ premier}} \frac{1}{p}$ diverge.
En particulier, l'ensemble des nombres premiers est infini.
\end{theoreme}

\begin{proof}
Pour $x \geq 2$, considérons le produit eulérien\index{produit eulérien}~:
\[
  \prod_{p \leq x} \frac{1}{1 - 1/p}
  = \prod_{p \leq x} \sum_{k=0}^{\infty} \frac{1}{p^k}
  \geq \sum_{n=1}^{\floor{x}} \frac{1}{n},
\]
la dernière inégalité résultant du fait que le membre de droite est une sous-somme du
développement du produit (par le théorème fondamental de l'arithmétique, chaque entier
$n \leq x$ a tous ses facteurs premiers $\leq x$).

Or la série harmonique diverge~: $\sum_{n=1}^{N} 1/n \geq \ln N$. Donc
\[
  \prod_{p \leq x} \frac{1}{1 - 1/p} \to +\infty \quad \text{quand } x \to +\infty.
\]

En prenant le logarithme~:
\[
  \sum_{p \leq x} -\ln\!\left(1 - \frac{1}{p}\right) \to +\infty.
\]
Or $-\ln(1 - t) = t + t^2/2 + t^3/3 + \cdots \leq t + t^2$ pour $0 < t \leq 1/2$,
donc
\[
  \sum_{p \leq x} \frac{1}{p} \geq \sum_{p \leq x} \left(-\ln\!\left(1-\frac{1}{p}\right)\right) - \sum_{p \leq x} \frac{1}{p^2}.
\]
Comme $\sum_{p} 1/p^2 \leq \sum_{n=1}^{\infty} 1/n^2 = \pi^2/6 < \infty$, et le
membre de gauche de l'inégalité tend vers $+\infty$, on conclut que
$\sum_{p} 1/p = +\infty$.
\end{proof}

\subsection{Troisième preuve~: par les nombres de Fermat}
\label{subsec:ch2:preuve3}

\index{Fermat, Pierre de!nombres de Fermat}

\begin{proof}[Preuve par les nombres de Fermat]
Considérons les nombres de Fermat $F_n = 2^{2^n} + 1$ pour $n \geq 0$~:
\[
  F_0 = 3, \quad F_1 = 5, \quad F_2 = 17, \quad F_3 = 257, \quad F_4 = 65537, \quad \ldots
\]

\textbf{Lemme.} Pour $m \neq n$, $\pgcd(F_m, F_n) = 1$.

\emph{Preuve du lemme.} Sans perte de généralité, supposons $m < n$. On montre
par récurrence que
\[
  F_0 F_1 \cdots F_{n-1} = F_n - 2.
\]
En effet, $F_0 = 3 = F_1 - 2$, et si la formule est vraie au rang $n-1$, alors
$F_0 \cdots F_{n-1} = F_0 \cdots F_{n-2} \cdot F_{n-1} = (F_{n-1}-2) F_{n-1}
= (2^{2^{n-1}}-1)(2^{2^{n-1}}+1) = 2^{2^n}-1 = F_n - 2$.

Si $d \mid F_m$ et $d \mid F_n$, alors $d \mid (F_n - F_0 \cdots F_{n-1}) = 2$.
Comme $F_m$ est impair (car $2^{2^m}$ est pair), $d$ est impair, donc $d = 1$.

\medskip

Puisque les $F_n$ sont deux à deux premiers entre eux et chacun $\geq 3$, chacun
possède au moins un diviseur premier, et ces diviseurs premiers sont distincts. Donc
il y a au moins autant de premiers que de nombres de Fermat, soit une infinité.
\end{proof}

% ============================================================
\section{Le crible d'Ératosthène}
\label{sec:ch2:eratosthene}

\index{crible d'Ératosthène}
\index{Eratosthene@Ératosthène}

Le crible d'Ératosthène est un algorithme ancien et élégant pour trouver tous les
nombres premiers jusqu'à un entier $N$ donné.

\subsection{Description de l'algorithme}
\label{subsec:ch2:crible_algo}

\begin{enumerate}[label=\textbf{Étape \arabic*.}]
  \item Écrire la liste des entiers de $2$ à $N$.
  \item Le plus petit nombre non barré ($p = 2$ au départ) est premier. Barrer tous
    ses multiples stricts $2p, 3p, 4p, \ldots$ qui sont $\leq N$.
  \item Passer au prochain nombre non barré et répéter l'étape~2.
  \item S'arrêter lorsque le prochain nombre non barré dépasse $\sqrt{N}$.
  \item Les nombres non barrés restants sont premiers.
\end{enumerate}

\begin{proposition}[Correction du crible]
\label{prop:crible:correct}
Le crible d'Ératosthène identifie correctement tous les nombres premiers jusqu'à $N$.
\end{proposition}

\begin{proof}
Un nombre $n \leq N$ est barré si et seulement s'il possède un diviseur premier
$p \leq \sqrt{N}$, ce qui équivaut à être composé (par le

ef{cor:diviseur_premier_sqrt}). Les nombres non barrés sont donc exactement les
premiers.
\end{proof}

\begin{proposition}[Complexité du crible]
\label{prop:crible:complexite}
\index{crible d'Ératosthène!complexité}
Le crible d'Ératosthène a une complexité en $O(N \log \log N)$ opérations.
\end{proposition}

\begin{proof}[Esquisse de preuve]
Le nombre d'opérations de barrage est
\[
  \sum_{p \leq \sqrt{N}} \frac{N}{p} = N \sum_{p \leq \sqrt{N}} \frac{1}{p}
  \sim N \log\log\sqrt{N} \sim N \log\log N,
\]
en utilisant le théorème de Mertens\index{Mertens, théorème de} selon lequel
$\sum_{p \leq x} 1/p \sim \log\log x$.
\end{proof}

\subsection{Visualisation du crible}
\label{subsec:ch2:crible_tikz}

\begin{center}
\begin{tikzpicture}[
    cell/.style={minimum size=7mm, draw, font=\small, anchor=center},
    prime/.style={cell, fill=blue!20, font=\small\bfseries},
    sieved/.style={cell, fill=gray!20, text=gray!50, font=\small},
  ]
  % Grille 6 x 5 pour les nombres 2 à 30 (on omet 1)
  \def\primes{2,3,5,7,11,13,17,19,23,29}
  \foreach \num in {2,...,30} {
    \pgfmathtruncatemacro{\col}{mod(\num-2,6)}
    \pgfmathtruncatemacro{\row}{div(\num-2,6)}
    \gdef\isprime{0}
    \foreach \p in \primes {
      \ifnum\num=\p\relax
        \gdef\isprime{1}
      \fi
    }
    \ifnum\isprime=1
      \node[prime] at (\col*0.9, -\row*0.9) {\num};
    \else
      \node[sieved] at (\col*0.9, -\row*0.9) {\num};
    \fi
  }
  % Légende
  \node[prime, label={right:Premier}] at (7, 0) {};
  \node[sieved, label={right:Composé}] at (7, -0.9) {};
\end{tikzpicture}
\captionof{figure}{Résultat du crible d'Ératosthène pour $N = 30$.}
\label{fig:crible_30}
\end{center}

% ============================================================
\section{La fonction de comptage $\pi(x)$}
\label{sec:ch2:pi_x}

\index{fonction de comptage des premiers}
\index{pi(x)@$\pi(x)$}

\begin{definition}[Fonction $\pi(x)$]
\label{def:pi_x}
Pour $x \geq 0$ réel, on note $\pi(x)$ le nombre de nombres premiers $p$ tels que
$p \leq x$~:
\[
  \pi(x) = \#\{p \leq x : p \text{ est premier}\}.
\]
\end{definition}

\begin{exemple}[Valeurs de $\pi(x)$]
\label{ex:pi_valeurs}
\begin{center}
\begin{tabular}{c*{10}{c}}
\toprule
$x$ & 10 & 20 & 50 & 100 & 200 & 500 & 1000 & $10^4$ & $10^6$ & $10^9$ \\
\midrule
$\pi(x)$ & 4 & 8 & 15 & 25 & 46 & 95 & 168 & 1229 & 78498 & 50847534 \\
\bottomrule
\end{tabular}
\end{center}
\end{exemple}

\begin{proposition}[Formule de Legendre]
\label{prop:legendre}
\index{Legendre, Adrien-Marie!formule de}
Pour tout entier $N \geq 2$, si $p_1, p_2, \ldots, p_k$ sont les nombres premiers
inférieurs ou égaux à $\sqrt{N}$, alors
\[
  \pi(N) - \pi(\sqrt{N}) + 1
  = N - \sum_{i} \floor{\frac{N}{p_i}}
  + \sum_{i < j} \floor{\frac{N}{p_i p_j}}
  - \cdots
  + (-1)^k \floor{\frac{N}{p_1 p_2 \cdots p_k}}.
\]
C'est une application directe du principe d'inclusion-exclusion\index{inclusion-exclusion, principe d'}.
\end{proposition}

\begin{proof}
Un entier $n$ avec $1 \leq n \leq N$ est soit $1$, soit premier $> \sqrt{N}$, soit
admet un diviseur premier $\leq \sqrt{N}$. Les entiers de cette troisième catégorie
sont exactement ceux divisibles par au moins un $p_i$. Le principe
d'inclusion-exclusion compte les entiers $\leq N$ non divisibles par aucun des $p_i$,
ce qui donne exactement $1$ plus les premiers dans $(\sqrt{N}, N]$, soit
$1 + \pi(N) - \pi(\sqrt{N})$.
\end{proof}

% ============================================================
\section{Le théorème des nombres premiers}
\label{sec:ch2:pnt}

\index{théorème des nombres premiers}

\begin{theoreme}[Théorème des nombres premiers]
\label{thm:pnt}
Quand $x \to +\infty$,
\[
  \pi(x) \sim \frac{x}{\ln x}.
\]
Autrement dit, $\displaystyle\lim_{x \to +\infty} \frac{\pi(x)}{x / \ln x} = 1$.
\end{theoreme}

Ce théorème, conjecturé par Legendre\index{Legendre, Adrien-Marie} et
Gauss\index{Gauss, Carl Friedrich} vers 1800, a été démontré indépendamment par
Hadamard\index{Hadamard, Jacques} et de la
Vallée-Poussin\index{Vallée-Poussin, Charles de la} en 1896, en utilisant les
propriétés analytiques de la fonction zêta de Riemann $\zeta(s) = \sum_{n=1}^{\infty}
n^{-s}$\index{Riemann!fonction zêta de}. La preuve repose sur le fait que $\zeta(s)$
n'a pas de zéros sur la droite $\operatorname{Re}(s) = 1$.

\begin{remarque}[Heuristique probabiliste]
\label{rem:pnt:heuristique}
On peut interpréter le théorème des nombres premiers de la manière suivante~: un
entier $n$ choisi \og au hasard \fg{} est premier avec \og probabilité \fg{}
$\approx 1/\ln n$. Plus précisément, la \og densité locale \fg{} des premiers
autour de $n$ est $\sim 1/\ln n$. Ainsi~:
\[
  \pi(x) \approx \int_2^x \frac{dt}{\ln t} = \operatorname{Li}(x),
\]
où $\operatorname{Li}(x)$\index{Li(x)@$\operatorname{Li}(x)$} est l'intégrale
logarithmique. En fait, $\operatorname{Li}(x)$ fournit une meilleure approximation de
$\pi(x)$ que $x/\ln x$.
\end{remarque}

\begin{theoreme}[Encadrement de Chebyshev]
\label{thm:chebyshev}
\index{Chebyshev, Pafnouti!encadrement}
Il existe des constantes $c_1, c_2 > 0$ telles que pour tout $x$ assez grand~:
\[
  c_1 \frac{x}{\ln x} \leq \pi(x) \leq c_2 \frac{x}{\ln x}.
\]
Plus précisément, Chebyshev a montré en 1852 que l'on peut prendre $c_1 = \ln 2 / 2
\approx 0{,}347$ et $c_2 = 6 \ln 2 / 5 \approx 0{,}831$, et qu'en fait, si la
limite $\lim_{x \to \infty} \pi(x) / (x/\ln x)$ existe, elle vaut $1$.
\end{theoreme}

\begin{remarque}
\label{rem:chebyshev:postulat}
Chebyshev a également démontré le \textbf{postulat de Bertrand}\index{Bertrand,
postulat de}~: pour tout $n \geq 1$, il existe au moins un nombre premier $p$ tel que
$n < p \leq 2n$. Ce résultat, conjecturé par Bertrand en 1845, avait été vérifié
numériquement par Bertrand jusqu'à $n = 6 \times 10^6$.
\end{remarque}

% ============================================================
\section{Nombres premiers remarquables}
\label{sec:ch2:speciaux}

\subsection{Nombres premiers de Mersenne}
\label{subsec:ch2:mersenne}

\index{Mersenne!nombres premiers de}

\begin{definition}[Nombre de Mersenne]
\label{def:mersenne}
Un \emph{nombre de Mersenne}\index{Mersenne!nombre de} est un entier de la forme
$M_n = 2^n - 1$, où $n \in \N^*$. Si $M_n$ est premier, on parle de
\emph{premier de Mersenne}.
\end{definition}

\begin{proposition}
\label{prop:mersenne:necessaire}
Si $M_n = 2^n - 1$ est premier, alors $n$ est premier.
\end{proposition}

\begin{proof}
Si $n = ab$ avec $1 < a, b < n$, alors
\[
  2^n - 1 = (2^a)^b - 1 = (2^a - 1)\bigl((2^a)^{b-1} + (2^a)^{b-2} + \cdots + 1\bigr),
\]
et $2^a - 1 > 1$, donc $M_n$ est composé.
\end{proof}

\begin{remarque}
\label{rem:mersenne:liste}
La réciproque est fausse~: $M_{11} = 2^{11} - 1 = 2047 = 23 \times 89$ n'est pas
premier. Les premiers de Mersenne connus en 2024 sont au nombre de $51$, le plus
grand étant $M_{82589933}$, un nombre à plus de $24$ millions de chiffres. La
recherche de nouveaux premiers de Mersenne est menée par le projet
GIMPS\index{GIMPS} (Great Internet Mersenne Prime Search).
\end{remarque}

\subsection{Nombres premiers de Fermat}
\label{subsec:ch2:fermat}

\index{Fermat, Pierre de!nombres premiers de}

\begin{definition}[Nombre de Fermat]
\label{def:fermat_number}
Le $n$-ième \emph{nombre de Fermat} est $F_n = 2^{2^n} + 1$.
\end{definition}

\begin{proposition}
\label{prop:fermat:necessaire}
Si $2^m + 1$ est premier, alors $m$ est une puissance de $2$.
\end{proposition}

\begin{proof}
Si $m$ a un facteur impair $d > 1$, disons $m = de$, alors
\[
  2^m + 1 = (2^e)^d + 1 = (2^e + 1)\bigl((2^e)^{d-1} - (2^e)^{d-2} + \cdots + 1\bigr),
\]
(car $d$ est impair), et $2^e + 1 > 1$, donc $2^m + 1$ est composé.
\end{proof}

\begin{remarque}
\label{rem:fermat:connus}
Fermat conjecturait que tous les $F_n$ sont premiers. Cela est vrai pour $n = 0, 1,
2, 3, 4$ ($F_4 = 65537$), mais Euler\index{Euler, Leonhard} a montré en 1732 que
$F_5 = 2^{32} + 1 = 4294967297 = 641 \times 6700417$ est composé. Aucun autre
premier de Fermat n'a été trouvé depuis~; il est conjecturé que seuls $F_0, \ldots,
F_4$ sont premiers.
\end{remarque}

\subsection{Nombres premiers jumeaux et conjecture de Goldbach}
\label{subsec:ch2:jumeaux_goldbach}

\begin{definition}[Nombres premiers jumeaux]
\label{def:jumeaux}
\index{nombres premiers jumeaux}
Deux nombres premiers $p$ et $p+2$ sont dits \emph{jumeaux}. Exemples~: $(3,5)$,
$(5,7)$, $(11,13)$, $(17,19)$, $(29,31)$, $(41,43)$, $\ldots$
\end{definition}

\begin{remarque}[Conjecture des premiers jumeaux]
\label{rem:jumeaux:conjecture}
\index{conjecture!des nombres premiers jumeaux}
Il est conjecturé qu'il existe une infinité de paires de premiers jumeaux. En 2013,
Yitang Zhang\index{Zhang, Yitang} a démontré qu'il existe une infinité de paires de
premiers $(p, p+k)$ avec $k \leq 70\,000\,000$. Depuis, le projet
Polymath\index{Polymath} a réduit cette borne à $246$~: il existe une infinité de
paires de premiers distants d'au plus $246$.
\end{remarque}

\begin{remarque}[Conjecture de Goldbach]
\label{rem:goldbach}
\index{conjecture!de Goldbach}
\index{Goldbach, Christian}
La \emph{conjecture de Goldbach} (1742) affirme que tout entier pair $n \geq 4$
est somme de deux nombres premiers. Elle a été vérifiée numériquement jusqu'à
$4 \times 10^{18}$, mais reste non démontrée.

Des résultats partiels importants existent~:
\begin{itemize}
  \item \textbf{Théorème de Vinogradov\index{Vinogradov, Ivan}} (1937)~: tout entier
    impair suffisamment grand est somme de trois nombres premiers.
  \item \textbf{Théorème de Chen\index{Chen Jingrun}} (1966)~: tout entier pair
    suffisamment grand est somme d'un premier et d'un produit d'au plus deux premiers
    (un \og $P_2$ \fg{}).
  \item \textbf{Théorème de Helfgott\index{Helfgott, Harald}} (2013)~: tout entier
    impair $\geq 7$ est somme de trois nombres premiers (version effective du
    théorème de Vinogradov).
\end{itemize}
\end{remarque}

% ============================================================
\section{Visualisations}
\label{sec:ch2:visualisations}

\subsection{Spirale d'Ulam}
\label{subsec:ch2:ulam}

\index{Ulam, spirale d'}

La spirale d'Ulam\index{spirale d'Ulam}, découverte par Stanislaw
Ulam\index{Ulam, Stanislaw} en 1963, est obtenue en écrivant les entiers en spirale à
partir du centre et en marquant les nombres premiers. Des structures diagonales
étonnantes apparaissent.

\begin{center}
\begin{tikzpicture}[scale=0.42]
  % Spirale d'Ulam 11x11 centrée sur 1
  % Coordonnées générées pour les nombres 1 à 121
  % Les premiers parmi 1..121 sont :
  % 2,3,5,7,11,13,17,19,23,29,31,37,41,43,47,53,59,61,67,71,73,79,83,89,97,101,103,107,109,113
  %
  % Positions dans la spirale (x,y) en partant de (0,0) pour n=1
  \def\spiralcoords{
    1/0/0, 2/1/0, 3/1/1, 4/0/1, 5/-1/1, 6/-1/0, 7/-1/-1, 8/0/-1, 9/1/-1,
    10/2/-1, 11/2/0, 12/2/1, 13/2/2, 14/1/2, 15/0/2, 16/-1/2, 17/-2/2,
    18/-2/1, 19/-2/0, 20/-2/-1, 21/-2/-2, 22/-1/-2, 23/0/-2, 24/1/-2,
    25/2/-2, 26/3/-2, 27/3/-1, 28/3/0, 29/3/1, 30/3/2, 31/3/3, 32/2/3,
    33/1/3, 34/0/3, 35/-1/3, 36/-2/3, 37/-3/3, 38/-3/2, 39/-3/1,
    40/-3/0, 41/-3/-1, 42/-3/-2, 43/-3/-3, 44/-2/-3, 45/-1/-3,
    46/0/-3, 47/1/-3, 48/2/-3, 49/3/-3
  }
  \def\primevals{2,3,5,7,11,13,17,19,23,29,31,37,41,43,47}
  % Dessiner les premiers
  \foreach \n/\x/\y in \spiralcoords {
    \gdef\isp{0}
    \foreach \p in \primevals {
      \ifnum\n=\p\relax \gdef\isp{1} \fi
    }
    \ifnum\isp=1
      \fill[blue!70!black] (\x,\y) rectangle ++(0.85,0.85);
    \else
      \fill[gray!15] (\x,\y) rectangle ++(0.85,0.85);
    \fi
    \draw[gray!50, very thin] (\x,\y) rectangle ++(0.85,0.85);
    \node[font=\tiny, text=black] at (\x+0.425,\y+0.425) {\n};
  }
\end{tikzpicture}
\captionof{figure}{Spirale d'Ulam (partielle)~: les cases bleues sont les nombres premiers.}
\label{fig:ulam}
\end{center}

\subsection{Graphe de la fonction $\pi(x)$}
\label{subsec:ch2:graphe_pi}

\begin{center}
\begin{tikzpicture}
\begin{axis}[
    width=12cm, height=7cm,
    xlabel={$x$},
    ylabel={$\pi(x)$},
    xmin=0, xmax=105,
    ymin=0, ymax=30,
    grid=major,
    grid style={gray!30},
    legend pos=north west,
    legend style={font=\small},
    axis lines=left,
    every axis label/.style={font=\small},
    tick label style={font=\small},
  ]
  % pi(x) en escalier
  \addplot[blue!70!black, thick, const plot mark left] coordinates {
    (0,0) (2,1) (3,2) (5,3) (7,4) (11,5) (13,6) (17,7) (19,8) (23,9)
    (29,10) (31,11) (37,12) (41,13) (43,14) (47,15) (53,16) (59,17)
    (61,18) (67,19) (71,20) (73,21) (79,22) (83,23) (89,24) (97,25) (101,26)
  };
  \addlegendentry{$\pi(x)$}
  % x / ln(x)
  \addplot[red!70!black, thick, dashed, domain=2.1:105, samples=200] {x/ln(x)};
  \addlegendentry{$x/\ln x$}
  % Li(x) approximation
  \addplot[green!50!black, thick, dotted, domain=2.1:105, samples=200]
    {x/ln(x) + x/(ln(x))^2};
  \addlegendentry{$\operatorname{Li}(x)$ (approx.)}
\end{axis}
\end{tikzpicture}
\captionof{figure}{La fonction $\pi(x)$ (escalier bleu) comparée à $x/\ln x$ (rouge
pointillé) et une approximation de $\operatorname{Li}(x)$ (vert pointillé).}
\label{fig:pi_graph}
\end{center}

% ============================================================
\section{Connexions avec la cryptographie~: le système RSA}
\label{sec:ch2:rsa}

\index{RSA!rôle des grands premiers}

La sécurité du système RSA\index{RSA} repose sur la difficulté de factoriser de grands
entiers. Le principe est le suivant~:

\begin{enumerate}[label=\textbf{\arabic*.}]
  \item \textbf{Génération des clés.} On choisit deux grands nombres premiers $p$ et
    $q$ (typiquement de $1024$ bits chacun) et on calcule $n = pq$. Le nombre $n$ est
    rendu public.
  \item \textbf{Chiffrement.} Pour chiffrer un message $m < n$, on calcule
    $c = m^e \bmod n$, où $e$ est l'exposant public.
  \item \textbf{Déchiffrement.} Le destinataire, connaissant $p$ et $q$, calcule
    $d = e^{-1} \bmod \varphi(n)$ et retrouve $m = c^d \bmod n$.
\end{enumerate}

Un adversaire souhaitant casser le système devrait factoriser $n$ pour retrouver $p$
et $q$. Le meilleur algorithme connu (le crible du corps de
nombres\index{crible du corps de nombres}) a une complexité sous-exponentielle
$L_n[1/3, c]$, ce qui le rend infaisable pour des $n$ de $2048$ bits ou plus.

C'est pourquoi la génération de grands nombres premiers est un enjeu crucial en
cryptographie. Les tests de primalité\index{test de primalité} (Miller--Rabin,
AKS)\index{Miller--Rabin}\index{AKS} permettent de vérifier efficacement qu'un nombre
est premier, et le théorème des nombres premiers garantit que les premiers sont
suffisamment denses~: en moyenne, il faut tester environ $\ln n \approx 710$ candidats
pour trouver un premier de $1024$ bits.

% ============================================================
\section{Exercices}
\label{sec:ch2:exercices}

\begin{exercice}[$\star$]
\label{exo:ch2:premiers_100}
Trouver tous les nombres premiers inférieurs à $100$ à l'aide du crible
d'Ératosthène.
\end{exercice}

\begin{exercice}[$\star$]
\label{exo:ch2:factoriser}
Factoriser en produit de facteurs premiers les entiers suivants~:
\begin{enumerate}[label=(\alph*)]
  \item $2024$,
  \item $10!$,
  \item $\binom{20}{10}$.
\end{enumerate}
\end{exercice}

\begin{exercice}[$\star$]
\label{exo:ch2:mersenne}
Vérifier que $M_7 = 2^7 - 1 = 127$ est premier et que $M_{11} = 2^{11} - 1 = 2047$
est composé.
\end{exercice}

\begin{exercice}[$\star\star$]
\label{exo:ch2:infinitude_4k3}
\index{nombre premier!de la forme $4k+3$}
Montrer qu'il existe une infinité de nombres premiers de la forme $4k + 3$.
\textit{(Indication~: adapter la preuve d'Euclide en considérant $N = 4p_1 p_2
\cdots p_n - 1$.)}
\end{exercice}

\begin{exercice}[$\star\star$]
\label{exo:ch2:infinitude_6k5}
Montrer qu'il existe une infinité de nombres premiers de la forme $6k + 5$.
\end{exercice}

\begin{exercice}[$\star\star$]
\label{exo:ch2:valuation}
\index{valuation p-adique@valuation $p$-adique!de $n"!$}
Montrer la formule de Legendre\index{Legendre, Adrien-Marie!formule de la valuation}~:
pour tout premier $p$ et tout $n \geq 1$,
\[
  v_p(n!) = \sum_{k=1}^{\infty} \floor{\frac{n}{p^k}}.
\]
En déduire que $v_p(n!) = \dfrac{n - s_p(n)}{p - 1}$, où $s_p(n)$ est la somme des
chiffres de $n$ en base $p$.
\end{exercice}

\begin{exercice}[$\star\star$]
\label{exo:ch2:produit_consecutifs}
Montrer que le produit de $n$ entiers consécutifs est divisible par $n!$.
\textit{(Indication~: utiliser l'exercice précédent et les coefficients binomiaux.)}
\end{exercice}

\begin{exercice}[$\star\star$]
\label{exo:ch2:bertrand}
\index{Bertrand, postulat de}
En admettant le postulat de Bertrand, montrer que pour tout $n \geq 1$, il existe
$n$ nombres composés consécutifs. \textit{(Indication~: considérer $(n+1)! + 2,
(n+1)! + 3, \ldots, (n+1)! + (n+1)$.)}
\end{exercice}

\begin{exercice}[$\star\star\star$]
\label{exo:ch2:bonse}
\index{Bonse, inégalité de}
Montrer l'\emph{inégalité de Bonse}~: si $p_1 < p_2 < \cdots < p_n$ sont les $n$
premiers nombres premiers et $n \geq 4$, alors $p_n^2 < p_1 p_2 \cdots p_{n-1}$.
\end{exercice}

\begin{exercice}[$\star\star\star$]
\label{exo:ch2:olympiade1}
Soit $p$ un nombre premier. Montrer que $\sqrt{p}$ est irrationnel.
\end{exercice}

\begin{exercice}[$\star\star\star$]
\label{exo:ch2:olympiade2}
\index{Sylvester--Schur, théorème de}
Montrer que parmi $n$ entiers consécutifs supérieurs à $n$, au moins l'un d'entre eux
possède un facteur premier $> n$. \textit{(C'est le théorème de
Sylvester\index{Sylvester, James Joseph}.)}
\end{exercice}

\begin{exercice}[$\star\star\star$]
\label{exo:ch2:olympiade3}
Montrer qu'il existe une infinité de nombres premiers $p$ tels que $p \equiv 1
\pmod{4}$. \textit{(Indication~: considérer $(2p_1 p_2 \cdots p_n)^2 + 1$ et
montrer que tout facteur premier de ce nombre est $\equiv 1 \pmod{4}$.)}
\end{exercice}

% ============================================================
\section{Résumé du chapitre}
\label{sec:ch2:resume}

\begin{tcolorbox}[enhanced,colback=blue!3,colframe=blue!50!black,
  title={\textbf{Chapitre 2 --- Résumé}},fonttitle=\large,
  breakable]
\begin{itemize}[leftmargin=1.5cm]
  \item \textbf{Nombre premier}~: $p \geq 2$ dont les seuls diviseurs positifs
    sont $1$ et $p$.
  \item \textbf{Théorème fondamental de l'arithmétique}~: tout entier $\geq 2$
    s'écrit de manière unique comme produit de premiers (à l'ordre près).
  \item \textbf{Infinitude des premiers}~: démontrée par Euclide (argument par
    l'absurde), Euler (divergence de $\sum 1/p$), et via les nombres de Fermat.
  \item \textbf{Crible d'Ératosthène}~: algorithme en $O(N \log\log N)$ pour
    lister tous les premiers $\leq N$.
  \item \textbf{Théorème des nombres premiers}~: $\pi(x) \sim x / \ln x$.
  \item \textbf{Encadrement de Chebyshev}~: $c_1\, x/\ln x \leq \pi(x) \leq
    c_2\, x / \ln x$ pour $x$ assez grand.
  \item \textbf{Premiers remarquables}~: premiers de Mersenne ($2^p - 1$), de
    Fermat ($2^{2^n}+1$), premiers jumeaux.
  \item \textbf{Conjectures ouvertes}~: conjecture des premiers jumeaux, conjecture
    de Goldbach.
  \item \textbf{Applications}~: la densité des premiers et la difficulté de la
    factorisation sont au c\oe ur de la cryptographie RSA.
\end{itemize}
\end{tcolorbox}

% === FIN DU CHUNK preamble+ch1-2 ===
% === DÉBUT DU CHUNK ch3-4 ===

%%%%%%%%%%%%%%%%%%%%%%%%%%%%%%%%%%%%%%%%%%%%%%%%%%%%%%%%%%%%%%%%%%%%
%%  CHAPITRE 3 — Congruences et Arithmétique Modulaire
%%%%%%%%%%%%%%%%%%%%%%%%%%%%%%%%%%%%%%%%%%%%%%%%%%%%%%%%%%%%%%%%%%%%
\chapter{Congruences et Arithmétique Modulaire}
\label{ch:congruences}
\index{congruence}
\index{arithmétique modulaire}

\vspace{-0.5em}
\begin{center}\itshape
  « Les recherches dont les résultats sont exposés dans cet ouvrage
  ont pour objet la science des nombres entiers. »\\[2pt]
  --- Carl Friedrich Gauss, \emph{Disquisitiones Arithmeticae} (1801)
\end{center}
\bigskip

%% ── Introduction historique ─────────────────────────────────────
\section{Introduction historique}
\label{sec:cong-hist}

La notion de congruence\index{congruence!histoire}, telle que nous la connaissons
aujourd'hui, a été formalisée par Carl Friedrich Gauss\index{Gauss, Carl Friedrich}
dans ses célèbres \emph{Disquisitiones Arithmeticae}\index{Disquisitiones Arithmeticae},
publiées en~1801, alors qu'il n'avait que vingt-quatre ans.
Gauss introduisit la notation $a \equiv b \pmod{n}$, qui révolutionna la manière
de raisonner sur la divisibilité. Bien que des résultats isolés sur les restes
fussent connus depuis l'Antiquité --- les Chinois étudiaient déjà les systèmes
de congruences simultanées --- c'est Gauss qui érigea l'arithmétique modulaire
en une théorie systématique et rigoureuse.

L'idée fondamentale est simple : deux nombres entiers sont considérés comme
\emph{équivalents} dès lors qu'ils laissent le même reste dans la division par
un entier fixé~$n$. Cette perspective transforme les questions de divisibilité
en un calcul algébrique dans un ensemble fini, ce qui simplifie considérablement
de nombreux problèmes.

%% ── Définition et propriétés élémentaires ───────────────────────
\section{Définition et propriétés élémentaires}
\label{sec:cong-def}

\begin{definition}[Congruence modulo $n$]
\label{def:congruence}
\index{congruence!définition}
Soient $a, b \in \Z$ et $n \in \N$ avec $n \geqslant 1$.
On dit que $a$ est \textbf{congru à $b$ modulo $n$}, et l'on note
\[
  a \equiv b \pmod{n},
\]
lorsque $n \mid (a - b)$, c'est-à-dire lorsqu'il existe $k \in \Z$ tel que
$a - b = kn$.
\end{definition}

\begin{remarque}
De manière équivalente, $a \equiv b \pmod{n}$ si et seulement si $a$ et $b$
ont le même reste dans la division euclidienne par~$n$.
\end{remarque}

\begin{exemple}
\label{ex:cong-basiques}
\leavevmode
\begin{enumerate}[label=(\roman*)]
  \item $17 \equiv 2 \pmod{5}$ car $5 \mid (17 - 2) = 15$.
  \item $-3 \equiv 8 \pmod{11}$ car $11 \mid (-3 - 8) = -11$.
  \item $100 \equiv 0 \pmod{4}$ car $4 \mid 100$.
  \item Pour tout $a \in \Z$, on a $a \equiv a \pmod{n}$.
\end{enumerate}
\end{exemple}

\begin{proposition}[Relation d'équivalence]
\label{prop:cong-equiv}
\index{congruence!relation d'équivalence}
La congruence modulo~$n$ est une relation d'équivalence sur~$\Z$ :
\begin{enumerate}[label=(\roman*)]
  \item \textbf{Réflexivité :} $a \equiv a \pmod{n}$ pour tout $a \in \Z$.
  \item \textbf{Symétrie :} si $a \equiv b \pmod{n}$, alors $b \equiv a \pmod{n}$.
  \item \textbf{Transitivité :} si $a \equiv b \pmod{n}$ et $b \equiv c \pmod{n}$,
    alors $a \equiv c \pmod{n}$.
\end{enumerate}
\end{proposition}

\begin{proof}
\leavevmode
\begin{enumerate}[label=(\roman*)]
  \item On a $a - a = 0 = 0 \cdot n$, donc $n \mid (a - a)$.
  \item Si $n \mid (a - b)$, alors $a - b = kn$ pour un certain $k \in \Z$, d'où
    $b - a = (-k)n$, et donc $n \mid (b - a)$.
  \item Si $n \mid (a - b)$ et $n \mid (b - c)$, alors
    $a - c = (a - b) + (b - c)$ est la somme de deux multiples de~$n$,
    donc $n \mid (a - c)$.
    \qedhere
\end{enumerate}
\end{proof}

Les classes d'équivalence pour cette relation sont appelées \emph{classes de
congruence}\index{classe de congruence} ou \emph{classes de restes}.
La classe de $a$ modulo~$n$ est
\[
  \bar{a} = [a]_n = \{ a + kn : k \in \Z \} = a + n\Z.
\]
Il y a exactement $n$ classes distinctes : $\bar{0}, \bar{1}, \ldots, \overline{n-1}$.

%% ── Arithmétique des congruences ────────────────────────────────
\section{Arithmétique des congruences}
\label{sec:cong-arith}

L'un des atouts majeurs de la congruence est sa \emph{compatibilité} avec les
opérations arithmétiques.

\begin{theoreme}[Compatibilité avec l'addition et la multiplication]
\label{thm:cong-compat}
\index{congruence!compatibilité}
Soient $a, b, c, d \in \Z$ et $n \geqslant 1$. Si $a \equiv b \pmod{n}$ et
$c \equiv d \pmod{n}$, alors :
\begin{enumerate}[label=(\roman*)]
  \item $a + c \equiv b + d \pmod{n}$,
  \item $a - c \equiv b - d \pmod{n}$,
  \item $ac \equiv bd \pmod{n}$.
\end{enumerate}
\end{theoreme}

\begin{proof}
Par hypothèse, $n \mid (a - b)$ et $n \mid (c - d)$.
\begin{enumerate}[label=(\roman*)]
  \item $(a + c) - (b + d) = (a - b) + (c - d)$, somme de deux multiples de~$n$.
  \item $(a - c) - (b - d) = (a - b) - (c - d)$, différence de deux multiples de~$n$.
  \item On écrit $ac - bd = ac - bc + bc - bd = c(a - b) + b(c - d)$.
    Or $n \mid c(a - b)$ et $n \mid b(c - d)$, donc $n \mid (ac - bd)$.
    \qedhere
\end{enumerate}
\end{proof}

\begin{corollaire}[Puissances]
\label{cor:cong-puiss}
\index{congruence!puissances}
Si $a \equiv b \pmod{n}$, alors pour tout $k \in \N$,
\[
  a^k \equiv b^k \pmod{n}.
\]
\end{corollaire}

\begin{proof}
Par récurrence sur~$k$. Le cas $k = 0$ est trivial ($1 \equiv 1$). Si
$a^k \equiv b^k \pmod{n}$, alors en multipliant par $a \equiv b \pmod{n}$
on obtient $a^{k+1} \equiv b^{k+1} \pmod{n}$ par le théorème~\ref{thm:cong-compat}(iii).
\end{proof}

\begin{exemple}
\label{ex:puissance-cong}
Calculons le reste de $3^{100}$ dans la division par~$7$.

On a $3^2 = 9 \equiv 2 \pmod{7}$, puis $3^3 \equiv 3 \cdot 2 = 6 \equiv -1 \pmod{7}$.
Donc $3^6 \equiv (-1)^2 = 1 \pmod{7}$.
Puisque $100 = 6 \times 16 + 4$, on obtient
\[
  3^{100} = (3^6)^{16} \cdot 3^4 \equiv 1^{16} \cdot 3^4 = 81 \equiv 81 - 11 \cdot 7 = 81 - 77 = 4 \pmod{7}.
\]
\end{exemple}

\begin{remarque}[Simplification prudente]
\label{rem:simplif-cong}
\index{congruence!simplification}
Attention : on ne peut pas toujours simplifier une congruence. Si $ca \equiv cb \pmod{n}$,
on ne peut conclure $a \equiv b \pmod{n}$ que si $\pgcd(c, n) = 1$. En effet,
$ca \equiv cb \pmod{n}$ signifie $n \mid c(a - b)$. Si $\pgcd(c, n) = 1$, le
lemme de Gauss donne $n \mid (a - b)$. En général, on a seulement
\[
  ca \equiv cb \pmod{n} \iff a \equiv b \pmod*{\frac{n}{\pgcd(c, n)}}.
\]
\end{remarque}

%% ── L'anneau Z/nZ ──────────────────────────────────────────────
\section{L'anneau $\Z/n\Z$}
\label{sec:ZnZ}
\index{Z/nZ@$\Z/n\Z$}

\begin{definition}[L'anneau $\Z/n\Z$]
\label{def:ZnZ}
L'ensemble quotient
\[
  \Z/n\Z = \{ \bar{0}, \bar{1}, \ldots, \overline{n-1} \}
\]
muni des opérations
\[
  \bar{a} + \bar{b} = \overline{a + b}, \qquad
  \bar{a} \cdot \bar{b} = \overline{a \cdot b}
\]
est un anneau commutatif unitaire\index{anneau!commutatif}. L'élément neutre pour
l'addition est $\bar{0}$ et l'élément neutre pour la multiplication est $\bar{1}$.
\end{definition}

\begin{remarque}
La bonne définition des opérations (c'est-à-dire leur indépendance par rapport
au choix du représentant) est précisément le contenu du
théorème~\ref{thm:cong-compat}.
\end{remarque}

\begin{definition}[Éléments inversibles et diviseurs de zéro]
\label{def:inversibles-ZnZ}
\index{Z/nZ@$\Z/n\Z$!inversibles}
\index{diviseur de zéro}
Soit $\bar{a} \in \Z/n\Z$.
\begin{enumerate}[label=(\roman*)]
  \item $\bar{a}$ est \textbf{inversible} (ou est une \textbf{unité}) s'il
    existe $\bar{b} \in \Z/n\Z$ tel que $\bar{a} \cdot \bar{b} = \bar{1}$.
  \item $\bar{a}$ est un \textbf{diviseur de zéro} si $\bar{a} \neq \bar{0}$
    et s'il existe $\bar{b} \neq \bar{0}$ tel que $\bar{a} \cdot \bar{b} = \bar{0}$.
\end{enumerate}
\end{definition}

\begin{proposition}[Caractérisation des inversibles]
\label{prop:inv-ZnZ}
\index{Z/nZ@$\Z/n\Z$!inversibles}
L'élément $\bar{a}$ est inversible dans $\Z/n\Z$ si et seulement si $\pgcd(a, n) = 1$.
De plus, lorsque $\bar{a}$ est inversible, son inverse peut être calculé à l'aide
de l'identité de Bézout.
\end{proposition}

\begin{proof}
$\bar{a}$ est inversible si et seulement s'il existe $b \in \Z$ tel que
$ab \equiv 1 \pmod{n}$, c'est-à-dire $ab - 1 = kn$ pour un $k \in \Z$, soit
$ab + n(-k) = 1$. Par le théorème de Bézout\index{Bézout, théorème de}, une
telle identité existe si et seulement si $\pgcd(a, n) = 1$.
\end{proof}

\begin{notation}
\label{not:groupe-unites}
Le groupe des éléments inversibles de $\Z/n\Z$ se note
$(\Z/n\Z)^\times$\index{Z/nZ@$(\Z/n\Z)^\times$}. Son cardinal est la
fonction indicatrice d'Euler $\varphi(n)$\index{Euler!fonction indicatrice $\varphi$},
que nous étudierons en détail au chapitre~\ref{ch:fermat-euler-wilson}.
\end{notation}

\begin{corollaire}[$\Z/n\Z$ est un corps si et seulement si $n$ est premier]
\label{cor:ZpZ-corps}
\index{Z/pZ@$\Z/p\Z$!corps}
L'anneau $\Z/n\Z$ est un corps si et seulement si $n$ est un nombre premier.
\end{corollaire}

\begin{proof}
Si $n = p$ est premier, alors pour tout $a$ avec $1 \leqslant a \leqslant p - 1$,
on a $\pgcd(a, p) = 1$, donc $\bar{a}$ est inversible. Réciproquement, si $n$
n'est pas premier, alors $n = ab$ avec $1 < a, b < n$, et $\bar{a} \cdot \bar{b} = \bar{0}$
dans $\Z/n\Z$ avec $\bar{a} \neq \bar{0}$ et $\bar{b} \neq \bar{0}$: l'anneau
possède des diviseurs de zéro et ne peut donc être un corps.
\end{proof}

\begin{exemple}
\label{ex:Z12Z}
Dans $\Z/12\Z$, les inversibles sont les classes $\bar{a}$ avec $\pgcd(a, 12) = 1$,
soit $(\Z/12\Z)^\times = \{\bar{1}, \bar{5}, \bar{7}, \bar{11}\}$,
de cardinal $\varphi(12) = 4$.

Les diviseurs de zéro sont $\bar{2}, \bar{3}, \bar{4}, \bar{6}, \bar{8}, \bar{9}, \bar{10}$.
Par exemple, $\bar{3} \cdot \bar{4} = \overline{12} = \bar{0}$.
\end{exemple}

%% ── Visualisation : arithmétique d'horloge ─────────────────────
\section{Visualisation de l'arithmétique modulaire}
\label{sec:cong-visu}

L'arithmétique modulo~$n$ peut être visualisée comme une \emph{arithmétique
d'horloge}\index{arithmétique d'horloge} : les entiers sont disposés sur un
cercle de $n$ positions, et les opérations « bouclent » après~$n$.

\begin{center}
\begin{tikzpicture}[scale=1.4]
  % Horloge mod 12
  \draw[thick] (0,0) circle (2);
  \foreach \i in {0,...,11} {
    \pgfmathsetmacro{\angle}{90 - \i*30}
    \draw[thick] (\angle:1.85) -- (\angle:2.15);
    \node at (\angle:2.45) {\small $\i$};
  }
  % Flèche : addition de 5 à partir de 3
  \draw[-{Stealth[length=3mm]}, thick, blue!70!black]
    (90-3*30:1.6) arc[start angle={90-3*30}, end angle={90-8*30}, radius=1.6];
  \node[blue!70!black] at (0, -0.3) {$+5$};
  \node[below] at (0,-2.5) {$3 + 5 \equiv 8 \pmod{12}$};
  % Titre
  \node[above] at (0,2.7) {\textbf{Arithmétique mod $12$}};
\end{tikzpicture}
\end{center}

\begin{center}
\begin{tikzpicture}[scale=0.85]
  % Table de multiplication mod 7
  \node[above] at (3.5, 0.8) {\textbf{Table de multiplication dans $\Z/7\Z$}};
  \foreach \i in {1,...,6} {
    \node[font=\small\bfseries] at (\i, 0) {$\i$};
    \node[font=\small\bfseries] at (0, -\i) {$\i$};
  }
  \node[font=\small\bfseries] at (0, 0) {$\times$};
  \draw[thick] (-0.4, -0.5) -- (6.6, -0.5);
  \draw[thick] (0.5, 0.4) -- (0.5, -6.6);
  \foreach \i in {1,...,6} {
    \foreach \j in {1,...,6} {
      \pgfmathsetmacro{\val}{int(Mod(\i*\j, 7))}
      \node[font=\small] at (\j, -\i) {$\val$};
    }
  }
\end{tikzpicture}
\end{center}

\begin{remarque}
Dans la table de multiplication modulo~$7$ (un nombre premier), chaque ligne et
chaque colonne contient une permutation de $\{1, 2, 3, 4, 5, 6\}$. Cette propriété
reflète le fait que tout élément non nul est inversible : $\Z/7\Z$ est un corps.
\end{remarque}

%% ── Congruences linéaires ───────────────────────────────────────
\section{Congruences linéaires}
\label{sec:cong-lin}
\index{congruence!linéaire}

\begin{theoreme}[Résolution de $ax \equiv b \pmod{n}$]
\label{thm:cong-lin}
Soient $a, b \in \Z$ et $n \geqslant 1$. Posons $d = \pgcd(a, n)$.
\begin{enumerate}[label=(\roman*)]
  \item La congruence $ax \equiv b \pmod{n}$ admet une solution si et seulement si
    $d \mid b$.
  \item Lorsque $d \mid b$, l'ensemble des solutions est une unique classe modulo
    $n/d$ : si $x_0$ est une solution particulière, les solutions sont
    \[
      x \equiv x_0 \pmod*{\frac{n}{d}}.
    \]
    Modulo~$n$, il y a donc exactement~$d$ solutions distinctes :
    $x_0, x_0 + \frac{n}{d}, x_0 + 2\cdot\frac{n}{d}, \ldots,
    x_0 + (d-1)\cdot\frac{n}{d}$.
\end{enumerate}
\end{theoreme}

\begin{proof}
L'équation $ax \equiv b \pmod{n}$ équivaut à l'existence de $y \in \Z$ tel que
$ax + ny = b$. Par le théorème de Bézout, cette équation diophantienne admet une
solution si et seulement si $d = \pgcd(a, n)$ divise~$b$.

Supposons $d \mid b$. En divisant par~$d$, on obtient
\[
  \frac{a}{d}\, x \equiv \frac{b}{d} \pmod*{\frac{n}{d}}.
\]
Or $\pgcd\!\left(\frac{a}{d}, \frac{n}{d}\right) = 1$, donc $\frac{a}{d}$ est
inversible modulo $\frac{n}{d}$, et il existe une unique solution
$x_0 \equiv \left(\frac{a}{d}\right)^{-1} \cdot \frac{b}{d} \pmod*{\frac{n}{d}}$.

Modulo~$n$, les valeurs $x_0, x_0 + \frac{n}{d}, \ldots,
x_0 + (d-1)\frac{n}{d}$ sont distinctes et forment l'ensemble complet des solutions.
\end{proof}

\begin{exemple}
\label{ex:cong-lin-1}
Résolvons $6x \equiv 15 \pmod{21}$.

On a $d = \pgcd(6, 21) = 3$ et $3 \mid 15$, donc il y a des solutions.
On divise par~$3$ : $2x \equiv 5 \pmod{7}$. L'inverse de $2$ modulo~$7$ est $4$
(car $2 \times 4 = 8 \equiv 1$). Donc $x \equiv 4 \times 5 = 20 \equiv 6 \pmod{7}$.

Modulo~$21$, les trois solutions sont $x \equiv 6, 13, 20 \pmod{21}$.
\end{exemple}

\begin{exemple}
\label{ex:cong-lin-sans-sol}
La congruence $4x \equiv 3 \pmod{6}$ n'a pas de solution car $\pgcd(4, 6) = 2$
et $2 \nmid 3$.
\end{exemple}

%% ── Systèmes de congruences (aperçu) ───────────────────────────
\section{Systèmes de congruences}
\label{sec:cong-systemes}
\index{congruence!système de}

Il est fréquent de devoir résoudre simultanément plusieurs congruences. Par exemple :
\[
  \begin{cases}
    x \equiv 2 \pmod{3}, \\
    x \equiv 3 \pmod{5}, \\
    x \equiv 2 \pmod{7}.
  \end{cases}
\]

La théorie complète de ces systèmes repose sur le \emph{théorème des restes chinois}
(CRT)\index{théorème des restes chinois}, que nous démontrerons au chapitre suivant.
Pour l'instant, notons le résultat fondamental :

\begin{theoreme}[Théorème des restes chinois --- énoncé préliminaire]
\label{thm:CRT-apercu}
Si $n_1, n_2, \ldots, n_k$ sont deux à deux premiers entre eux et si
$a_1, a_2, \ldots, a_k \in \Z$, alors le système
\[
  x \equiv a_i \pmod{n_i}, \qquad i = 1, \ldots, k,
\]
admet une solution, unique modulo $N = n_1 n_2 \cdots n_k$.
\end{theoreme}

\begin{exemple}
\label{ex:CRT-apercu}
Le système ci-dessus donne $x \equiv 23 \pmod{105}$ (on vérifie : $23 = 7 \times 3 + 2$,
$23 = 4 \times 5 + 3$, $23 = 3 \times 7 + 2$).
\end{exemple}

%% ── Critères de divisibilité ────────────────────────────────────
\section{Critères de divisibilité par congruences}
\label{sec:cong-divis}
\index{divisibilité!critères}

Les critères classiques de divisibilité s'expliquent élégamment par les congruences.
Soit $N = a_k a_{k-1} \cdots a_1 a_0$ un entier écrit en base~$10$, c'est-à-dire
$N = \sum_{i=0}^{k} a_i \cdot 10^i$.

\begin{proposition}[Critère de divisibilité par $3$ et par $9$]
\label{prop:div-3-9}
\index{divisibilité!par 3}\index{divisibilité!par 9}
Un entier $N$ est divisible par~$3$ (resp.\ par~$9$) si et seulement si la somme
de ses chiffres $a_0 + a_1 + \cdots + a_k$ est divisible par~$3$ (resp.\ par~$9$).
\end{proposition}

\begin{proof}
Puisque $10 \equiv 1 \pmod{9}$, on a $10^i \equiv 1 \pmod{9}$ pour tout~$i$, d'où
\[
  N = \sum_{i=0}^{k} a_i \cdot 10^i \equiv \sum_{i=0}^{k} a_i \cdot 1
  = a_0 + a_1 + \cdots + a_k \pmod{9}.
\]
Comme $3 \mid 9$, le même raisonnement vaut modulo~$3$.
\end{proof}

\begin{proposition}[Critère de divisibilité par $11$]
\label{prop:div-11}
\index{divisibilité!par 11}
Un entier $N$ est divisible par~$11$ si et seulement si la somme alternée de ses
chiffres $a_0 - a_1 + a_2 - \cdots + (-1)^k a_k$ est divisible par~$11$.
\end{proposition}

\begin{proof}
On a $10 \equiv -1 \pmod{11}$, donc $10^i \equiv (-1)^i \pmod{11}$, d'où
$N \equiv \sum_{i=0}^{k} (-1)^i a_i \pmod{11}$.
\end{proof}

\begin{proposition}[Critère de divisibilité par $7$]
\label{prop:div-7}
\index{divisibilité!par 7}
Les puissances de $10$ modulo~$7$ suivent le cycle
$10^0 \equiv 1$, $10^1 \equiv 3$, $10^2 \equiv 2$, $10^3 \equiv 6$,
$10^4 \equiv 4$, $10^5 \equiv 5$, $10^6 \equiv 1 \pmod{7}$,
de période~$6$. Ainsi,
\[
  N \equiv a_0 + 3a_1 + 2a_2 + 6a_3 + 4a_4 + 5a_5 + a_6 + \cdots \pmod{7}.
\]
\end{proposition}

\begin{exemple}
\label{ex:div-criteres}
Soit $N = 123\,456\,789$.
\begin{itemize}
  \item Somme des chiffres : $1+2+3+4+5+6+7+8+9 = 45$. Or $9 \mid 45$,
    donc $9 \mid N$ (et a fortiori $3 \mid N$).
  \item Somme alternée : $9 - 8 + 7 - 6 + 5 - 4 + 3 - 2 + 1 = 5$. Or $11 \nmid 5$,
    donc $11 \nmid N$.
\end{itemize}
\end{exemple}

%% ── Exponentiation modulaire rapide ─────────────────────────────
\section{Exponentiation modulaire rapide}
\label{sec:cong-exp-rapide}
\index{exponentiation modulaire}
\index{exponentiation rapide}

Le calcul de $a^k \bmod n$ est une opération fondamentale en théorie des nombres
et en cryptographie. L'algorithme naïf (multiplier $a$ par lui-même $k$ fois)
nécessite $k - 1$ multiplications. L'algorithme d'\textbf{exponentiation rapide}
\index{exponentiation rapide!algorithme}(\emph{binary exponentiation} ou
\emph{square-and-multiply}) réduit ce nombre à $O(\log k)$ multiplications.

\begin{definition}[Algorithme d'exponentiation rapide]
\label{def:exp-rapide}
Pour calculer $a^k \bmod n$ :
\begin{enumerate}
  \item Écrire $k$ en binaire : $k = (b_s b_{s-1} \cdots b_1 b_0)_2$.
  \item Initialiser $r \leftarrow 1$.
  \item Pour $i = s, s-1, \ldots, 0$ :
    \begin{itemize}
      \item $r \leftarrow r^2 \bmod n$ \quad (élever au carré),
      \item si $b_i = 1$, alors $r \leftarrow r \cdot a \bmod n$ \quad (multiplier).
    \end{itemize}
  \item Retourner $r$.
\end{enumerate}
\end{definition}

\begin{exemple}
\label{ex:exp-rapide}
Calculons $3^{45} \bmod 17$. On a $45 = (101101)_2$.

\medskip
\begin{center}
\begin{tabular}{cccc}
  \toprule
  Étape & Bit $b_i$ & Opération & Résultat $r \bmod 17$ \\
  \midrule
  Init & --- & $r = 1$ & $1$ \\
  $i=5$ & $1$ & $1^2 = 1$, puis $1 \cdot 3 = 3$ & $3$ \\
  $i=4$ & $0$ & $3^2 = 9$ & $9$ \\
  $i=3$ & $1$ & $9^2 = 81 \equiv 13$, puis $13 \cdot 3 = 39 \equiv 5$ & $5$ \\
  $i=2$ & $1$ & $5^2 = 25 \equiv 8$, puis $8 \cdot 3 = 24 \equiv 7$ & $7$ \\
  $i=1$ & $0$ & $7^2 = 49 \equiv 15$ & $15$ \\
  $i=0$ & $1$ & $15^2 = 225 \equiv 4$, puis $4 \cdot 3 = 12$ & $12$ \\
  \bottomrule
\end{tabular}
\end{center}
\medskip

Donc $3^{45} \equiv 12 \pmod{17}$, obtenu en seulement $9$ multiplications modulaires
au lieu de~$44$.
\end{exemple}

%% ── Lien avec la cryptographie ──────────────────────────────────
\section{Connexion avec la cryptographie}
\label{sec:cong-crypto}
\index{cryptographie}

L'arithmétique modulaire est la pierre angulaire de la cryptographie moderne.
Le système RSA\index{RSA}, inventé par Rivest, Shamir et Adleman en~1977,
repose entièrement sur les propriétés des congruences et de l'exponentiation
modulaire. L'idée centrale est que le calcul de $a^k \bmod n$ est efficace
(grâce à l'exponentiation rapide), tandis que le problème inverse ---
trouver $k$ connaissant $a$, $a^k \bmod n$ et $n$ --- est supposé
calculatoirement difficile (c'est le problème du \emph{logarithme
discret}\index{logarithme discret}). Nous détaillerons le RSA au
chapitre~\ref{ch:fermat-euler-wilson} après avoir établi les théorèmes
de Fermat et d'Euler.

%% ── Exercices ────────────────────────────────────────────────────
\section{Exercices}
\label{sec:cong-exo}

\begin{exercice}[$\star$]
\label{exo:cong-1}
Déterminer le reste de la division euclidienne de $2^{100}$ par~$7$.
\end{exercice}

\begin{exercice}[$\star$]
\label{exo:cong-2}
Résoudre la congruence $5x \equiv 3 \pmod{13}$.
\end{exercice}

\begin{exercice}[$\star$]
\label{exo:cong-3}
Montrer que $n^3 - n$ est divisible par~$6$ pour tout $n \in \Z$.
\end{exercice}

\begin{exercice}[$\star$]
\label{exo:cong-4}
Calculer les quatre derniers chiffres de $7^{2024}$ (c'est-à-dire $7^{2024} \bmod 10000$).
\end{exercice}

\begin{exercice}[$\star\star$]
\label{exo:cong-5}
Résoudre la congruence $12x \equiv 18 \pmod{30}$. Donner toutes les solutions modulo~$30$.
\end{exercice}

\begin{exercice}[$\star\star$]
\label{exo:cong-6}
\index{Z/nZ@$\Z/n\Z$!structure}
Dresser la table de multiplication de $\Z/8\Z$ et identifier les inversibles
et les diviseurs de zéro.
\end{exercice}

\begin{exercice}[$\star\star$]
\label{exo:cong-7}
Montrer que si $p$ est un nombre premier impair, alors
$1^2 + 2^2 + \cdots + (p-1)^2 \equiv 0 \pmod{p}$.

\emph{Indication : utiliser l'appariement $k \leftrightarrow p - k$.}
\end{exercice}

\begin{exercice}[$\star\star$]
\label{exo:cong-8}
Soit $n \geqslant 2$. Montrer que $\bar{a}$ est un diviseur de zéro dans $\Z/n\Z$
si et seulement si $\pgcd(a, n) > 1$ et $a \not\equiv 0 \pmod{n}$.
\end{exercice}

\begin{exercice}[$\star\star\star$]
\label{exo:cong-9}
Résoudre le système de congruences
\[
  \begin{cases}
    x \equiv 1 \pmod{2}, \\
    x \equiv 2 \pmod{3}, \\
    x \equiv 3 \pmod{5}, \\
    x \equiv 4 \pmod{7}.
  \end{cases}
\]
\end{exercice}

\begin{exercice}[$\star\star\star$]
\label{exo:cong-10}
Montrer que pour tout $n \geqslant 1$, l'anneau $\Z/n\Z$ est un anneau intègre
si et seulement si $n$ est premier.

\emph{Indication : un anneau intègre est un anneau commutatif unitaire sans diviseur
de zéro non nul.}
\end{exercice}

\begin{exercice}[$\star\star\star$]
\label{exo:cong-11}
Soit $p$ un nombre premier. Montrer que les seules solutions de $x^2 \equiv 1 \pmod{p}$
sont $x \equiv \pm 1 \pmod{p}$.
\end{exercice}

%% ── Résumé du chapitre ──────────────────────────────────────────
\section*{Résumé du chapitre}
\addcontentsline{toc}{section}{Résumé du chapitre}
\label{sec:cong-resume}

\begin{enumerate}
  \item $a \equiv b \pmod{n}$ si et seulement si $n \mid (a - b)$. C'est une
    relation d'équivalence compatible avec l'addition et la multiplication.
  \item L'ensemble quotient $\Z/n\Z$ est un anneau commutatif unitaire ; c'est un
    corps si et seulement si $n$ est premier.
  \item L'élément $\bar{a}$ est inversible dans $\Z/n\Z$ si et seulement si
    $\pgcd(a, n) = 1$.
  \item La congruence $ax \equiv b \pmod{n}$ a des solutions si et seulement si
    $\pgcd(a, n) \mid b$, auquel cas il y a exactement $\pgcd(a, n)$ solutions
    modulo~$n$.
  \item Les critères de divisibilité classiques découlent des congruences de~$10$
    modulo $3, 7, 9, 11$.
  \item L'exponentiation rapide permet de calculer $a^k \bmod n$ en $O(\log k)$
    multiplications.
  \item L'arithmétique modulaire est le fondement de la cryptographie moderne (RSA,
    échange de Diffie--Hellman).
\end{enumerate}

\newpage

%%%%%%%%%%%%%%%%%%%%%%%%%%%%%%%%%%%%%%%%%%%%%%%%%%%%%%%%%%%%%%%%%%%%
%%  CHAPITRE 4 — Théorèmes de Fermat, Euler, Wilson
%%%%%%%%%%%%%%%%%%%%%%%%%%%%%%%%%%%%%%%%%%%%%%%%%%%%%%%%%%%%%%%%%%%%
\chapter{Théorèmes de Fermat, Euler et Wilson}
\label{ch:fermat-euler-wilson}
\index{Fermat, petit théorème de}
\index{Euler, théorème d'}
\index{Wilson, théorème de}

\vspace{-0.5em}
\begin{center}\itshape
  « J'ai trouvé un grand nombre de très belles propositions que je vous
  communiquerai une autre fois \ldots{} Tout nombre premier mesure
  infailliblement une des puissances $-1$ de quelque progression que ce soit,
  et l'exposant de la dite puissance est sous-multiple du nombre premier
  donné $-1$ [\ldots]. »\\[2pt]
  --- Pierre de Fermat, lettre à Frénicle de Bessy (18 octobre 1640)
\end{center}
\bigskip

%% ── Introduction historique ─────────────────────────────────────
\section{Introduction historique}
\label{sec:few-hist}

En 1640, dans une lettre adressée à Bernard Frénicle de Bessy\index{Frénicle de Bessy},
Pierre de Fermat\index{Fermat, Pierre de} énonça sans démonstration ce qui allait
devenir l'un des résultats les plus célèbres de la théorie des nombres : si $p$
est premier et si $a$ n'est pas divisible par~$p$, alors $a^{p-1} \equiv 1 \pmod{p}$.
La première démonstration publiée est due à Euler\index{Euler, Leonhard} (1736),
qui, trente ans plus tard, généralisa le résultat à tout module~$n$ en introduisant
sa fameuse fonction $\varphi$.

Quant au théorème de Wilson, il fut conjecturé par Edward Waring\index{Waring, Edward}
en 1770 (l'attribuant à son élève John Wilson\index{Wilson, John}) et démontré pour
la première fois par Lagrange\index{Lagrange, Joseph-Louis} la même année.

%% ── La fonction indicatrice d'Euler ─────────────────────────────
\section{La fonction indicatrice d'Euler $\varphi(n)$}
\label{sec:euler-phi}
\index{Euler!fonction indicatrice $\varphi$}

\begin{definition}[Fonction indicatrice d'Euler]
\label{def:euler-phi}
Pour $n \geqslant 1$, on définit la \textbf{fonction indicatrice d'Euler}\index{fonction indicatrice d'Euler|textbf}
\[
  \varphi(n) = \#\{ a \in \{1, 2, \ldots, n\} : \pgcd(a, n) = 1 \}
  = \abs{(\Z/n\Z)^\times}.
\]
\end{definition}

\begin{exemple}
\label{ex:phi-petits}
\leavevmode
\begin{enumerate}[label=(\roman*)]
  \item $\varphi(1) = 1$.
  \item Si $p$ est premier, $\varphi(p) = p - 1$.
  \item $\varphi(6) = \#\{1, 5\} = 2$.
  \item $\varphi(12) = \#\{1, 5, 7, 11\} = 4$.
  \item $\varphi(8) = \#\{1, 3, 5, 7\} = 4$.
\end{enumerate}
\end{exemple}

\begin{proposition}[Calcul de $\varphi$ pour les puissances de premiers]
\label{prop:phi-prime-power}
\index{Euler!fonction indicatrice $\varphi$!puissances de premiers}
Si $p$ est premier et $k \geqslant 1$, alors
\[
  \varphi(p^k) = p^k - p^{k-1} = p^{k-1}(p - 1).
\]
\end{proposition}

\begin{proof}
Parmi les entiers $1, 2, \ldots, p^k$, ceux qui ne sont \emph{pas} premiers avec
$p^k$ sont exactement les multiples de~$p$, à savoir $p, 2p, 3p, \ldots, p^{k-1} \cdot p = p^k$.
Il y en a $p^{k-1}$. Donc $\varphi(p^k) = p^k - p^{k-1}$.
\end{proof}

\begin{theoreme}[Multiplicativité de $\varphi$]
\label{thm:phi-mult}
\index{Euler!fonction indicatrice $\varphi$!multiplicativité}
Si $\pgcd(m, n) = 1$, alors $\varphi(mn) = \varphi(m)\,\varphi(n)$.
\end{theoreme}

\begin{proof}
Par le théorème des restes chinois (théorème~\ref{thm:CRT-apercu}), l'application
\[
  \Z/mn\Z \longrightarrow \Z/m\Z \times \Z/n\Z, \qquad
  \bar{a} \longmapsto (\bar{a}, \bar{a})
\]
est un isomorphisme d'anneaux lorsque $\pgcd(m, n) = 1$. En passant aux groupes
d'unités, on obtient un isomorphisme
$(\Z/mn\Z)^\times \simeq (\Z/m\Z)^\times \times (\Z/n\Z)^\times$,
d'où $\varphi(mn) = \varphi(m)\,\varphi(n)$.
\end{proof}

\begin{corollaire}[Formule générale pour $\varphi(n)$]
\label{cor:phi-formule}
\index{Euler!fonction indicatrice $\varphi$!formule}
Si $n = p_1^{k_1} p_2^{k_2} \cdots p_r^{k_r}$ est la décomposition en facteurs
premiers, alors
\[
  \varphi(n) = n \prod_{p \mid n} \left(1 - \frac{1}{p}\right)
  = n \cdot \frac{p_1 - 1}{p_1} \cdot \frac{p_2 - 1}{p_2} \cdots
    \frac{p_r - 1}{p_r}.
\]
\end{corollaire}

\begin{proof}
C'est la combinaison de la proposition~\ref{prop:phi-prime-power} et du
théorème~\ref{thm:phi-mult} :
\[
  \varphi(n) = \prod_{i=1}^{r} \varphi(p_i^{k_i})
  = \prod_{i=1}^{r} p_i^{k_i-1}(p_i - 1)
  = \prod_{i=1}^{r} p_i^{k_i} \cdot \frac{p_i - 1}{p_i}
  = n \prod_{p \mid n} \left(1 - \frac{1}{p}\right).
  \qedhere
\]
\end{proof}

\begin{exemple}
\label{ex:phi-calculs}
\leavevmode
\begin{enumerate}[label=(\roman*)]
  \item $\varphi(60) = 60 \cdot \left(1 - \frac{1}{2}\right)\left(1 - \frac{1}{3}\right)
    \left(1 - \frac{1}{5}\right) = 60 \cdot \frac{1}{2} \cdot \frac{2}{3} \cdot
    \frac{4}{5} = 16$.
  \item $\varphi(100) = 100 \cdot \left(1 - \frac{1}{2}\right)\left(1 - \frac{1}{5}\right)
    = 100 \cdot \frac{1}{2} \cdot \frac{4}{5} = 40$.
  \item $\varphi(1001) = \varphi(7 \cdot 11 \cdot 13) = 6 \cdot 10 \cdot 12 = 720$.
\end{enumerate}
\end{exemple}

\begin{proposition}[Identité de Gauss pour $\varphi$]
\label{prop:gauss-phi}
\index{Euler!fonction indicatrice $\varphi$!identité de Gauss}
Pour tout $n \geqslant 1$,
\[
  \sum_{d \mid n} \varphi(d) = n.
\]
\end{proposition}

\begin{proof}
Considérons les fractions $\frac{1}{n}, \frac{2}{n}, \ldots, \frac{n}{n}$.
Pour chaque diviseur $d$ de $n$, les fractions qui se simplifient en une fraction
irréductible de dénominateur~$d$ sont exactement les $\frac{a}{d}$ avec
$1 \leqslant a \leqslant d$ et $\pgcd(a, d) = 1$. Il y en a $\varphi(d)$.
Puisque chacune des $n$ fractions se simplifie en une unique fraction irréductible
dont le dénominateur divise~$n$, on a $\sum_{d \mid n} \varphi(d) = n$.
\end{proof}

%% ── Théorème d'Euler ────────────────────────────────────────────
\section{Le théorème d'Euler}
\label{sec:euler-thm}

\begin{theoreme}[Euler]
\label{thm:euler}
\index{Euler, théorème d'|textbf}
Pour tout $n \geqslant 1$ et tout $a \in \Z$ avec $\pgcd(a, n) = 1$,
\[
  a^{\varphi(n)} \equiv 1 \pmod{n}.
\]
\end{theoreme}

\begin{proof}
Soit $(\Z/n\Z)^\times = \{u_1, u_2, \ldots, u_{\varphi(n)}\}$ le groupe des unités
de $\Z/n\Z$. Puisque $\pgcd(a, n) = 1$, l'élément $\bar{a}$ appartient à
$(\Z/n\Z)^\times$, et la multiplication par~$\bar{a}$ est une bijection de
$(\Z/n\Z)^\times$ sur lui-même. Donc l'ensemble
$\{\bar{a} u_1, \bar{a} u_2, \ldots, \bar{a} u_{\varphi(n)}\}$
est une permutation de $\{u_1, u_2, \ldots, u_{\varphi(n)}\}$.

En prenant le produit de tous ces éléments :
\[
  (\bar{a} u_1)(\bar{a} u_2) \cdots (\bar{a} u_{\varphi(n)})
  = u_1 u_2 \cdots u_{\varphi(n)}.
\]
Le membre de gauche vaut $\bar{a}^{\varphi(n)} \cdot u_1 u_2 \cdots u_{\varphi(n)}$.
En simplifiant par le produit $u_1 u_2 \cdots u_{\varphi(n)}$ (qui est inversible
dans $(\Z/n\Z)^\times$), on obtient
\[
  \bar{a}^{\varphi(n)} = \bar{1},
\]
c'est-à-dire $a^{\varphi(n)} \equiv 1 \pmod{n}$.
\end{proof}

\begin{remarque}
Cette preuve est parfois appelée \emph{preuve par permutation}\index{preuve par permutation},
car elle repose sur le fait que la multiplication par un élément inversible permute
les éléments du groupe. C'est essentiellement l'argument utilisé par Euler lui-même.
\end{remarque}

%% ── Petit théorème de Fermat ────────────────────────────────────
\section{Le petit théorème de Fermat}
\label{sec:fermat-petit}

\begin{corollaire}[Petit théorème de Fermat]
\label{cor:petit-fermat}
\index{Fermat, petit théorème de|textbf}
Soit $p$ un nombre premier. Pour tout $a \in \Z$ avec $p \nmid a$,
\[
  a^{p-1} \equiv 1 \pmod{p}.
\]
De manière équivalente, pour tout $a \in \Z$,
\[
  a^p \equiv a \pmod{p}.
\]
\end{corollaire}

\begin{proof}
Si $p \nmid a$, alors $\pgcd(a, p) = 1$ et $\varphi(p) = p - 1$. Le théorème
d'Euler donne $a^{p-1} \equiv 1 \pmod{p}$, d'où $a^p = a \cdot a^{p-1} \equiv a
\pmod{p}$.

Si $p \mid a$, alors $a \equiv 0 \pmod{p}$, et $a^p \equiv 0 \equiv a \pmod{p}$.
Ainsi, la seconde formulation vaut pour tout $a \in \Z$.
\end{proof}

\begin{exemple}
\label{ex:fermat-calculs}
\leavevmode
\begin{enumerate}[label=(\roman*)]
  \item Calculons $2^{340} \bmod 341$. Si $341$ était premier, on aurait
    $2^{340} \equiv 1 \pmod{341}$. On vérifie que $2^{340} \equiv 1 \pmod{341}$
    (en effet, cela est vrai). Pourtant, $341 = 11 \times 31$ n'est pas premier !
    C'est un \emph{pseudo-premier de Fermat}\index{pseudo-premier de Fermat}
    en base~$2$.
  \item Calculons $7^{222} \bmod 11$. Par Fermat, $7^{10} \equiv 1 \pmod{11}$.
    Or $222 = 10 \times 22 + 2$, donc
    $7^{222} = (7^{10})^{22} \cdot 7^2 \equiv 1 \cdot 49 \equiv 5 \pmod{11}$.
  \item Quel est le dernier chiffre de $3^{1000}$ ? On travaille modulo~$10$ :
    $\varphi(10) = 4$, et $\pgcd(3, 10) = 1$, donc $3^4 \equiv 1 \pmod{10}$.
    Puisque $1000 = 4 \times 250$, on a $3^{1000} \equiv 1 \pmod{10}$.
    Le dernier chiffre est~$1$.
\end{enumerate}
\end{exemple}

%% ── Applications : calcul de puissances et test de primalité ────
\section{Applications du petit théorème de Fermat}
\label{sec:fermat-appli}

\subsection{Calcul de grandes puissances modulaires}
\label{subsec:grandes-puiss}

Le théorème de Fermat (ou d'Euler) permet de réduire l'exposant avant d'appliquer
l'exponentiation rapide.

\begin{exemple}
\label{ex:grande-puissance}
Calculons $3^{10^{100}} \bmod 17$.

Par Fermat, $3^{16} \equiv 1 \pmod{17}$. Il suffit donc de trouver
$10^{100} \bmod 16$. On a $\varphi(16) = 8$, et $\pgcd(10, 16) = 2 \neq 1$,
donc on ne peut pas appliquer Euler directement. Calculons :
$10^2 = 100 \equiv 4 \pmod{16}$, $10^3 \equiv 40 \equiv 8 \pmod{16}$,
$10^4 \equiv 80 \equiv 0 \pmod{16}$. Donc pour $k \geqslant 4$,
$10^k \equiv 0 \pmod{16}$.

Ainsi $10^{100} \equiv 0 \pmod{16}$, et
$3^{10^{100}} \equiv 3^0 = 1 \pmod{17}$.
\end{exemple}

\subsection{Test de primalité de Fermat}
\label{subsec:test-fermat}
\index{test de primalité!de Fermat}

\begin{definition}[Test de Fermat]
\label{def:test-fermat}
Soit $n > 1$ un entier impair. Si l'on trouve $a$ avec $1 < a < n$ et
$\pgcd(a, n) = 1$ tel que $a^{n-1} \not\equiv 1 \pmod{n}$, alors $n$ est
\emph{certainement} composé. On dit que $a$ est un \textbf{témoin de Fermat}
\index{témoin de Fermat} pour la non-primalité de~$n$.
\end{definition}

\begin{remarque}
Le test de Fermat est un test probabiliste : il peut déclarer un nombre composé
comme « probablement premier ». Les nombres composés $n$ tels que $a^{n-1} \equiv 1
\pmod{n}$ pour tout $a$ premier avec~$n$ sont appelés \emph{nombres de
Carmichael}\index{Carmichael, nombres de} (par exemple, $561 = 3 \times 11 \times 17$).
Des tests plus puissants, comme le test de Miller--Rabin\index{Miller--Rabin, test de},
surmontent cette limitation.
\end{remarque}

%% ── Théorème de Wilson ──────────────────────────────────────────
\section{Le théorème de Wilson}
\label{sec:wilson}

\begin{theoreme}[Wilson]
\label{thm:wilson}
\index{Wilson, théorème de|textbf}
Un entier $n \geqslant 2$ est premier si et seulement si
\[
  (n - 1)! \equiv -1 \pmod{n}.
\]
\end{theoreme}

\begin{proof}
\textbf{Sens direct.} Supposons $p$ premier. Nous travaillons dans le corps
$\Z/p\Z$. Pour $a \in \{1, 2, \ldots, p-1\}$, l'élément $\bar{a}$ est inversible
dans $(\Z/p\Z)^\times$. Nous cherchons les éléments qui sont leur propre inverse,
c'est-à-dire les solutions de $x^2 \equiv 1 \pmod{p}$, soit $p \mid (x-1)(x+1)$.
Puisque $p$ est premier, $p \mid (x - 1)$ ou $p \mid (x + 1)$, ce qui donne
$x \equiv 1$ ou $x \equiv -1 \pmod{p}$.

Ainsi, dans le produit $(p-1)! = 1 \cdot 2 \cdots (p-1)$, les éléments $\bar{1}$
et $\overline{p-1} = \overline{-1}$ sont les seuls à être leur propre inverse. Les
$p - 3$ autres éléments de $\{2, 3, \ldots, p-2\}$ se regroupent en
$\frac{p-3}{2}$ paires $(a, a^{-1})$ avec $a \neq a^{-1}$. Le produit de chaque
paire vaut~$1$, donc
\[
  (p-1)! \equiv 1 \cdot \underbrace{(2 \cdot 2^{-1})(3 \cdot 3^{-1}) \cdots}_{\text{produit} = 1}
  \cdot (p-1) \equiv 1 \cdot 1 \cdot (p-1) \equiv -1 \pmod{p}.
\]

\textbf{Réciproque.} Supposons $n$ composé. Si $n = ab$ avec $1 < a, b < n$ et
$a \neq b$, alors $a$ et $b$ apparaissent tous deux comme facteurs dans $(n-1)!$,
donc $n = ab \mid (n-1)!$, d'où $(n-1)! \equiv 0 \pmod{n}$, et en particulier
$(n-1)! \not\equiv -1 \pmod{n}$.

Le cas $n = p^2$ avec $p$ premier nécessite un argument séparé. Si $p \geqslant 3$,
alors $p$ et $2p$ sont deux facteurs distincts dans $\{1, \ldots, p^2 - 1\}$,
donc $p^2 \mid (p^2 - 1)!$. Si $n = 4 = 2^2$, on vérifie directement :
$3! = 6 \equiv 2 \not\equiv -1 \pmod{4}$.
\end{proof}

\begin{exemple}
\label{ex:wilson}
\leavevmode
\begin{enumerate}[label=(\roman*)]
  \item $4! = 24 \equiv 4 \equiv -1 \pmod{5}$. Confirmé : $5$ est premier.
  \item $6! = 720 \equiv 720 - 102 \times 7 = 720 - 714 = 6 \equiv -1 \pmod{7}$.
    Confirmé.
  \item $9! = 362880 \equiv 0 \pmod{10}$. Or $-1 \equiv 9 \pmod{10}$, donc
    $9! \not\equiv -1 \pmod{10}$. En effet, $10$ n'est pas premier.
\end{enumerate}
\end{exemple}

\begin{remarque}
Le théorème de Wilson est théoriquement élégant mais pratiquement inutilisable
comme test de primalité, car le calcul de $(n-1)!$ est de complexité exponentielle
en la taille de~$n$.
\end{remarque}

%% ── Racines primitives ──────────────────────────────────────────
\section{Racines primitives}
\label{sec:racines-prim}
\index{racine primitive}

\subsection{Ordre d'un élément}
\label{subsec:ordre-elt}

\begin{definition}[Ordre multiplicatif]
\label{def:ordre-mult}
\index{ordre multiplicatif}
Soit $n \geqslant 1$ et $a \in \Z$ avec $\pgcd(a, n) = 1$. L'\textbf{ordre
multiplicatif}\index{ordre multiplicatif|textbf} de $a$ modulo~$n$, noté
$\ord_n(a)$\index{$\ord_n(a)$}, est le plus petit entier $d \geqslant 1$ tel que
$a^d \equiv 1 \pmod{n}$.
\end{definition}

\begin{proposition}[Propriétés de l'ordre]
\label{prop:ordre-props}
Soient $a \in \Z$ avec $\pgcd(a, n) = 1$ et $d = \ord_n(a)$.
\begin{enumerate}[label=(\roman*)]
  \item $a^k \equiv 1 \pmod{n}$ si et seulement si $d \mid k$.
  \item En particulier, $d \mid \varphi(n)$.
  \item Les éléments $a^0, a^1, \ldots, a^{d-1}$ sont distincts modulo~$n$.
\end{enumerate}
\end{proposition}

\begin{proof}
\leavevmode
\begin{enumerate}[label=(\roman*)]
  \item Si $d \mid k$, alors $k = dq$ et $a^k = (a^d)^q \equiv 1 \pmod{n}$.
    Réciproquement, écrivons $k = dq + r$ avec $0 \leqslant r < d$. Alors
    $a^r = a^k \cdot (a^d)^{-q} \equiv 1 \pmod{n}$. Par minimalité de~$d$,
    on a $r = 0$, donc $d \mid k$.
  \item Par Euler, $a^{\varphi(n)} \equiv 1 \pmod{n}$, donc $d \mid \varphi(n)$
    par~(i).
  \item Si $a^i \equiv a^j \pmod{n}$ avec $0 \leqslant i < j < d$, alors
    $a^{j-i} \equiv 1 \pmod{n}$ avec $0 < j - i < d$, ce qui contredit la
    minimalité de~$d$.
    \qedhere
\end{enumerate}
\end{proof}

\subsection{Définition et existence des racines primitives}
\label{subsec:rac-prim-def}

\begin{definition}[Racine primitive]
\label{def:racine-prim}
\index{racine primitive!définition}
Un entier $g$ est une \textbf{racine primitive modulo~$n$} si $\pgcd(g, n) = 1$ et
$\ord_n(g) = \varphi(n)$, c'est-à-dire si $g$ engendre le groupe
$(\Z/n\Z)^\times$ :
\[
  (\Z/n\Z)^\times = \{ \bar{g}^0, \bar{g}^1, \ldots, \bar{g}^{\varphi(n)-1} \}.
\]
\end{definition}

\begin{theoreme}[Existence des racines primitives pour les premiers]
\label{thm:rac-prim-existe}
\index{racine primitive!existence}
Pour tout nombre premier~$p$, le groupe $(\Z/p\Z)^\times$ est cyclique
d'ordre $p - 1$. Autrement dit, il existe une racine primitive modulo~$p$.
\end{theoreme}

\begin{proof}[Esquisse de la preuve]
Pour chaque diviseur~$d$ de $p - 1$, notons $\psi(d)$ le nombre d'éléments
d'ordre exactement~$d$ dans $(\Z/p\Z)^\times$. Nous allons montrer que
$\psi(d) = \varphi(d)$ pour tout $d \mid (p-1)$, ce qui en particulier donne
$\psi(p-1) = \varphi(p-1) \geqslant 1$.

\emph{Étape 1.} Si $(\Z/p\Z)^\times$ contient un élément $a$ d'ordre~$d$, alors
$a, a^2, \ldots, a^d$ sont $d$ racines distinctes du polynôme $X^d - 1$ dans le
corps $\Z/p\Z$. Puisqu'un polynôme de degré~$d$ sur un corps a au plus $d$ racines,
ce sont les \emph{seules} racines. Tout élément d'ordre~$d$ est une puissance
$a^k$ avec $\pgcd(k, d) = 1$, et il y en a $\varphi(d)$. Donc
$\psi(d) \in \{0, \varphi(d)\}$.

\emph{Étape 2.} Chaque élément de $(\Z/p\Z)^\times$ a un ordre qui divise $p - 1$.
Donc $\sum_{d \mid (p-1)} \psi(d) = p - 1$. Or, par l'identité de
Gauss (proposition~\ref{prop:gauss-phi}), $\sum_{d \mid (p-1)} \varphi(d) = p - 1$.
Puisque $\psi(d) \leqslant \varphi(d)$ pour tout~$d$ et que les sommes sont égales,
on conclut $\psi(d) = \varphi(d)$ pour tout $d \mid (p-1)$.

En particulier, $\psi(p-1) = \varphi(p-1) \geqslant 1$, ce qui prouve l'existence
d'une racine primitive modulo~$p$.
\end{proof}

\begin{remarque}[Existence générale]
\label{rem:rac-prim-generale}
\index{racine primitive!existence générale}
Plus généralement, des racines primitives existent modulo~$n$ si et seulement si
$n \in \{1, 2, 4, p^k, 2p^k\}$ où $p$ est un premier impair et $k \geqslant 1$.
Nous admettrons ce résultat.
\end{remarque}

\subsection{Dénombrement des racines primitives}
\label{subsec:comptage-rac-prim}

\begin{corollaire}[Nombre de racines primitives]
\label{cor:nb-rac-prim}
\index{racine primitive!nombre de}
Si $p$ est premier, le nombre de racines primitives modulo~$p$ est
$\varphi(p - 1)$.
\end{corollaire}

\begin{proof}
C'est exactement $\psi(p-1) = \varphi(p-1)$ établi dans la preuve ci-dessus.
\end{proof}

\begin{exemple}
\label{ex:rac-prim-7}
Déterminons les racines primitives modulo~$7$. On a $\varphi(6) = 2$, donc il y a
$2$ racines primitives. Les ordres possibles divisent $6$ : $\{1, 2, 3, 6\}$.

\medskip
\begin{center}
\begin{tabular}{c|cccccc}
  $a$ & $a^1$ & $a^2$ & $a^3$ & $a^4$ & $a^5$ & $a^6$ \\
  \hline
  $1$ & $1$ & & & & & \\
  $2$ & $2$ & $4$ & $1$ & & & \\
  $3$ & $3$ & $2$ & $6$ & $4$ & $5$ & $1$ \\
  $4$ & $4$ & $2$ & $1$ & & & \\
  $5$ & $5$ & $4$ & $6$ & $2$ & $3$ & $1$ \\
  $6$ & $6$ & $1$ & & & & \\
\end{tabular}
\end{center}
\medskip

Les ordres sont : $\ord_7(1)=1$, $\ord_7(2)=3$, $\ord_7(3)=6$, $\ord_7(4)=3$,
$\ord_7(5)=6$, $\ord_7(6)=2$. Les racines primitives sont $3$ et $5$.
\end{exemple}

%% ── Visualisation des groupes multiplicatifs ────────────────────
\subsection{Visualisation du groupe multiplicatif $(\Z/p\Z)^\times$}
\label{subsec:visu-groupe-mult}

\begin{center}
\begin{tikzpicture}[scale=1.3]
  % Puissances de 3 mod 7 sur un cercle
  \def\n{6}
  \node[above] at (0,2.7) {\textbf{Puissances de $g = 3$ dans $(\Z/7\Z)^\times$}};
  \draw[gray!40, thick] (0,0) circle (2);
  % Place elements at positions corresponding to powers
  \foreach \k/\val in {0/1, 1/3, 2/2, 3/6, 4/4, 5/5} {
    \pgfmathsetmacro{\angle}{90 - \k*60}
    \node[circle, draw, fill=blue!15, inner sep=3pt, font=\small\bfseries] (n\k) at (\angle:2) {$\val$};
    \node[font=\tiny, gray] at (\angle:2.5) {$3^{\k}$};
  }
  % Draw arrows for multiplication by 3
  \foreach \k [evaluate=\k as \knext using {int(mod(\k+1,6))}] in {0,...,5} {
    \draw[-{Stealth[length=2mm]}, thick, blue!60!black]
      (n\k) -- (n\knext);
  }
\end{tikzpicture}
\end{center}

\begin{center}
\begin{tikzpicture}[scale=1.3]
  % Puissances de 2 mod 7 (sous-groupe d'ordre 3)
  \def\n{3}
  \node[above] at (0,2.7) {\textbf{Puissances de $a = 2$ dans $(\Z/7\Z)^\times$ (ordre 3)}};
  \draw[gray!40, thick] (0,0) circle (2);
  \foreach \k/\val in {0/1, 1/2, 2/4} {
    \pgfmathsetmacro{\angle}{90 - \k*120}
    \node[circle, draw, fill=red!15, inner sep=3pt, font=\small\bfseries] (m\k) at (\angle:2) {$\val$};
    \node[font=\tiny, gray] at (\angle:2.5) {$2^{\k}$};
  }
  \foreach \k [evaluate=\k as \knext using {int(mod(\k+1,3))}] in {0,1,2} {
    \draw[-{Stealth[length=2mm]}, thick, red!60!black]
      (m\k) -- (m\knext);
  }
\end{tikzpicture}
\end{center}

%% ── Le cryptosystème RSA ────────────────────────────────────────
\section{Le cryptosystème RSA}
\label{sec:RSA}
\index{RSA|textbf}
\index{cryptographie!RSA}

Le système RSA\index{RSA!description}, inventé par Rivest, Shamir et
Adleman\index{Rivest}\index{Shamir}\index{Adleman} en~1977, est l'un des
cryptosystèmes à clé publique les plus utilisés au monde. Sa sécurité repose sur
la difficulté de factoriser de grands entiers, et sa correction sur le théorème
d'Euler.

\subsection{Génération des clés}
\label{subsec:RSA-cles}
\index{RSA!génération des clés}

\begin{enumerate}
  \item Choisir deux grands nombres premiers distincts $p$ et $q$ (typiquement
    de 1024 bits chacun).
  \item Calculer $n = pq$ et $\varphi(n) = (p-1)(q-1)$.
  \item Choisir $e$ avec $1 < e < \varphi(n)$ et $\pgcd(e, \varphi(n)) = 1$.
    Le choix classique est $e = 65537 = 2^{16} + 1$.
  \item Calculer $d \equiv e^{-1} \pmod{\varphi(n)}$ (par l'algorithme d'Euclide
    étendu).
  \item La \textbf{clé publique}\index{RSA!clé publique} est $(n, e)$.
    La \textbf{clé privée}\index{RSA!clé privée} est $(n, d)$.
\end{enumerate}

\subsection{Chiffrement et déchiffrement}
\label{subsec:RSA-chiffrement}
\index{RSA!chiffrement}

\begin{definition}[Chiffrement RSA]
\label{def:RSA-operations}
Soit $m \in \{0, 1, \ldots, n-1\}$ le message en clair (représenté comme un
entier).
\begin{itemize}
  \item \textbf{Chiffrement :} $c \equiv m^e \pmod{n}$.
  \item \textbf{Déchiffrement :} $m \equiv c^d \pmod{n}$.
\end{itemize}
\end{definition}

\subsection{Preuve de correction}
\label{subsec:RSA-preuve}
\index{RSA!preuve de correction}

\begin{theoreme}[Correction du RSA]
\label{thm:RSA-correct}
Avec les notations ci-dessus, pour tout $m \in \{0, 1, \ldots, n-1\}$,
\[
  (m^e)^d \equiv m \pmod{n}.
\]
\end{theoreme}

\begin{proof}
Puisque $ed \equiv 1 \pmod{\varphi(n)}$, on a $ed = 1 + k\varphi(n)$ pour un
certain $k \in \N$.

\textbf{Cas 1 :} $\pgcd(m, n) = 1$. Par le théorème d'Euler,
$m^{\varphi(n)} \equiv 1 \pmod{n}$, d'où
\[
  m^{ed} = m^{1 + k\varphi(n)} = m \cdot (m^{\varphi(n)})^k \equiv m \cdot 1^k
  = m \pmod{n}.
\]

\textbf{Cas 2 :} $\pgcd(m, n) > 1$. Puisque $n = pq$ et $m < n$, on a soit
$p \mid m$ (mais $q \nmid m$), soit $q \mid m$ (mais $p \nmid m$). Traitons le
premier cas (l'autre est symétrique).

Si $p \mid m$, alors $m^{ed} \equiv 0 \equiv m \pmod{p}$.
D'autre part, $\pgcd(m, q) = 1$, et $\varphi(n) = (p-1)(q-1)$, donc
$(q-1) \mid \varphi(n)$. Par Fermat, $m^{q-1} \equiv 1 \pmod{q}$, d'où
\[
  m^{ed} = m \cdot (m^{\varphi(n)})^k = m \cdot ((m^{q-1})^{(p-1)})^k
  \equiv m \cdot 1 = m \pmod{q}.
\]
Puisque $p \mid (m^{ed} - m)$ et $q \mid (m^{ed} - m)$, et que $\pgcd(p, q) = 1$,
on conclut $pq = n \mid (m^{ed} - m)$.
\end{proof}

\subsection{Exemple numérique complet}
\label{subsec:RSA-exemple}

\begin{exemple}
\label{ex:RSA}
\index{RSA!exemple}
Prenons $p = 61$ et $q = 53$.
\begin{enumerate}
  \item $n = 61 \times 53 = 3233$ et $\varphi(n) = 60 \times 52 = 3120$.
  \item Choisissons $e = 17$. Vérifions : $\pgcd(17, 3120) = 1$. \checkmark
  \item Calculons $d \equiv 17^{-1} \pmod{3120}$.
    Par l'algorithme d'Euclide étendu : $17 \times 2753 = 46801 = 15 \times 3120 + 1$,
    donc $d = 2753$.
  \item Clé publique : $(3233, 17)$. Clé privée : $(3233, 2753)$.
  \item Chiffrement de $m = 65$ : $c = 65^{17} \bmod 3233$.

    Par exponentiation rapide : $17 = (10001)_2$.
    \begin{align*}
      65^1 &\equiv 65 \pmod{3233} \\
      65^2 &\equiv 4225 \equiv 992 \pmod{3233} \\
      65^4 &\equiv 992^2 = 984064 \equiv 984064 - 304 \times 3233 = 1 + 2432 \equiv 2433 \pmod{3233}\\
      65^8 &\equiv 2433^2 = 5919489 \equiv 5919489 - 1831 \times 3233 = 5919489 - 5921723\ldots
    \end{align*}
    En poursuivant le calcul (ou par ordinateur), on trouve $c = 2790$.
  \item Déchiffrement : $m = 2790^{2753} \bmod 3233 = 65$. \checkmark
\end{enumerate}
\end{exemple}

\subsection{Considérations pratiques}
\label{subsec:RSA-pratique}
\index{RSA!considérations pratiques}

\begin{remarque}
En pratique, le RSA nécessite :
\begin{itemize}
  \item Des nombres premiers de grande taille (au moins $2048$ bits pour $n$,
    soit environ $617$ chiffres décimaux) pour résister aux algorithmes de
    factorisation connus.
  \item La génération de grands nombres premiers, réalisée en pratique par des
    tests de primalité probabilistes (Miller--Rabin\index{Miller--Rabin, test de}).
  \item Un \emph{padding} cryptographique (par exemple OAEP\index{OAEP})
    pour éviter certaines attaques.
  \item La destruction sécurisée de $p$, $q$ et $\varphi(n)$ après la génération
    des clés.
\end{itemize}
\end{remarque}

%% ── Échange de clés de Diffie–Hellman ───────────────────────────
\section{L'échange de clés de Diffie--Hellman}
\label{sec:diffie-hellman}
\index{Diffie--Hellman}

Le protocole de Diffie--Hellman\index{Diffie--Hellman|textbf} (1976) permet à deux
parties, Alice et Bob, d'établir un secret partagé sur un canal non sécurisé.

\begin{enumerate}
  \item Alice et Bob choisissent publiquement un grand premier~$p$ et une racine
    primitive $g$ modulo~$p$.
  \item Alice choisit un secret $a \in \{2, \ldots, p-2\}$ et envoie
    $A = g^a \bmod p$ à Bob.
  \item Bob choisit un secret $b \in \{2, \ldots, p-2\}$ et envoie
    $B = g^b \bmod p$ à Alice.
  \item Alice calcule $s = B^a \bmod p = g^{ab} \bmod p$.
  \item Bob calcule $s = A^b \bmod p = g^{ab} \bmod p$.
\end{enumerate}

Les deux parties partagent maintenant le secret $s = g^{ab} \bmod p$. Un
espion qui observe $g$, $p$, $A = g^a \bmod p$ et $B = g^b \bmod p$
devrait résoudre le problème du logarithme discret\index{logarithme discret}
pour retrouver $a$ ou $b$, ce qui est supposé calculatoirement difficile.

\begin{exemple}
\label{ex:DH}
Avec $p = 23$ et $g = 5$ (racine primitive modulo~$23$):
\begin{itemize}
  \item Alice choisit $a = 6$, envoie $A = 5^6 \bmod 23 = 15625 \bmod 23 = 8$.
  \item Bob choisit $b = 15$, envoie $B = 5^{15} \bmod 23 = 19$.
  \item Alice calcule $s = 19^6 \bmod 23 = 2$.
  \item Bob calcule $s = 8^{15} \bmod 23 = 2$.
  \item Le secret partagé est $s = 2$.
\end{itemize}
\end{exemple}

%% ── Exercices ────────────────────────────────────────────────────
\section{Exercices}
\label{sec:few-exo}

\begin{exercice}[$\star$]
\label{exo:few-1}
Calculer $\varphi(n)$ pour $n = 36, 72, 100, 250, 360$.
\end{exercice}

\begin{exercice}[$\star$]
\label{exo:few-2}
En utilisant le petit théorème de Fermat, calculer $2^{1000} \bmod 13$.
\end{exercice}

\begin{exercice}[$\star$]
\label{exo:few-3}
Déterminer le dernier chiffre de $7^{9999}$.
\end{exercice}

\begin{exercice}[$\star$]
\label{exo:few-4}
Vérifier le théorème de Wilson pour $p = 11$ en calculant $10! \bmod 11$.
\end{exercice}

\begin{exercice}[$\star\star$]
\label{exo:few-5}
Trouver toutes les racines primitives modulo~$11$.
\end{exercice}

\begin{exercice}[$\star\star$]
\label{exo:few-6}
\index{Euler, théorème d'!application}
En utilisant le théorème d'Euler, calculer $3^{340} \bmod 341$.
En déduire que $341$ n'est pas nécessairement premier.
\end{exercice}

\begin{exercice}[$\star\star$]
\label{exo:few-7}
Soit $p$ premier. Montrer que si $g$ est une racine primitive modulo~$p$, alors
$g^k$ est aussi une racine primitive si et seulement si $\pgcd(k, p-1) = 1$.
\end{exercice}

\begin{exercice}[$\star\star$]
\label{exo:few-8}
Montrer que pour tout premier $p \geqslant 5$, on a
$1^{p-2} + 2^{p-2} + \cdots + (p-1)^{p-2} \equiv 0 \pmod{p}$.

\emph{Indication : utiliser le petit théorème de Fermat et la somme des inverses
dans $(\Z/p\Z)^\times$.}
\end{exercice}

\begin{exercice}[$\star\star$]
\label{exo:few-9}
\index{Euler!fonction indicatrice $\varphi$!propriétés}
Montrer que $\varphi(n)$ est pair pour tout $n \geqslant 3$.
\end{exercice}

\begin{exercice}[$\star\star\star$]
\label{exo:few-10}
\index{RSA!exercice}
Alice utilise RSA avec $n = 2537$, $e = 13$. Elle reçoit le message chiffré $c = 2081$.
Factoriser $n$, calculer $d$ et déchiffrer le message.

\emph{Indication : $2537 = 43 \times 59$.}
\end{exercice}

\begin{exercice}[$\star\star\star$]
\label{exo:few-11}
\index{Wilson, théorème de!application}
En utilisant le théorème de Wilson, montrer que pour tout premier
$p \equiv 1 \pmod{4}$, l'entier $\left(\frac{p-1}{2}\right)!$
satisfait $\left[\left(\frac{p-1}{2}\right)!\right]^2 \equiv -1 \pmod{p}$.

\emph{Indication : dans le produit $(p-1)! = 1 \cdot 2 \cdots (p-1)$, apparier
chaque $k > \frac{p-1}{2}$ avec $p - k$.}
\end{exercice}

\begin{exercice}[$\star\star\star$]
\label{exo:few-12}
Soit $p$ un premier impair et $g$ une racine primitive modulo~$p$. Montrer que
$g^{(p-1)/2} \equiv -1 \pmod{p}$.
\end{exercice}

\begin{exercice}[$\star\star\star$]
\label{exo:few-13}
\index{Carmichael, nombres de}
Montrer que $561 = 3 \times 11 \times 17$ est un nombre de Carmichael, c'est-à-dire
que $a^{560} \equiv 1 \pmod{561}$ pour tout $a$ avec $\pgcd(a, 561) = 1$.

\emph{Indication : vérifier que $560$ est divisible par $\varphi(3) = 2$,
$\varphi(11) = 10$ et $\varphi(17) = 16$, puis appliquer Fermat modulo chaque
facteur premier et conclure par le CRT.}
\end{exercice}

%% ── Résumé du chapitre ──────────────────────────────────────────
\section*{Résumé du chapitre}
\addcontentsline{toc}{section}{Résumé du chapitre}
\label{sec:few-resume}

\begin{enumerate}
  \item La fonction indicatrice d'Euler $\varphi(n)$ compte les entiers de $\{1, \ldots, n\}$
    premiers avec~$n$. Elle est multiplicative et se calcule par
    $\varphi(n) = n \prod_{p \mid n} (1 - 1/p)$.
  \item \textbf{Théorème d'Euler :} si $\pgcd(a, n) = 1$, alors
    $a^{\varphi(n)} \equiv 1 \pmod{n}$.
  \item \textbf{Petit théorème de Fermat :} si $p$ est premier et $p \nmid a$,
    alors $a^{p-1} \equiv 1 \pmod{p}$.
  \item \textbf{Théorème de Wilson :} $n$ est premier si et seulement si
    $(n-1)! \equiv -1 \pmod{n}$.
  \item Une \textbf{racine primitive} modulo~$n$ est un générateur du groupe
    $(\Z/n\Z)^\times$. Le groupe $(\Z/p\Z)^\times$ est cyclique pour tout
    premier~$p$, et possède $\varphi(p-1)$ racines primitives.
  \item Le \textbf{cryptosystème RSA} repose sur le théorème d'Euler : on chiffre
    par $c = m^e \bmod n$ et on déchiffre par $m = c^d \bmod n$, où
    $ed \equiv 1 \pmod{\varphi(n)}$.
  \item L'\textbf{échange de Diffie--Hellman} permet d'établir un secret partagé
    en exploitant la difficulté du logarithme discret dans $(\Z/p\Z)^\times$.
\end{enumerate}

% === FIN DU CHUNK ch3-4 ===% === DÉBUT DU CHUNK ch5-6 ===

%%%%%%%%%%%%%%%%%%%%%%%%%%%%%%%%%%%%%%%%%%%%%%%%%%%%%%%%%%%%%%%%%%%%
%  CHAPITRE 5 — RESTES CHINOIS ET APPLICATIONS
%%%%%%%%%%%%%%%%%%%%%%%%%%%%%%%%%%%%%%%%%%%%%%%%%%%%%%%%%%%%%%%%%%%%
\chapter{Restes Chinois et Applications}
\label{chap:crt}

\section*{Perspective historique}
\addcontentsline{toc}{section}{Perspective historique}

Le \emph{théorème des restes chinois}\index{théorème des restes chinois} constitue l'un des
résultats les plus anciens de la théorie des nombres. Sa première formulation connue remonte
au traité \emph{Sunzi Suanjing}\index{Sunzi Suanjing} (<<~Manuel de calcul du Maître Sun~>>),
rédigé entre le \textsc{iii}\textsuperscript{e} et le \textsc{v}\textsuperscript{e}~siècle de
notre ère. Le problème fondateur est le suivant~:

\begin{quote}
\itshape
<<~Un nombre donne un reste de $2$ quand on le divise par $3$, un reste de $3$ quand
on le divise par $5$, et un reste de $2$ quand on le divise par $7$. Quel est ce nombre~?~>>
\end{quote}

Au \textsc{xiii}\textsuperscript{e}~siècle, le mathématicien chinois
\textbf{Qin Jiushao}\index{Qin Jiushao} (1202--1261) formula une méthode
algorithmique générale dans son ouvrage \emph{Shushu Jiuzhang} (1247),
anticipant de plusieurs siècles les travaux européens. En Occident,
le résultat fut redécouvert et généralisé par Euler, puis placé dans
un cadre algébrique par Gauss dans ses \emph{Disquisitiones Arithmeticae} (1801).

\medskip

Dans ce chapitre, nous énonçons et démontrons le théorème des restes chinois
sous sa forme classique et algébrique, puis nous en présentons de nombreuses
applications, allant du calcul calendaire à la cryptographie RSA.

%--------------------------------------------------------------------
\section{Systèmes de congruences simultanées}
\label{sec:crt:systemes}
%--------------------------------------------------------------------

\begin{definition}[Système de congruences]
\label{def:systeme-congruences}
\index{système de congruences}
Un \emph{système de congruences simultanées} est un ensemble d'équations
de la forme
\[
  \begin{cases}
    x \equiv a_1 \pmod{m_1} \\
    x \equiv a_2 \pmod{m_2} \\
    \quad \vdots \\
    x \equiv a_k \pmod{m_k}
  \end{cases}
\]
où les $a_i \in \Z$ sont les \emph{restes} prescrits et les $m_i \geq 2$
sont les \emph{modules}.
\end{definition}

\begin{remarque}
Sans hypothèse supplémentaire sur les modules, un tel système peut n'avoir
aucune solution. Par exemple, le système $x \equiv 0 \pmod{2}$,
$x \equiv 1 \pmod{4}$ est incompatible, car la première congruence impose
que $x$ est pair, tandis que la seconde impose $x \equiv 1 \pmod{4}$,
donc $x$ impair.
\end{remarque}

%--------------------------------------------------------------------
\section{Le théorème des restes chinois}
\label{sec:crt:enonce}
%--------------------------------------------------------------------

\begin{theoreme}[Théorème des restes chinois --- CRT]
\label{thm:crt}
\index{théorème des restes chinois!énoncé}
Soient $m_1, m_2, \ldots, m_k$ des entiers $\geq 2$ deux à deux premiers
entre eux, c'est-à-dire $\pgcd(m_i, m_j) = 1$ pour tout $i \neq j$.
Posons $M = m_1 m_2 \cdots m_k$. Alors, pour tout choix d'entiers
$a_1, a_2, \ldots, a_k$, le système
\[
  \begin{cases}
    x \equiv a_1 \pmod{m_1} \\
    x \equiv a_2 \pmod{m_2} \\
    \quad \vdots \\
    x \equiv a_k \pmod{m_k}
  \end{cases}
\]
admet une solution, et cette solution est \emph{unique modulo $M$}.
Autrement dit, l'ensemble des solutions est exactement une classe de
congruence modulo $M$.
\end{theoreme}

\begin{proof}
\index{théorème des restes chinois!preuve}
La preuve procède en deux étapes~: construction d'une solution, puis unicité.

\medskip
\noindent\textbf{Existence (construction explicite).}
Pour chaque $i \in \{1, \ldots, k\}$, posons
\[
  M_i = \frac{M}{m_i} = \prod_{j \neq i} m_j.
\]
Par hypothèse, $\pgcd(M_i, m_i) = 1$ puisque $m_i$ est premier avec chacun
des $m_j$ ($j \neq i$), donc premier avec leur produit. Il existe donc,
par le théorème de Bézout\index{Bézout!théorème de}, un entier $N_i$ tel que
\[
  M_i N_i \equiv 1 \pmod{m_i}.
\]
L'entier $N_i$ est l'inverse de $M_i$ modulo $m_i$, calculable par
l'algorithme d'Euclide étendu\index{algorithme d'Euclide étendu}.

Posons alors
\begin{equation}
\label{eq:crt-solution}
  x_0 = \sum_{i=1}^{k} a_i M_i N_i.
\end{equation}

Vérifions que $x_0$ résout le système. Fixons un indice $j$. Pour $i \neq j$,
le terme $a_i M_i N_i$ est divisible par $m_j$ (car $M_i$ contient le facteur
$m_j$). Ainsi
\[
  x_0 \equiv a_j M_j N_j \pmod{m_j}.
\]
Or $M_j N_j \equiv 1 \pmod{m_j}$ par construction, donc
$x_0 \equiv a_j \pmod{m_j}$.

\medskip
\noindent\textbf{Unicité.}
Supposons que $x_0$ et $x_1$ soient deux solutions. Pour tout $i$, on a
$x_0 \equiv x_1 \pmod{m_i}$, donc $m_i \mid (x_0 - x_1)$.
Comme les $m_i$ sont deux à deux premiers entre eux, leur produit
$M = m_1 \cdots m_k$ divise aussi $x_0 - x_1$. Ainsi $x_0 \equiv x_1 \pmod{M}$.
\end{proof}

\begin{remarque}
\label{rem:crt-construction}
La formule~\eqref{eq:crt-solution} fournit un algorithme effectif~:
\begin{enumerate}[label=(\roman*)]
  \item Calculer $M = m_1 \cdots m_k$ et les quotients partiels $M_i = M / m_i$.
  \item Pour chaque $i$, déterminer $N_i \equiv M_i^{-1} \pmod{m_i}$ par
        l'algorithme d'Euclide étendu.
  \item Former la somme $x_0 = \sum a_i M_i N_i$ et la réduire modulo $M$.
\end{enumerate}
\end{remarque}

%--------------------------------------------------------------------
\section{Exemples détaillés}
\label{sec:crt:exemples}
%--------------------------------------------------------------------

\begin{exemple}[Problème de Sun Tzu]
\label{ex:crt-sun-tzu}
\index{Sun Tzu!problème de}
Résolvons le système
\[
  x \equiv 2 \pmod{3}, \qquad x \equiv 3 \pmod{5}, \qquad x \equiv 2 \pmod{7}.
\]
Les modules $3, 5, 7$ sont deux à deux premiers entre eux, et $M = 105$.
\begin{itemize}
  \item $M_1 = 105/3 = 35$. On cherche $N_1$ tel que $35 N_1 \equiv 1 \pmod{3}$.
        Or $35 \equiv 2 \pmod{3}$ et $2 \cdot 2 = 4 \equiv 1 \pmod{3}$,
        donc $N_1 = 2$.
  \item $M_2 = 105/5 = 21$. On cherche $N_2$ tel que $21 N_2 \equiv 1 \pmod{5}$.
        Or $21 \equiv 1 \pmod{5}$, donc $N_2 = 1$.
  \item $M_3 = 105/7 = 15$. On cherche $N_3$ tel que $15 N_3 \equiv 1 \pmod{7}$.
        Or $15 \equiv 1 \pmod{7}$, donc $N_3 = 1$.
\end{itemize}
Ainsi,
\[
  x_0 = 2 \cdot 35 \cdot 2 + 3 \cdot 21 \cdot 1 + 2 \cdot 15 \cdot 1
      = 140 + 63 + 30 = 233 \equiv 23 \pmod{105}.
\]
Vérification~: $23 = 7 \cdot 3 + 2$, $23 = 4 \cdot 5 + 3$,
$23 = 3 \cdot 7 + 2$. \checkmark
\end{exemple}

\begin{exemple}[Système à quatre modules]
\label{ex:crt-quatre-modules}
Résolvons
\[
  x \equiv 1 \pmod{2}, \quad x \equiv 2 \pmod{3}, \quad
  x \equiv 3 \pmod{5}, \quad x \equiv 4 \pmod{7}.
\]
Ici $M = 2 \cdot 3 \cdot 5 \cdot 7 = 210$.
\begin{itemize}
  \item $M_1 = 105$, $105 \equiv 1 \pmod{2}$, $N_1 = 1$.
        Contribution~: $1 \cdot 105 \cdot 1 = 105$.
  \item $M_2 = 70$, $70 \equiv 1 \pmod{3}$, $N_2 = 1$.
        Contribution~: $2 \cdot 70 \cdot 1 = 140$.
  \item $M_3 = 42$, $42 \equiv 2 \pmod{5}$, $2 \cdot 3 = 6 \equiv 1 \pmod{5}$,
        $N_3 = 3$. Contribution~: $3 \cdot 42 \cdot 3 = 378$.
  \item $M_4 = 30$, $30 \equiv 2 \pmod{7}$, $2 \cdot 4 = 8 \equiv 1 \pmod{7}$,
        $N_4 = 4$. Contribution~: $4 \cdot 30 \cdot 4 = 480$.
\end{itemize}
Total~: $x_0 = 105 + 140 + 378 + 480 = 1103 \equiv 53 \pmod{210}$.

Vérification~: $53 = 26 \cdot 2 + 1$, $53 = 17 \cdot 3 + 2$,
$53 = 10 \cdot 5 + 3$, $53 = 7 \cdot 7 + 4$. \checkmark
\end{exemple}

%--------------------------------------------------------------------
\section{Interprétation algébrique}
\label{sec:crt:algebrique}
%--------------------------------------------------------------------

Le théorème des restes chinois admet une reformulation élégante en termes
d'isomorphisme d'anneaux\index{isomorphisme d'anneaux}.

\begin{theoreme}[Forme algébrique du CRT]
\label{thm:crt-algebrique}
\index{théorème des restes chinois!forme algébrique}
Soient $m, n \geq 2$ des entiers tels que $\pgcd(m, n) = 1$. L'application
\[
  \Phi : \Z/mn\Z \longrightarrow \Z/m\Z \times \Z/n\Z, \qquad
  \overline{x} \longmapsto (\overline{x}_m,\, \overline{x}_n)
\]
(où $\overline{x}_m$ désigne la classe de $x$ modulo $m$) est un
\emph{isomorphisme d'anneaux}.
\end{theoreme}

\begin{proof}
\label{proof:crt-algebrique}
L'application $\Phi$ est clairement un morphisme d'anneaux (elle préserve
l'addition et la multiplication). Son noyau est
\[
  \ker \Phi = \{ \overline{x} \in \Z/mn\Z : m \mid x \text{ et } n \mid x \}
            = \{ \overline{x} \in \Z/mn\Z : mn \mid x \}
            = \{ \overline{0} \},
\]
la dernière égalité résultant de $\pgcd(m, n) = 1$. Ainsi $\Phi$ est injectif.
Comme les deux ensembles ont le même cardinal $mn$, l'application $\Phi$ est
bijective, donc un isomorphisme.
\end{proof}

\begin{corollaire}[Forme générale]
\label{cor:crt-general}
Si $m_1, \ldots, m_k$ sont deux à deux premiers entre eux et $M = m_1 \cdots m_k$,
alors
\[
  \Z/M\Z \;\simeq\; \Z/m_1\Z \times \Z/m_2\Z \times \cdots \times \Z/m_k\Z
\]
en tant qu'anneaux.
\end{corollaire}

\medskip

La figure~\ref{fig:crt-iso} illustre l'isomorphisme du CRT pour $m = 3$, $n = 5$.

\begin{figure}[ht]
\centering
\begin{tikzpicture}[
    every node/.style={font=\small},
    >=Stealth
  ]
  % Left: Z/15Z as a column
  \node[font=\bfseries] at (0, 5.5) {$\Z/15\Z$};
  \foreach \k in {0,1,...,14} {
    \pgfmathsetmacro{\ypos}{5 - \k * 0.33}
    \node (L\k) at (0, \ypos) {$\overline{\k}$};
  }

  % Right: Z/3Z x Z/5Z as a grid
  \node[font=\bfseries] at (6, 5.5) {$\Z/3\Z \times \Z/5\Z$};
  \foreach \i in {0,1,2} {
    \foreach \j in {0,1,2,3,4} {
      \pgfmathsetmacro{\xpos}{4.5 + \i}
      \pgfmathsetmacro{\ypos}{5 - \j * 0.8}
      \node (R\i\j) at (\xpos, \ypos) {$(\overline{\i},\overline{\j})$};
    }
  }

  % Arrows for a few elements
  \draw[->,blue!70!black,thick] (L0.east) -- (R00.west) node[midway,above,font=\tiny] {};
  \draw[->,blue!70!black,thick] (L1.east) -- (R11.west);
  \draw[->,blue!70!black,thick] (L2.east) -- (R22.west);
  \draw[->,blue!70!black,thick] (L3.east) -- (R03.west);
  \draw[->,blue!70!black,thick] (L5.east) -- (R20.west);

  % Label
  \node[font=\Large] at (2.2, 2.5) {$\Phi$};
  \draw[->,very thick] (1, 3) -- (3.5, 3);
  \node[below,font=\footnotesize] at (2.2, 2.8) {isomorphisme};
\end{tikzpicture}
\caption{Isomorphisme du CRT~: $\Z/15\Z \simeq \Z/3\Z \times \Z/5\Z$.
Quelques correspondances sont indiquées.}
\label{fig:crt-iso}
\end{figure}

%--------------------------------------------------------------------
\section{Conséquences pour la fonction d'Euler}
\label{sec:crt:euler}
%--------------------------------------------------------------------

\begin{corollaire}[Multiplicativité de $\varphi$]
\label{cor:phi-multiplicative}
\index{fonction d'Euler!multiplicativité}
Si $\pgcd(m, n) = 1$, alors $\varphi(mn) = \varphi(m) \cdot \varphi(n)$.
\end{corollaire}

\begin{proof}
Par le théorème~\ref{thm:crt-algebrique}, $\Z/mn\Z \simeq \Z/m\Z \times \Z/n\Z$.
Cet isomorphisme envoie les éléments inversibles sur les éléments inversibles~:
\[
  (\Z/mn\Z)^\times \;\simeq\; (\Z/m\Z)^\times \times (\Z/n\Z)^\times.
\]
En effet, un couple $(a, b)$ est inversible dans le produit si et seulement si
$a$ est inversible dans $\Z/m\Z$ et $b$ est inversible dans $\Z/n\Z$.
En prenant les cardinaux~:
\[
  \varphi(mn) = \abs{(\Z/mn\Z)^\times}
             = \abs{(\Z/m\Z)^\times} \cdot \abs{(\Z/n\Z)^\times}
             = \varphi(m) \cdot \varphi(n). \qedhere
\]
\end{proof}

\begin{corollaire}[Formule d'Euler revisitée]
\label{cor:formule-euler}
Pour tout $n \geq 2$ ayant la décomposition $n = p_1^{\alpha_1} \cdots p_r^{\alpha_r}$,
\[
  \varphi(n) = n \prod_{p \mid n} \left(1 - \frac{1}{p}\right).
\]
\end{corollaire}

\begin{proof}
Par la multiplicativité (Corollaire~\ref{cor:phi-multiplicative}),
$\varphi(n) = \prod_{i=1}^{r} \varphi(p_i^{\alpha_i})$.
Or $\varphi(p^k) = p^k - p^{k-1} = p^k(1 - 1/p)$ pour tout premier $p$ et
$k \geq 1$. Le résultat suit.
\end{proof}

%--------------------------------------------------------------------
\section{Applications}
\label{sec:crt:applications}
%--------------------------------------------------------------------

\subsection{Calcul calendaire}
\label{subsec:crt:calendrier}
\index{calcul calendaire}

\begin{exemple}[Jour de la semaine]
\label{ex:crt-calendrier}
Supposons que nous sachions qu'un événement a lieu un nombre de jours $x$ après
un lundi de référence, avec les contraintes~:
\[
  x \equiv 3 \pmod{7} \quad \text{(c'est un jeudi)}, \qquad
  x \equiv 0 \pmod{4} \quad \text{(multiple de 4 jours)}.
\]
Comme $\pgcd(7, 4) = 1$, le CRT donne une solution unique modulo $28$.
On a $M_1 = 4$, $N_1 \equiv 4^{-1} \equiv 2 \pmod{7}$ (car $4 \cdot 2 = 8 \equiv 1$),
et $M_2 = 7$, $N_2 \equiv 7^{-1} \equiv 3 \pmod{4}$ (car $7 \cdot 3 = 21 \equiv 1$).
Donc $x_0 = 3 \cdot 4 \cdot 2 + 0 \cdot 7 \cdot 3 = 24 \equiv 24 \pmod{28}$.
Le $24$\textsuperscript{e}~jour après le lundi tombe bien un jeudi, et c'est un
multiple de $4$.
\end{exemple}

\subsection{Partage de secret}
\label{subsec:crt:secret}
\index{partage de secret}

\begin{exemple}[Partage de secret par le CRT]
\label{ex:crt-secret}
Un secret $s = 47$ est encodé modulo des nombres premiers entre eux $m_1 = 11$
et $m_2 = 13$. On distribue~:
\begin{itemize}
  \item À la personne A~: le reste $r_1 = 47 \bmod 11 = 3$.
  \item À la personne B~: le reste $r_2 = 47 \bmod 13 = 8$.
\end{itemize}
Aucune des deux personnes ne peut reconstituer $s$ seule, car il existe $13$
(resp.\ $11$) solutions possibles modulo $143$. En revanche, la combinaison
des deux parts permet de retrouver $s$ par le CRT~:
\[
  x \equiv 3 \pmod{11}, \quad x \equiv 8 \pmod{13}
  \implies x \equiv 47 \pmod{143}.
\]
\end{exemple}

\subsection{Codes détecteurs d'erreurs}
\label{subsec:crt:codes}
\index{codes détecteurs d'erreurs}

\begin{exemple}[Codage par résidus]
\label{ex:crt-codes}
Un système de \emph{codage par résidus}\index{codage par résidus} (Residue Number
System, RNS) représente un entier $x \in \{0, 1, \ldots, M-1\}$ par le $k$-uplet
$(x \bmod m_1, \ldots, x \bmod m_k)$ où les $m_i$ sont deux à deux premiers entre eux
et $M = m_1 \cdots m_k$. Pour détecter des erreurs, on ajoute un module redondant
$m_{k+1}$ premier avec les précédents. Si le vecteur reçu ne correspond à aucun
$x < M$ via le CRT, une erreur est détectée.
\end{exemple}

\subsection{Accélération CRT de RSA}
\label{subsec:crt:rsa}
\index{RSA!accélération CRT}

\begin{proposition}[Déchiffrement RSA par le CRT]
\label{prop:rsa-crt}
Soit $(n, d)$ la clé privée RSA avec $n = pq$. Pour déchiffrer un message $c$,
au lieu de calculer $c^d \bmod n$ directement, on peut~:
\begin{enumerate}[label=(\roman*)]
  \item Calculer $d_p = d \bmod (p-1)$ et $d_q = d \bmod (q-1)$.
  \item Calculer $m_p = c^{d_p} \bmod p$ et $m_q = c^{d_q} \bmod q$.
  \item Reconstituer $m = c^d \bmod n$ par le CRT à partir de $m_p$ et $m_q$.
\end{enumerate}
Cette méthode est environ \emph{quatre fois plus rapide} que l'exponentiation
directe modulo~$n$, car les exponentiations se font modulo des nombres de taille moitié.
\end{proposition}

\begin{proof}
Par le petit théorème de Fermat, pour tout $a$ non divisible par $p$~:
$a^{p-1} \equiv 1 \pmod{p}$. Écrivons $d = q_p(p-1) + d_p$. Alors
\[
  c^d = c^{q_p(p-1) + d_p} = (c^{p-1})^{q_p} \cdot c^{d_p}
      \equiv c^{d_p} \pmod{p}.
\]
De même $c^d \equiv c^{d_q} \pmod{q}$. Le CRT reconstitue alors
$c^d \bmod n$ à partir de ces deux résidus.

Pour le gain de vitesse, une exponentiation modulaire modulo un nombre de $b$ bits
coûte $O(b^3)$ opérations binaires (par élévation au carré itérée avec multiplications
de nombres de $b$ bits). Modulo $p$ ou $q$ de taille $b/2$, chaque exponentiation
coûte $O((b/2)^3) = O(b^3/8)$. Deux telles exponentiations coûtent $O(b^3/4)$,
soit un facteur~$4$ de gain.
\end{proof}

%--------------------------------------------------------------------
\section{Exercices}
\label{sec:crt:exercices}
%--------------------------------------------------------------------

\begin{exercice}[$\star$]
\label{exo:crt-1}
Résoudre le système
$x \equiv 1 \pmod{5}$, $x \equiv 2 \pmod{7}$, $x \equiv 3 \pmod{11}$.
\end{exercice}

\begin{exercice}[$\star$]
\label{exo:crt-2}
Trouver le plus petit entier positif $x$ tel que
$x \equiv 3 \pmod{4}$, $x \equiv 5 \pmod{9}$, $x \equiv 7 \pmod{11}$.
\end{exercice}

\begin{exercice}[$\star$]
\label{exo:crt-3}
Montrer que si $\pgcd(m, n) = 1$, alors le système
$x \equiv a \pmod{m}$, $x \equiv b \pmod{n}$ a exactement une solution
modulo $mn$, directement à partir de Bézout (sans invoquer le CRT).
\end{exercice}

\begin{exercice}[$\star\star$]
\label{exo:crt-4}
\index{fonction d'Euler!formule}
Soit $n = p_1^{\alpha_1} \cdots p_r^{\alpha_r}$. En utilisant le CRT,
montrer que $(\Z/n\Z)^\times \simeq \prod_{i=1}^r (\Z/p_i^{\alpha_i}\Z)^\times$,
et en déduire la formule $\varphi(n) = n \prod_{p \mid n}(1 - 1/p)$.
\end{exercice}

\begin{exercice}[$\star\star$]
\label{exo:crt-5}
Soit $f : \Z/mn\Z \to \Z/m\Z \times \Z/n\Z$ l'isomorphisme du CRT
avec $\pgcd(m,n) = 1$. Déterminer explicitement $f^{-1}(1, 0)$ et
$f^{-1}(0, 1)$. Montrer que ces deux éléments sont des idempotents
orthogonaux dans $\Z/mn\Z$.
\end{exercice}

\begin{exercice}[$\star\star$]
\label{exo:crt-6}
Montrer que si $\pgcd(m, n) = d > 1$, le système
$x \equiv a \pmod{m}$, $x \equiv b \pmod{n}$ a une solution si et seulement
si $a \equiv b \pmod{d}$, et que dans ce cas la solution est unique modulo
$\ppcm(m, n)$.
\end{exercice}

\begin{exercice}[$\star\star\star$]
\label{exo:crt-7}
\index{nombres de Carmichael}
On dit que $n$ est un \emph{nombre de Carmichael} si $n$ est composé et
$a^{n-1} \equiv 1 \pmod{n}$ pour tout $a$ avec $\pgcd(a, n) = 1$.
En utilisant le CRT, montrer le \emph{critère de Korselt}~: $n$ est un nombre
de Carmichael si et seulement si $n$ est sans facteur carré et, pour tout
premier $p$ divisant $n$, on a $(p - 1) \mid (n - 1)$.
\end{exercice}

\begin{exercice}[$\star\star\star$]
\label{exo:crt-8}
Soit $p$ un nombre premier impair et $a$ un entier non divisible par $p$.
En utilisant le CRT avec les modules $p$ et $p-1$, montrer qu'il existe
un entier $x$ tel que $x^x \equiv a \pmod{p}$.
\textit{Indication~: chercher $x$ tel que $x \equiv g^k \pmod{p}$ et
$x \equiv k \pmod{p-1}$ pour un choix convenable de $k$.}
\end{exercice}

%--------------------------------------------------------------------
\section*{Résumé du chapitre}
\addcontentsline{toc}{section}{Résumé du chapitre}
%--------------------------------------------------------------------

\begin{itemize}
  \item Le \textbf{théorème des restes chinois} affirme que si les modules
        $m_1, \ldots, m_k$ sont deux à deux premiers entre eux, un système de
        congruences simultanées admet une unique solution modulo
        $M = m_1 \cdots m_k$.
  \item La solution se construit explicitement via la formule
        $x_0 = \sum a_i M_i N_i$ où $M_i = M/m_i$ et $N_i \equiv M_i^{-1} \pmod{m_i}$.
  \item \textbf{Forme algébrique}~: $\Z/M\Z \simeq \Z/m_1\Z \times \cdots \times \Z/m_k\Z$.
  \item La \textbf{multiplicativité} de la fonction d'Euler $\varphi$ est une
        conséquence directe du CRT.
  \item Les \textbf{applications} incluent le calcul calendaire, le partage de secret,
        les codes par résidus et l'accélération RSA.
\end{itemize}


%%%%%%%%%%%%%%%%%%%%%%%%%%%%%%%%%%%%%%%%%%%%%%%%%%%%%%%%%%%%%%%%%%%%
%  CHAPITRE 6 — RÉSIDUS QUADRATIQUES ET RÉCIPROCITÉ QUADRATIQUE
%%%%%%%%%%%%%%%%%%%%%%%%%%%%%%%%%%%%%%%%%%%%%%%%%%%%%%%%%%%%%%%%%%%%
\chapter{Résidus Quadratiques et Réciprocité Quadratique}
\label{chap:qr}

\section*{Perspective historique}
\addcontentsline{toc}{section}{Perspective historique}

La question de savoir quand une congruence $x^2 \equiv a \pmod{p}$ admet une
solution est l'un des problèmes centraux de la théorie des nombres.
\textbf{Leonhard Euler}\index{Euler} (1707--1783) étudia le premier cette question
de manière systématique et découvrit le critère qui porte son nom.
\textbf{Adrien-Marie Legendre}\index{Legendre} (1752--1833) introduisit le symbole
$\legendre{a}{p}$ et formula la loi de réciprocité quadratique, mais sa
<<~preuve~>> comportait une lacune.

C'est \textbf{Carl Friedrich Gauss}\index{Gauss} (1777--1855) qui donna la première
démonstration complète en 1796, à l'âge de dix-huit ans. Gauss qualifia ce
résultat de \emph{theorema aureum}\index{theorema aureum} (<<~théorème d'or~>>)
et en donna pas moins de \emph{huit démonstrations différentes} au cours de sa
vie, témoignant de la richesse du sujet. La preuve que nous présenterons utilise
le lemme de Gauss et le comptage de points du réseau, dans l'esprit de la
démonstration d'Eisenstein\index{Eisenstein}.

%--------------------------------------------------------------------
\section{Résidus quadratiques modulo un premier}
\label{sec:qr:definition}
%--------------------------------------------------------------------

Dans tout ce chapitre, $p$ désigne un nombre premier \emph{impair} (sauf mention
contraire).

\begin{definition}[Résidu quadratique]
\label{def:qr}
\index{résidu quadratique}
Soit $a \in \Z$ avec $p \nmid a$. On dit que $a$ est un \emph{résidu quadratique}
modulo $p$ (noté QR\index{QR}) s'il existe $x \in \Z$ tel que $x^2 \equiv a \pmod{p}$.
Sinon, on dit que $a$ est un \emph{non-résidu quadratique} (noté QNR\index{QNR}).
\end{definition}

\begin{proposition}[Dénombrement des QR]
\label{prop:denombrement-qr}
\index{résidu quadratique!dénombrement}
Parmi les éléments non nuls de $\Z/p\Z$, il y a exactement $\frac{p-1}{2}$ résidus
quadratiques et $\frac{p-1}{2}$ non-résidus quadratiques.
\end{proposition}

\begin{proof}
Considérons l'application $\sigma : (\Z/p\Z)^\times \to (\Z/p\Z)^\times$
définie par $\sigma(\overline{x}) = \overline{x}^2$. Son image est exactement
l'ensemble des QR. Montrons que $\sigma$ est deux-pour-un~: si $x^2 \equiv y^2 \pmod{p}$,
alors $p \mid (x-y)(x+y)$, donc $x \equiv y$ ou $x \equiv -y \pmod{p}$.
Comme $p$ est impair, $x \not\equiv -x \pmod{p}$ pour $x \not\equiv 0$,
donc chaque QR a exactement deux antécédents. Il y a donc
$\frac{p-1}{2}$ résidus quadratiques et $p - 1 - \frac{p-1}{2} = \frac{p-1}{2}$
non-résidus.
\end{proof}

\begin{exemple}
\label{ex:qr-mod-13}
\index{résidu quadratique!exemple}
Pour $p = 13$, les carrés modulo $13$ sont~:
\[
  1^2 \equiv 1, \quad 2^2 \equiv 4, \quad 3^2 \equiv 9, \quad
  4^2 \equiv 3, \quad 5^2 \equiv 12, \quad 6^2 \equiv 10.
\]
Les QR modulo $13$ sont donc $\{1, 3, 4, 9, 10, 12\}$, et les QNR sont
$\{2, 5, 6, 7, 8, 11\}$. On a bien $6$ éléments dans chaque catégorie.
\end{exemple}

%--------------------------------------------------------------------
\section{Critère d'Euler}
\label{sec:qr:euler}
%--------------------------------------------------------------------

\begin{theoreme}[Critère d'Euler]
\label{thm:euler-criterion}
\index{critère d'Euler}
\index{Euler!critère d'}
Soit $p$ un premier impair et $a \in \Z$ avec $p \nmid a$. Alors
\[
  a^{(p-1)/2} \equiv
  \begin{cases}
    +1 \pmod{p} & \text{si } a \text{ est un QR modulo } p, \\
    -1 \pmod{p} & \text{si } a \text{ est un QNR modulo } p.
  \end{cases}
\]
\end{theoreme}

\begin{proof}
\index{critère d'Euler!preuve}
Par le petit théorème de Fermat, $a^{p-1} \equiv 1 \pmod{p}$, c'est-à-dire
$p \mid (a^{(p-1)/2} - 1)(a^{(p-1)/2} + 1)$.
Comme $p$ est premier, $a^{(p-1)/2} \equiv 1$ ou $a^{(p-1)/2} \equiv -1 \pmod{p}$.

\medskip
\noindent\textbf{Cas QR.} Supposons $a \equiv x^2 \pmod{p}$. Alors
\[
  a^{(p-1)/2} \equiv x^{p-1} \equiv 1 \pmod{p}
\]
par le petit théorème de Fermat.

\medskip
\noindent\textbf{Cas QNR.} Le polynôme $X^{(p-1)/2} - 1$ a au plus $(p-1)/2$ racines
dans $\Z/p\Z$ (corps). Or les $(p-1)/2$ résidus quadratiques sont des racines
(par le cas précédent). Donc les $(p-1)/2$ non-résidus ne sont \emph{pas} racines
de ce polynôme, et vérifient nécessairement $a^{(p-1)/2} \equiv -1 \pmod{p}$.
\end{proof}

\begin{exemple}
\label{ex:euler-5-mod-13}
Vérifions que $5$ est un QNR modulo $13$~:
$5^{6} = 15625$. Or $15625 = 1201 \cdot 13 + 12$, donc $5^6 \equiv 12 \equiv -1 \pmod{13}$. \checkmark

En pratique, on calcule par élévation au carré~:
$5^2 = 25 \equiv 12$, $5^4 \equiv 12^2 = 144 \equiv 1$, $5^6 \equiv 5^4 \cdot 5^2 \equiv 1 \cdot 12 \equiv -1$. \checkmark
\end{exemple}

%--------------------------------------------------------------------
\section{Le symbole de Legendre}
\label{sec:qr:legendre}
%--------------------------------------------------------------------

\begin{definition}[Symbole de Legendre]
\label{def:legendre}
\index{symbole de Legendre}
\index{Legendre!symbole de}
Pour $p$ premier impair et $a \in \Z$, le \emph{symbole de Legendre} est
\[
  \legendre{a}{p} =
  \begin{cases}
    0  & \text{si } p \mid a, \\
    +1 & \text{si } a \text{ est un QR modulo } p, \\
    -1 & \text{si } a \text{ est un QNR modulo } p.
  \end{cases}
\]
\end{definition}

Le critère d'Euler se réécrit alors de manière concise.

\begin{corollaire}
\label{cor:euler-legendre}
Pour tout $a \in \Z$ et tout premier impair $p$ avec $p \nmid a$~:
\[
  \legendre{a}{p} \equiv a^{(p-1)/2} \pmod{p}.
\]
\end{corollaire}

\begin{proposition}[Propriétés du symbole de Legendre]
\label{prop:legendre-proprietes}
\index{symbole de Legendre!propriétés}
Soient $p$ un premier impair et $a, b \in \Z$.
\begin{enumerate}[label=(\roman*)]
  \item \textbf{Périodicité}~: $\legendre{a}{p} = \legendre{a + p}{p}$.
  \item \textbf{Multiplicativité}\index{symbole de Legendre!multiplicativité}~:
        $\legendre{ab}{p} = \legendre{a}{p} \cdot \legendre{b}{p}$.
  \item \textbf{Carrés}~: $\legendre{a^2}{p} = 1$ si $p \nmid a$.
  \item $\legendre{-1}{p} = (-1)^{(p-1)/2}$.
\end{enumerate}
\end{proposition}

\begin{proof}
\begin{enumerate}[label=(\roman*)]
  \item Immédiat car le symbole ne dépend que de la classe de $a$ modulo $p$.
  \item Par le critère d'Euler,
        $\legendre{ab}{p} \equiv (ab)^{(p-1)/2} = a^{(p-1)/2} b^{(p-1)/2}
        \equiv \legendre{a}{p} \cdot \legendre{b}{p} \pmod{p}$.
        Comme les deux membres sont dans $\{-1, 0, 1\}$ et $p \geq 3$,
        la congruence est une égalité.
  \item $\legendre{a^2}{p} \equiv (a^2)^{(p-1)/2} = a^{p-1} \equiv 1 \pmod{p}$.
  \item $\legendre{-1}{p} \equiv (-1)^{(p-1)/2} \pmod{p}$, et comme
        $(-1)^{(p-1)/2} \in \{-1, 1\}$ et $p \geq 3$, c'est une égalité.
        Concrètement~: $-1$ est QR modulo $p$ si et seulement si
        $p \equiv 1 \pmod{4}$. \qedhere
\end{enumerate}
\end{proof}

\begin{exemple}
\label{ex:legendre-calculs}
Calculons $\legendre{-3}{13}$.
Par multiplicativité~: $\legendre{-3}{13} = \legendre{-1}{13} \cdot \legendre{3}{13}$.

$\legendre{-1}{13} = (-1)^{6} = 1$ car $13 \equiv 1 \pmod{4}$.

Pour $\legendre{3}{13}$~: $3^6 = 729$, $729 = 56 \cdot 13 + 1$, donc
$\legendre{3}{13} = 1$.

Ainsi $\legendre{-3}{13} = 1 \cdot 1 = 1$, et $-3 \equiv 10$ est bien un
QR modulo $13$ (on vérifie~: $6^2 = 36 \equiv 10 \pmod{13}$). \checkmark
\end{exemple}

%--------------------------------------------------------------------
\section{Le lemme de Gauss}
\label{sec:qr:gauss-lemma}
%--------------------------------------------------------------------

\begin{notation}
\label{not:representants}
Pour $p$ premier impair, on note
$S = \{1, 2, \ldots, (p-1)/2\}$ l'ensemble des \emph{représentants positifs}
des classes non nulles modulo $p$.
\end{notation}

\begin{lemme}[Lemme de Gauss]
\label{lem:gauss}
\index{lemme de Gauss}
\index{Gauss!lemme de}
Soient $p$ un premier impair et $a \in \Z$ avec $p \nmid a$. Pour chaque
$s \in S$, réduisons $as \bmod p$ dans l'intervalle $\{1, \ldots, p-1\}$
et écrivons le résultat sous la forme $\varepsilon_s \cdot r_s$ avec
$r_s \in S$ et $\varepsilon_s \in \{+1, -1\}$ (c'est-à-dire $r_s = as \bmod p$
si $as \bmod p \leq (p-1)/2$, et $r_s = p - (as \bmod p)$ sinon). Alors~:
\begin{enumerate}[label=(\alph*)]
  \item Les $r_s$ pour $s \in S$ forment une permutation de $S$.
  \item $\legendre{a}{p} = \prod_{s \in S} \varepsilon_s = (-1)^\nu$,
        où $\nu$ est le nombre d'éléments $s \in S$ tels que le résidu
        $as \bmod p$ tombe dans $\{(p+1)/2, \ldots, p-1\}$.
\end{enumerate}
\end{lemme}

\begin{proof}
\index{lemme de Gauss!preuve}
\textbf{(a)} Montrons que les $r_s$ sont deux à deux distincts. Supposons
$r_s = r_t$ avec $s \neq t$. Cela signifie que $as \equiv \pm at \pmod{p}$.
Si $as \equiv at$, alors $s \equiv t$ (car $\pgcd(a, p) = 1$), contradiction.
Si $as \equiv -at$, alors $a(s + t) \equiv 0$, donc $p \mid (s + t)$.
Mais $2 \leq s + t \leq p - 1$, contradiction. Ainsi les $(p-1)/2$ valeurs
$r_s$ sont distinctes et appartiennent à $S$ (qui a $(p-1)/2$ éléments),
donc elles forment une permutation de $S$.

\medskip
\noindent\textbf{(b)} Par construction, $as \equiv \varepsilon_s \cdot r_s \pmod{p}$
pour chaque $s \in S$. En prenant le produit sur tous les $s \in S$~:
\[
  \prod_{s \in S} (as) \equiv \prod_{s \in S} (\varepsilon_s \cdot r_s) \pmod{p}.
\]
Le membre de gauche vaut $a^{(p-1)/2} \prod_{s \in S} s$. Le membre de droite
vaut $\left(\prod \varepsilon_s\right) \cdot \prod_{s \in S} r_s$.
Par (a), $\prod r_s = \prod s$, et ce produit est non nul modulo $p$.
En simplifiant~:
\[
  a^{(p-1)/2} \equiv \prod_{s \in S} \varepsilon_s = (-1)^\nu \pmod{p}.
\]
Par le critère d'Euler, $a^{(p-1)/2} \equiv \legendre{a}{p} \pmod{p}$, d'où le résultat.
\end{proof}

\begin{exemple}[Application du lemme de Gauss]
\label{ex:gauss-lemma}
Calculons $\legendre{5}{11}$. Ici $p = 11$, $S = \{1, 2, 3, 4, 5\}$.
\[
\begin{array}{c|ccccc}
  s       & 1 & 2 & 3 & 4 & 5 \\ \hline
  5s \bmod 11 & 5 & 10 & 4 & 9 & 3 \\
  r_s     & 5 & 1 & 4 & 2 & 3 \\
  \varepsilon_s & +1 & -1 & +1 & -1 & +1
\end{array}
\]
Le nombre $\nu$ de signes négatifs est $\nu = 2$, donc
$\legendre{5}{11} = (-1)^2 = +1$. Ainsi $5$ est un QR modulo $11$.

Vérification~: $4^2 = 16 \equiv 5 \pmod{11}$. \checkmark
\end{exemple}

%--------------------------------------------------------------------
\section{Premier supplément~: $\legendre{-1}{p}$}
\label{sec:qr:supplement1}
%--------------------------------------------------------------------

\begin{theoreme}[Premier supplément]
\label{thm:premier-supplement}
\index{premier supplément}
\index{réciprocité quadratique!premier supplément}
Pour tout premier impair $p$~:
\[
  \legendre{-1}{p} = (-1)^{(p-1)/2} =
  \begin{cases}
    +1 & \text{si } p \equiv 1 \pmod{4}, \\
    -1 & \text{si } p \equiv 3 \pmod{4}.
  \end{cases}
\]
\end{theoreme}

\begin{proof}
C'est la propriété (iv) de la Proposition~\ref{prop:legendre-proprietes},
conséquence immédiate du critère d'Euler.
\end{proof}

\begin{corollaire}
\label{cor:somme-deux-carres-premier}
Un premier impair $p$ est somme de deux carrés ($p = a^2 + b^2$) seulement si
$p \equiv 1 \pmod{4}$.
\end{corollaire}

\begin{proof}
Si $p = a^2 + b^2$ avec $p$ impair, alors $a^2 + b^2 \equiv 0 \pmod{p}$, et
$b \not\equiv 0$, donc $a^2 \equiv -b^2$, d'où $(ab^{-1})^2 \equiv -1 \pmod{p}$.
Ainsi $-1$ est un QR, ce qui impose $p \equiv 1 \pmod{4}$.
\end{proof}

%--------------------------------------------------------------------
\section{Second supplément~: $\legendre{2}{p}$}
\label{sec:qr:supplement2}
%--------------------------------------------------------------------

\begin{theoreme}[Second supplément]
\label{thm:second-supplement}
\index{second supplément}
\index{réciprocité quadratique!second supplément}
Pour tout premier impair $p$~:
\[
  \legendre{2}{p} = (-1)^{(p^2-1)/8} =
  \begin{cases}
    +1 & \text{si } p \equiv \pm 1 \pmod{8}, \\
    -1 & \text{si } p \equiv \pm 3 \pmod{8}.
  \end{cases}
\]
\end{theoreme}

\begin{proof}
\index{second supplément!preuve}
Appliquons le lemme de Gauss avec $a = 2$. Pour $s \in S = \{1, \ldots, (p-1)/2\}$,
le produit $2s$ modulo $p$ est simplement $2s$ si $2s \leq p-1$
(c'est-à-dire $s \leq (p-1)/2$, toujours vrai), mais il faut déterminer
si $2s$ tombe dans $\{1, \ldots, (p-1)/2\}$ ou dans $\{(p+1)/2, \ldots, p-1\}$.
On a $\varepsilon_s = -1$ si et seulement si $2s > (p-1)/2$,
c'est-à-dire $s > (p-1)/4$. Donc
\[
  \nu = \abs{\left\{ s \in \{1, \ldots, \tfrac{p-1}{2}\} : s > \tfrac{p-1}{4} \right\}}
      = \tfrac{p-1}{2} - \floor{\tfrac{p-1}{4}}.
\]
Un calcul par cas selon $p \bmod 8$ montre que $\nu \equiv (p^2 - 1)/8 \pmod{2}$~:
\begin{itemize}
  \item $p = 8k + 1$~: $\nu = 4k - 2k = 2k$, $(p^2-1)/8 = 8k^2 + 2k$, parité identique.
  \item $p = 8k + 3$~: $\nu = 4k+1 - (2k) = 2k+1$, $(p^2-1)/8 = 8k^2+6k+1$, parité identique.
  \item $p = 8k + 5$~: $\nu = 4k+2 - (2k+1) = 2k+1$, $(p^2-1)/8 = 8k^2+10k+3$, parité identique.
  \item $p = 8k + 7$~: $\nu = 4k+3 - (2k+1) = 2k+2$, $(p^2-1)/8 = 8k^2+14k+6$, parité identique.
\end{itemize}
Donc $\legendre{2}{p} = (-1)^\nu = (-1)^{(p^2-1)/8}$.
\end{proof}

%--------------------------------------------------------------------
\section{La loi de réciprocité quadratique}
\label{sec:qr:reciprocite}
%--------------------------------------------------------------------

Nous arrivons au résultat central de ce chapitre, qualifié de <<~théorème d'or~>>
par Gauss.

\begin{theoreme}[Loi de réciprocité quadratique]
\label{thm:reciprocite-quadratique}
\index{réciprocité quadratique!loi de}
\index{loi de réciprocité quadratique}
\index{theorema aureum}
Soient $p$ et $q$ deux premiers impairs distincts. Alors
\[
  \legendre{p}{q} \cdot \legendre{q}{p}
  = (-1)^{\frac{p-1}{2} \cdot \frac{q-1}{2}}.
\]
De manière équivalente~:
\[
  \legendre{p}{q} = \legendre{q}{p} \quad \text{sauf si } p \equiv q \equiv 3 \pmod{4},
  \text{ auquel cas } \legendre{p}{q} = -\legendre{q}{p}.
\]
\end{theoreme}

Avant de démontrer ce théorème, établissons un lemme combinatoire clé.

\begin{lemme}[Lemme d'Eisenstein]
\label{lem:eisenstein}
\index{Eisenstein!lemme d'}
\index{lemme d'Eisenstein}
Avec les notations du lemme de Gauss (Lemme~\ref{lem:gauss}), le nombre $\nu$
de changements de signe est
\[
  \nu = \sum_{s=1}^{(p-1)/2} \floor{\frac{as}{p}}.
\]
\end{lemme}

\begin{proof}
Pour chaque $s \in \{1, \ldots, (p-1)/2\}$, effectuons la division euclidienne
de $as$ par $p$~:
\[
  as = \floor{\frac{as}{p}} \cdot p + (as \bmod p).
\]
Le reste $as \bmod p$ appartient à $\{1, \ldots, p-1\}$ (car $p \nmid as$).
Si $as \bmod p \leq (p-1)/2$, alors $\varepsilon_s = +1$ et $r_s = as \bmod p$.
Si $as \bmod p > (p-1)/2$, alors $\varepsilon_s = -1$ et $r_s = p - (as \bmod p)$.
Dans le second cas~:
\[
  as = \floor{\frac{as}{p}} \cdot p + (p - r_s)
     = \left(\floor{\frac{as}{p}} + 1\right) p - r_s.
\]
Sommons sur tout $s \in S$~:
\[
  a \sum_{s=1}^{(p-1)/2} s
  = p \sum_{s=1}^{(p-1)/2} \floor{\frac{as}{p}}
    + \sum_{\varepsilon_s = +1} r_s + \sum_{\varepsilon_s = -1} (p - r_s).
\]
La dernière somme se réécrit~:
$\sum_{\varepsilon_s = -1} p - \sum_{\varepsilon_s = -1} r_s = \nu p - \sum_{\varepsilon_s = -1} r_s$.
Donc~:
\[
  a \sum s = p \left(\sum \floor{\frac{as}{p}} + \nu\right)
           + \sum_{\varepsilon_s = +1} r_s - \sum_{\varepsilon_s = -1} r_s.
\]
Or, comme les $r_s$ forment une permutation de $S$~:
$\sum r_s = \sum s$ et
$\sum_{\varepsilon_s = +1} r_s - \sum_{\varepsilon_s = -1} r_s$
plus $2\sum_{\varepsilon_s = -1} r_s = \sum s$. On obtient, en travaillant
modulo $2$~:
\[
  a \sum s \equiv \sum s + 2\left(\text{termes}\right) + \nu \cdot p \pmod{2}.
\]
Comme $\sum s$ et $p$ sont impairs (on peut les simplifier dans le calcul de parité),
une analyse modulo $2$ donne $\nu \equiv \sum \floor{as/p} \pmod{2}$.
Puisque seule la parité de $\nu$ intervient dans $\legendre{a}{p} = (-1)^\nu$,
le résultat est établi.
\end{proof}

Nous pouvons maintenant démontrer la loi de réciprocité quadratique.

\begin{proof}[Preuve de la loi de réciprocité quadratique (méthode d'Eisenstein)]
\label{proof:qr}
\index{réciprocité quadratique!preuve}
\index{Eisenstein!preuve d'}

Par le lemme de Gauss et le Lemme~\ref{lem:eisenstein}, on a
\[
  \legendre{q}{p} = (-1)^{A}, \quad \text{où} \quad
  A = \sum_{i=1}^{(p-1)/2} \floor{\frac{iq}{p}},
\]
et de même
\[
  \legendre{p}{q} = (-1)^{B}, \quad \text{où} \quad
  B = \sum_{j=1}^{(q-1)/2} \floor{\frac{jp}{q}}.
\]
Il suffit donc de montrer que
\begin{equation}
\label{eq:AB}
  A + B = \frac{p-1}{2} \cdot \frac{q-1}{2}.
\end{equation}

\textbf{Comptage de points du réseau.}
\index{comptage de points du réseau}
Considérons le rectangle
$R = \{(i, j) \in \Z^2 : 1 \leq i \leq (p-1)/2, \; 1 \leq j \leq (q-1)/2\}$.
Ce rectangle contient $\frac{p-1}{2} \cdot \frac{q-1}{2}$ points entiers.

Considérons la droite $D : j = \frac{q}{p} \, i$ (qui passe par l'origine et
le point $(p, q)$). Comme $p \nmid q$ et les deux sont premiers, aucun point
entier $(i, j)$ de $R$ ne se trouve exactement sur $D$. En effet, $j = qi/p$
entier impliquerait $p \mid i$, impossible pour $1 \leq i \leq (p-1)/2$.

Comptons les points de $R$ situés \emph{en dessous} de $D$ (c'est-à-dire
$j < qi/p$)~:
\[
  \abs{\{(i, j) \in R : j < qi/p\}} = \sum_{i=1}^{(p-1)/2} \floor{\frac{qi}{p}} = A.
\]

Par symétrie, les points situés \emph{au-dessus} de $D$ (c'est-à-dire $j > qi/p$,
soit $i < pj/q$) sont comptés par~:
\[
  \abs{\{(i, j) \in R : i < pj/q\}} = \sum_{j=1}^{(q-1)/2} \floor{\frac{pj}{q}} = B.
\]

Puisque aucun point n'est sur $D$, chaque point de $R$ est soit au-dessus, soit
en dessous. Donc $A + B = \frac{p-1}{2} \cdot \frac{q-1}{2}$,
ce qui établit~\eqref{eq:AB} et achève la preuve.
\end{proof}

\bigskip

La figure~\ref{fig:eisenstein} illustre ce comptage pour $p = 5$, $q = 7$.

\begin{figure}[ht]
\centering
\begin{tikzpicture}[scale=0.85]
  % Grid
  \draw[gray!30, very thin] (0,0) grid (3,4);

  % Axes
  \draw[->,thick] (-0.3,0) -- (3.5,0) node[right] {$i$};
  \draw[->,thick] (0,-0.3) -- (0,4.5) node[above] {$j$};

  % Diagonal line D: j = (7/5)i
  \draw[red, thick, dashed] (0,0) -- (3, 4.2)
    node[right, font=\small] {$D : j = \frac{7}{5}i$};

  % Rectangle boundary
  \draw[blue!60!black, thick] (0.5,-0.1) rectangle (2.5, 3.1);

  % Points below D (counted by A)
  % i=1: floor(7/5) = 1 -> j=1
  \fill[blue!70!black] (1,1) circle (3pt);
  % i=2: floor(14/5) = 2 -> j=1,2
  \fill[blue!70!black] (2,1) circle (3pt);
  \fill[blue!70!black] (2,2) circle (3pt);

  % Points above D (counted by B)
  % i=1: j=2,3 (since 7/5 = 1.4, so j>=2)
  \fill[green!50!black] (1,2) circle (3pt);
  \fill[green!50!black] (1,3) circle (3pt);
  % i=2: j=3 (since 14/5 = 2.8, so j>=3)
  \fill[green!50!black] (2,3) circle (3pt);

  % Labels
  \node[below, font=\footnotesize] at (1,0) {$1$};
  \node[below, font=\footnotesize] at (2,0) {$2$};
  \node[left, font=\footnotesize] at (0,1) {$1$};
  \node[left, font=\footnotesize] at (0,2) {$2$};
  \node[left, font=\footnotesize] at (0,3) {$3$};

  % Legend
  \fill[blue!70!black] (4.5, 3.5) circle (3pt);
  \node[right, font=\small] at (4.8, 3.5) {En dessous de $D$ ($A = 3$)};
  \fill[green!50!black] (4.5, 2.8) circle (3pt);
  \node[right, font=\small] at (4.8, 2.8) {Au-dessus de $D$ ($B = 3$)};
  \node[right, font=\small] at (4.5, 2.1) {$A + B = 6 = \frac{4}{2}\cdot\frac{6}{2}$};
\end{tikzpicture}
\caption{Comptage de points du réseau pour la preuve d'Eisenstein avec
$p = 5$ et $q = 7$. Le rectangle $R$ contient $2 \times 3 = 6$ points
entiers, répartis en $A = 3$ sous la diagonale et $B = 3$ au-dessus.}
\label{fig:eisenstein}
\end{figure}

%--------------------------------------------------------------------
\section{Exemples de calcul du symbole de Legendre}
\label{sec:qr:exemples-calcul}
%--------------------------------------------------------------------

\begin{exemple}[Calcul de $\legendre{23}{59}$]
\label{ex:legendre-23-59}
Par réciprocité quadratique (puisque $59 \equiv 3 \pmod{4}$ et $23 \equiv 3 \pmod{4}$,
on est dans le cas exceptionnel)~:
\[
  \legendre{23}{59} = -\legendre{59}{23} = -\legendre{59 \bmod 23}{23}
  = -\legendre{13}{23}.
\]
Maintenant $23 \equiv 3 \pmod{4}$ et $13 \equiv 1 \pmod{4}$, donc~:
\[
  \legendre{13}{23} = \legendre{23}{13} = \legendre{23 \bmod 13}{13}
  = \legendre{10}{13}.
\]
Décomposons~: $\legendre{10}{13} = \legendre{2}{13} \cdot \legendre{5}{13}$.

$\legendre{2}{13}$~: $13 \equiv 5 \pmod{8}$, donc $\legendre{2}{13} = (-1)^{(169-1)/8} = (-1)^{21} = -1$.

$\legendre{5}{13}$~: $13 \equiv 1 \pmod{4}$, donc $\legendre{5}{13} = \legendre{13}{5} = \legendre{3}{5}$.
Or $5 \equiv 1 \pmod{4}$, donc $\legendre{3}{5} = \legendre{5}{3} = \legendre{2}{3}$.
Et $3 \equiv 3 \pmod{8}$, donc $\legendre{2}{3} = -1$.

Rassemblons~: $\legendre{10}{13} = (-1)(-1) = 1$, puis
$\legendre{13}{23} = 1$, et enfin $\legendre{23}{59} = -1$.

Donc $23$ est un QNR modulo $59$.
\end{exemple}

\begin{exemple}[Calcul de $\legendre{-6}{31}$]
\label{ex:legendre--6-31}
$\legendre{-6}{31} = \legendre{-1}{31} \cdot \legendre{2}{31} \cdot \legendre{3}{31}$.

$\legendre{-1}{31} = (-1)^{15} = -1$ car $31 \equiv 3 \pmod{4}$.

$\legendre{2}{31}$~: $31 \equiv 7 \pmod{8}$, donc $\legendre{2}{31} = +1$.

$\legendre{3}{31}$~: $31 \equiv 3 \pmod{4}$ et $3 \equiv 3 \pmod{4}$, cas exceptionnel~:
$\legendre{3}{31} = -\legendre{31}{3} = -\legendre{1}{3} = -1$.

Bilan~: $\legendre{-6}{31} = (-1)(+1)(-1) = +1$.

Vérification~: $-6 \equiv 25 \pmod{31}$, et $5^2 = 25$. \checkmark
\end{exemple}

%--------------------------------------------------------------------
\section{Le symbole de Jacobi}
\label{sec:qr:jacobi}
%--------------------------------------------------------------------

\begin{definition}[Symbole de Jacobi]
\label{def:jacobi}
\index{symbole de Jacobi}
\index{Jacobi!symbole de}
Soit $n$ un entier impair $\geq 3$ avec la factorisation $n = p_1^{e_1} \cdots p_r^{e_r}$.
Pour $a \in \Z$ avec $\pgcd(a, n) = 1$, le \emph{symbole de Jacobi} est
\[
  \jacobi{a}{n} = \legendre{a}{p_1}^{e_1} \cdots \legendre{a}{p_r}^{e_r}.
\]
\end{definition}

\begin{remarque}
\label{rem:jacobi-attention}
\textbf{Attention}~: $\jacobi{a}{n} = 1$ n'implique \emph{pas} que $a$ soit un QR
modulo $n$. Par exemple, $\jacobi{2}{15} = \legendre{2}{3} \cdot \legendre{2}{5}
= (-1)(-1) = 1$, mais $x^2 \equiv 2 \pmod{15}$ n'a aucune solution
(car $x^2 \equiv 2 \pmod{3}$ n'en a pas). En revanche, $\jacobi{a}{n} = -1$
garantit que $a$ est un QNR modulo $n$.
\end{remarque}

\begin{proposition}[Propriétés du symbole de Jacobi]
\label{prop:jacobi-proprietes}
\index{symbole de Jacobi!propriétés}
Soient $m, n$ des entiers impairs $\geq 3$ et $a, b \in \Z$ premiers avec $m$ et $n$.
\begin{enumerate}[label=(\roman*)]
  \item \textbf{Multiplicativité en haut}~: $\jacobi{ab}{n} = \jacobi{a}{n} \cdot \jacobi{b}{n}$.
  \item \textbf{Multiplicativité en bas}~: $\jacobi{a}{mn} = \jacobi{a}{m} \cdot \jacobi{a}{n}$.
  \item \textbf{Périodicité}~: $\jacobi{a}{n} = \jacobi{a + n}{n}$.
  \item \textbf{Premier supplément}~: $\jacobi{-1}{n} = (-1)^{(n-1)/2}$.
  \item \textbf{Second supplément}~: $\jacobi{2}{n} = (-1)^{(n^2-1)/8}$.
\end{enumerate}
\end{proposition}

\begin{proof}
Les points (i), (ii) et (iii) découlent directement de la définition et des
propriétés correspondantes du symbole de Legendre. Pour (iv), on utilise le fait que
$(n-1)/2 \equiv \sum e_i (p_i - 1)/2 \pmod{2}$ et de même pour (v) avec $(n^2-1)/8$.
\end{proof}

\begin{theoreme}[Réciprocité quadratique pour le symbole de Jacobi]
\label{thm:jacobi-reciprocite}
\index{symbole de Jacobi!réciprocité}
\index{réciprocité quadratique!symbole de Jacobi}
Soient $m, n$ des entiers impairs $\geq 3$ avec $\pgcd(m, n) = 1$. Alors
\[
  \jacobi{m}{n} \cdot \jacobi{n}{m} = (-1)^{\frac{m-1}{2} \cdot \frac{n-1}{2}}.
\]
\end{theoreme}

\begin{proof}
Écrivons $m = p_1 \cdots p_s$ et $n = q_1 \cdots q_t$ (avec répétitions).
Alors
\[
  \jacobi{m}{n} \cdot \jacobi{n}{m}
  = \prod_{i=1}^{s} \prod_{j=1}^{t} \legendre{p_i}{q_j} \cdot \legendre{q_j}{p_i}
  = \prod_{i,j} (-1)^{\frac{p_i - 1}{2} \cdot \frac{q_j - 1}{2}}.
\]
L'exposant total est
\[
  \sum_{i,j} \frac{p_i - 1}{2} \cdot \frac{q_j - 1}{2}
  = \left(\sum_{i} \frac{p_i - 1}{2}\right) \left(\sum_{j} \frac{q_j - 1}{2}\right)
  \equiv \frac{m-1}{2} \cdot \frac{n-1}{2} \pmod{2},
\]
la dernière congruence résultant du fait que
$\sum_i (p_i - 1)/2 \equiv (m-1)/2 \pmod{2}$.
\end{proof}

\subsection{Calcul efficace du symbole de Jacobi}
\label{subsec:qr:jacobi-algorithme}
\index{symbole de Jacobi!algorithme}

Grâce aux propriétés ci-dessus, le symbole de Jacobi $\jacobi{a}{n}$ se calcule
en temps $O(\log^2 n)$ sans connaître la factorisation de $n$, par un algorithme
analogue à celui d'Euclide~:

\begin{enumerate}[label=\textbf{Étape~\arabic*.}]
  \item Réduire $a$ modulo $n$ pour que $0 \leq a < n$.
  \item Si $a = 0$, retourner $0$. Si $a = 1$, retourner $1$.
  \item Écrire $a = 2^e \cdot a'$ avec $a'$ impair. Calculer
        $\jacobi{2}{n}^e = (-1)^{e(n^2-1)/8}$.
  \item Appliquer la réciprocité~: $\jacobi{a'}{n} = (-1)^{(a'-1)(n-1)/4} \cdot \jacobi{n}{a'}$.
  \item Réduire $n$ modulo $a'$ et recommencer.
\end{enumerate}

\begin{exemple}[Calcul de $\jacobi{1001}{9907}$]
\label{ex:jacobi-algorithme}
\index{symbole de Jacobi!exemple de calcul}
Ici $1001 = 7 \times 11 \times 13$ et $9907$ est premier, mais nous n'utilisons
pas ces factorisations.

$\jacobi{1001}{9907}$~: $1001$ et $9907$ sont impairs, $1001 \equiv 1 \pmod{4}$, donc
$\jacobi{1001}{9907} = \jacobi{9907}{1001}$.

$9907 \bmod 1001 = 9907 - 9 \cdot 1001 = 9907 - 9009 = 898$.
$\jacobi{898}{1001} = \jacobi{2}{1001} \cdot \jacobi{449}{1001}$.

$\jacobi{2}{1001}$~: $1001 \equiv 1 \pmod{8}$, donc $\jacobi{2}{1001} = 1$.

$\jacobi{449}{1001}$~: $449 \equiv 1 \pmod{4}$, donc $\jacobi{449}{1001} = \jacobi{1001}{449}$.
$1001 \bmod 449 = 1001 - 2 \cdot 449 = 103$.
$\jacobi{103}{449}$~: $103 \equiv 3 \pmod{4}$, $449 \equiv 1 \pmod{4}$, donc
$\jacobi{103}{449} = \jacobi{449}{103}$.
$449 \bmod 103 = 449 - 4 \cdot 103 = 37$.
$\jacobi{37}{103}$~: $37 \equiv 1 \pmod{4}$, donc $\jacobi{37}{103} = \jacobi{103}{37}$.
$103 \bmod 37 = 103 - 2 \cdot 37 = 29$.
$\jacobi{29}{37}$~: $29 \equiv 1 \pmod{4}$, donc $\jacobi{29}{37} = \jacobi{37}{29}$.
$37 \bmod 29 = 8 = 2^3$.
$\jacobi{8}{29} = \jacobi{2}{29}^3$. Or $29 \equiv 5 \pmod{8}$, donc
$\jacobi{2}{29} = -1$, et $\jacobi{2}{29}^3 = -1$.

Conclusion~: $\jacobi{1001}{9907} = 1 \cdot (-1) = -1$.

Donc $1001$ est un QNR modulo $9907$.
\end{exemple}

%--------------------------------------------------------------------
\section{Visualisation des résidus quadratiques}
\label{sec:qr:visualisation}
%--------------------------------------------------------------------

\begin{figure}[ht]
\centering
\begin{tikzpicture}[scale=0.42]
  % QR pattern mod 29
  \node[font=\bfseries] at (14.5, 15.5) {Résidus quadratiques modulo $29$};

  % QR mod 29: 1,4,5,6,7,9,13,16,20,22,23,24,25,28
  \foreach \k in {1,...,28} {
    \pgfmathsetmacro{\col}{mod(\k-1,7)}
    \pgfmathsetmacro{\row}{int((\k-1)/7)}
    \pgfmathsetmacro{\xpos}{\col * 4 + 2}
    \pgfmathsetmacro{\ypos}{13 - \row * 3}
    % Check if QR
    \pgfmathtruncatemacro{\isQR}{%
      (\k==1 || \k==4 || \k==5 || \k==6 || \k==7 ||
       \k==9 || \k==13 || \k==16 || \k==20 || \k==22 ||
       \k==23 || \k==24 || \k==25 || \k==28) ? 1 : 0}
    \ifnum\isQR=1
      \fill[blue!25] (\xpos-1.3, \ypos-0.8) rectangle (\xpos+1.3, \ypos+0.8);
      \node[font=\small\bfseries, blue!70!black] at (\xpos, \ypos) {\pgfmathprintnumber[int detect]{\k}};
    \else
      \fill[red!15] (\xpos-1.3, \ypos-0.8) rectangle (\xpos+1.3, \ypos+0.8);
      \node[font=\small, red!60!black] at (\xpos, \ypos) {\pgfmathprintnumber[int detect]{\k}};
    \fi
  }
  % Legend
  \fill[blue!25] (1, -1.5) rectangle (2.5, -0.5);
  \node[right, font=\small] at (2.8, -1) {QR};
  \fill[red!15] (7, -1.5) rectangle (8.5, -0.5);
  \node[right, font=\small] at (8.8, -1) {QNR};
\end{tikzpicture}
\caption{Résidus quadratiques et non-résidus modulo $29$.
Les $14$ cases bleues correspondent aux QR, les $14$ rouges aux QNR.}
\label{fig:qr-pattern-29}
\end{figure}

%--------------------------------------------------------------------
\section{Connexion avec les formes quadratiques}
\label{sec:qr:formes-quadratiques}
%--------------------------------------------------------------------

\begin{proposition}[QR et représentation par des formes quadratiques]
\label{prop:qr-formes}
\index{formes quadratiques!et résidus quadratiques}
Soit $p$ un premier impair et $d \in \Z$ avec $p \nmid d$. L'équation
$x^2 + dy^2 \equiv 0 \pmod{p}$ admet une solution non triviale (c'est-à-dire
$p \nmid x$ ou $p \nmid y$) si et seulement si $\legendre{-d}{p} = 1$.
\end{proposition}

\begin{proof}
Si $p \nmid y$, l'équation se réécrit $(xy^{-1})^2 \equiv -d \pmod{p}$,
ce qui équivaut à $\legendre{-d}{p} = 1$. Si $p \mid y$ et $p \nmid x$,
alors $x^2 \equiv 0$, contradiction. L'existence d'une solution non triviale
est donc équivalente à ce que $-d$ soit un QR modulo $p$.
\end{proof}

\begin{remarque}
\label{rem:fermat-deux-carres}
\index{Fermat!théorème des deux carrés de}
En prenant $d = 1$, on retrouve le fait que $x^2 + y^2 \equiv 0 \pmod{p}$
admet une solution non triviale si et seulement si $\legendre{-1}{p} = 1$,
soit $p \equiv 1 \pmod{4}$. C'est la première étape de la preuve du théorème de
Fermat sur les sommes de deux carrés.
\end{remarque}

%--------------------------------------------------------------------
\section{Exercices}
\label{sec:qr:exercices}
%--------------------------------------------------------------------

\begin{exercice}[$\star$]
\label{exo:qr-1}
Dresser la liste des résidus quadratiques modulo $17$.
\end{exercice}

\begin{exercice}[$\star$]
\label{exo:qr-2}
Calculer $\legendre{7}{29}$ par le critère d'Euler, puis vérifier par la
réciprocité quadratique.
\end{exercice}

\begin{exercice}[$\star$]
\label{exo:qr-3}
Calculer les symboles de Legendre suivants en utilisant la réciprocité~:
\[
  \legendre{3}{43}, \quad \legendre{-5}{37}, \quad \legendre{11}{71}, \quad \legendre{15}{97}.
\]
\end{exercice}

\begin{exercice}[$\star$]
\label{exo:qr-4}
Montrer que $-1$ est un QR modulo tout premier de la forme $4k + 1$
en exhibant explicitement une racine carrée (indication~: considérer
$((2k)!)^2$, ou plus directement $((p-1)/2)!$, modulo $p$).
\end{exercice}

\begin{exercice}[$\star\star$]
\label{exo:qr-5}
Soit $p$ un premier impair. Montrer que la somme des résidus quadratiques
modulo $p$ dans $\{1, \ldots, p-1\}$ est $p(p-1)/4$.
\end{exercice}

\begin{exercice}[$\star\star$]
\label{exo:qr-6}
Montrer que le produit de deux QNR modulo $p$ est un QR, et que le produit
d'un QR et d'un QNR est un QNR. En déduire que les QR modulo $p$ forment
un sous-groupe d'indice $2$ de $(\Z/p\Z)^\times$.
\end{exercice}

\begin{exercice}[$\star\star$]
\label{exo:qr-7}
Soit $p$ un premier $\geq 5$. Montrer qu'il existe deux résidus quadratiques
consécutifs modulo $p$ (c'est-à-dire $a$ et $a + 1$ tous deux QR).
\textit{Indication~: utiliser le fait que $\sum_{a=1}^{p-2} (1 + \legendre{a}{p})(1 + \legendre{a+1}{p}) > 0$.}
\end{exercice}

\begin{exercice}[$\star\star$]
\label{exo:qr-8}
Calculer $\jacobi{713}{1009}$ par l'algorithme du symbole de Jacobi, en
détaillant chaque étape.
\end{exercice}

\begin{exercice}[$\star\star\star$]
\label{exo:qr-9}
\index{somme de Gauss}
Soit $p$ un premier impair. La \emph{somme de Gauss} est
$g = \sum_{a=1}^{p-1} \legendre{a}{p} \, e^{2\pi i a / p}$.
Montrer que $g^2 = (-1)^{(p-1)/2} \, p = \legendre{-1}{p} \, p$.
\textit{Indication~: calculer $\abs{g}^2 = g \overline{g}$ en échangeant les
sommes.}
\end{exercice}

\begin{exercice}[$\star\star\star$]
\label{exo:qr-10}
\index{nombres premiers!infinité de $p \equiv 1 \pmod{4}$}
En utilisant le premier supplément, montrer qu'il existe une infinité de
nombres premiers $p \equiv 1 \pmod{4}$.
\textit{Indication~: considérer $(2 n!)^2 + 1$ pour un entier $n$ convenable
et examiner ses facteurs premiers.}
\end{exercice}

\begin{exercice}[$\star\star\star$]
\label{exo:qr-11}
\index{réciprocité quadratique!preuve par les sommes de Gauss}
Donner une preuve de la loi de réciprocité quadratique en utilisant
les sommes de Gauss.
\textit{Indication~: montrer que $g^p \equiv \legendre{p}{q} \cdot g \pmod{p}$
et utiliser le fait que $g^2 = (-1)^{(q-1)/2} q$.}
\end{exercice}

%--------------------------------------------------------------------
\section*{Résumé du chapitre}
\addcontentsline{toc}{section}{Résumé du chapitre}
%--------------------------------------------------------------------

\begin{itemize}
  \item Un entier $a$ non divisible par $p$ est un \textbf{résidu quadratique}
        (QR) modulo $p$ s'il existe $x$ tel que $x^2 \equiv a \pmod{p}$.
        Il y a exactement $(p-1)/2$ QR et $(p-1)/2$ QNR modulo $p$.

  \item Le \textbf{critère d'Euler} affirme que
        $a^{(p-1)/2} \equiv \legendre{a}{p} \pmod{p}$.

  \item Le \textbf{symbole de Legendre} $\legendre{a}{p}$ vaut $0$, $+1$ ou $-1$
        selon que $p \mid a$, $a$ est QR, ou $a$ est QNR.

  \item Le \textbf{lemme de Gauss} relie $\legendre{a}{p}$ au nombre de
        <<~dépassements~>> dans la multiplication par $a$ des éléments de
        $\{1, \ldots, (p-1)/2\}$.

  \item Les \textbf{suppléments} donnent~:
        $\legendre{-1}{p} = (-1)^{(p-1)/2}$ et
        $\legendre{2}{p} = (-1)^{(p^2-1)/8}$.

  \item La \textbf{loi de réciprocité quadratique}~:
        $\legendre{p}{q} \legendre{q}{p} = (-1)^{\frac{p-1}{2}\cdot\frac{q-1}{2}}$
        pour $p, q$ premiers impairs distincts.

  \item Le \textbf{symbole de Jacobi} $\jacobi{a}{n}$ généralise le symbole de
        Legendre aux modules composés impairs et satisfait la même loi de
        réciprocité. Il se calcule efficacement sans factorisation.
\end{itemize}

% === FIN DU CHUNK ch5-6 ===
% === DÉBUT DU CHUNK ch7-8 ===

%%%%%%%%%%%%%%%%%%%%%%%%%%%%%%%%%%%%%%%%%%%%%%%%%%%%%%%%%%%%%%%%%%%%
%  CHAPITRE 7 — Représentations par des Formes Quadratiques
%%%%%%%%%%%%%%%%%%%%%%%%%%%%%%%%%%%%%%%%%%%%%%%%%%%%%%%%%%%%%%%%%%%%
\chapter{Représentations par des Formes Quadratiques}
\label{chap:formes-quadratiques}
\index{forme quadratique}

\section*{Introduction historique}
\addcontentsline{toc}{section}{Introduction historique}

L'étude des représentations d'entiers par des formes quadratiques constitue l'un des
chapitres les plus anciens et les plus riches de la théorie des nombres. Dès le
\textsc{iii}\textsuperscript{e} siècle, \textbf{Diophante d'Alexandrie}\index{Diophante}
s'interrogeait sur les entiers pouvant s'écrire comme sommes de deux carrés. La question
resta largement ouverte jusqu'au \textsc{xvii}\textsuperscript{e} siècle, lorsque
\textbf{Pierre de Fermat}\index{Fermat} annonça, sans démonstration complète, que tout
nombre premier $p \equiv 1 \pmod{4}$ est somme de deux carrés. C'est
\textbf{Leonhard Euler}\index{Euler} qui, après des décennies d'efforts, fournit en 1749
la première preuve rigoureuse de ce résultat fondamental.

\textbf{Joseph-Louis Lagrange}\index{Lagrange} étendit considérablement la théorie~: il
démontra en 1770 que tout entier naturel est somme de quatre carrés, et développa une
théorie systématique des formes quadratiques binaires. Cette théorie fut ensuite portée à
un haut degré de perfection par \textbf{Carl Friedrich Gauss}\index{Gauss} dans ses
\emph{Disquisitiones Arithmeticae} (1801), où il introduisit les notions de classe
d'équivalence de formes et de composition.

%% ─── 7.1 Formes quadratiques binaires ─────────────────────────────
\section{Formes quadratiques binaires}
\label{sec:formes-binaires}

\begin{definition}[Forme quadratique binaire]
\label{def:forme-quadratique-binaire}
\index{forme quadratique!binaire}
Une \emph{forme quadratique binaire} est un polynôme de la forme
\[
  f(x,y) = ax^2 + bxy + cy^2
\]
où $a, b, c \in \Z$. On note cette forme $(a,b,c)$. Le \emph{discriminant}\index{discriminant!d'une forme quadratique} de $f$ est
\[
  \Delta = b^2 - 4ac.
\]
On dit que $f$ \emph{représente} un entier $n$ s'il existe $x_0, y_0 \in \Z$ tels que
$f(x_0, y_0) = n$. Si $\pgcd(x_0, y_0) = 1$, la représentation est dite
\emph{propre}\index{representation@représentation!propre}.
\end{definition}

\begin{exemple}[Formes quadratiques classiques]
\label{ex:formes-classiques}
\begin{enumerate}[label=(\alph*)]
  \item La forme $f(x,y) = x^2 + y^2$ a pour discriminant $\Delta = -4$. Elle représente
    $n$ si et seulement si $n$ est somme de deux carrés.
  \item La forme $f(x,y) = x^2 + xy + y^2$ a pour discriminant $\Delta = -3$. Elle
    représente les entiers de la forme $x^2 + xy + y^2$, liés aux nombres premiers
    $p \equiv 1 \pmod{3}$.
  \item La forme $f(x,y) = x^2 - y^2 = (x+y)(x-y)$ a pour discriminant $\Delta = 4$.
    Elle représente toute différence de deux carrés.
\end{enumerate}
\end{exemple}

\begin{definition}[Forme définie positive]
\label{def:forme-definie-positive}
\index{forme quadratique!définie positive}
Une forme quadratique binaire $(a,b,c)$ est dite \emph{définie positive} si $a > 0$ et
$\Delta = b^2 - 4ac < 0$. Dans ce cas, $f(x,y) \geq 0$ pour tout $(x,y) \in \Z^2$, avec
égalité si et seulement si $(x,y) = (0,0)$.
\end{definition}

\begin{remarque}
\label{rem:discriminant-negatif}
Pour une forme définie positive $(a,b,c)$ avec $\Delta < 0$, on a nécessairement $c > 0$
puisque $4ac = b^2 - \Delta > 0$ et $a > 0$.
\end{remarque}

%% ─── 7.2 Sommes de deux carrés ─────────────────────────────────────
\section{Sommes de deux carrés}
\label{sec:sommes-deux-carres}
\index{somme de deux carrés}

\subsection{Identité de Brahmagupta}

L'étude des sommes de deux carrés repose sur une identité multiplicative fondamentale.

\begin{lemme}[Identité de Brahmagupta--Fibonacci]
\label{lem:brahmagupta}
\index{identité!de Brahmagupta--Fibonacci}
Pour tous $a, b, c, d \in \Z$,
\[
  (a^2 + b^2)(c^2 + d^2) = (ac - bd)^2 + (ad + bc)^2
    = (ac + bd)^2 + (ad - bc)^2.
\]
\end{lemme}

\begin{proof}
Par développement direct~:
\begin{align*}
  (ac - bd)^2 + (ad + bc)^2
    &= a^2c^2 - 2abcd + b^2d^2 + a^2d^2 + 2abcd + b^2c^2 \\
    &= a^2(c^2 + d^2) + b^2(c^2 + d^2) \\
    &= (a^2 + b^2)(c^2 + d^2). \qedhere
\end{align*}
\end{proof}

\begin{corollaire}
\label{cor:produit-sommes-carres}
Le produit de deux sommes de deux carrés est une somme de deux carrés.
\end{corollaire}

\begin{remarque}
\label{rem:brahmagupta-gaussiens}
L'identité de Brahmagupta s'interprète naturellement dans l'anneau des entiers de
Gauss\index{entiers de Gauss} $\Z[i]$~: si $\alpha = a + bi$ et $\beta = c + di$, alors
$\abs{\alpha}^2 \abs{\beta}^2 = \abs{\alpha\beta}^2$, ce qui n'est autre que la
multiplicativité de la norme.
\end{remarque}

\subsection{Entiers de Gauss}

Pour la preuve du théorème de Fermat sur les sommes de deux carrés, nous utiliserons
l'anneau des entiers de Gauss.

\begin{definition}[Anneau des entiers de Gauss]
\label{def:entiers-gauss}
\index{entiers de Gauss}
L'\emph{anneau des entiers de Gauss} est
\[
  \Z[i] = \{a + bi : a, b \in \Z\} \subset \C.
\]
La \emph{norme}\index{norme!dans $\Z[i]$} d'un élément $\alpha = a + bi$ est
$N(\alpha) = a^2 + b^2 = \alpha \bar{\alpha} \in \N$.
\end{definition}

\begin{proposition}[{Propriétés de $\Z[i]$}]
\label{prop:proprietes-Z-i}
\begin{enumerate}[label=(\roman*)]
  \item La norme est multiplicative~: $N(\alpha\beta) = N(\alpha)N(\beta)$.
  \item Les inversibles de $\Z[i]$ sont $\{1, -1, i, -i\}$, c'est-à-dire les éléments de
    norme~$1$.
  \item L'anneau $\Z[i]$ est un anneau euclidien pour la norme $N$, et donc un anneau
    principal et factoriel.
\end{enumerate}
\end{proposition}

\begin{proof}
\begin{enumerate}[label=(\roman*)]
  \item $N(\alpha\beta) = \alpha\beta\overline{\alpha\beta}
    = \alpha\bar\alpha\cdot\beta\bar\beta = N(\alpha)N(\beta)$.
  \item Si $\alpha$ est inversible, $N(\alpha)N(\alpha^{-1}) = N(1) = 1$, donc
    $N(\alpha) = 1$, d'où $a^2 + b^2 = 1$, ce qui donne $(a,b) \in \{(\pm 1, 0),
    (0, \pm 1)\}$.
  \item Soient $\alpha, \beta \in \Z[i]$ avec $\beta \neq 0$. On écrit
    $\alpha/\beta = u + vi$ avec $u, v \in \Q$. Soient $m, n \in \Z$ les entiers les plus
    proches de $u$ et $v$. Posons $\gamma = m + ni$ et $\rho = \alpha - \beta\gamma$.
    Alors $N(\rho) = N(\beta)N(\alpha/\beta - \gamma) = N(\beta)\bigl((u-m)^2 +
    (v-n)^2\bigr) \leq N(\beta)(1/4 + 1/4) = N(\beta)/2 < N(\beta)$. \qedhere
\end{enumerate}
\end{proof}

\begin{lemme}[{Factorisation des premiers dans $\Z[i]$}]
\label{lem:premiers-gaussiens}
\index{nombre premier!dans $\Z[i]$}
Soit $p$ un nombre premier (de $\Z$).
\begin{enumerate}[label=(\roman*)]
  \item Si $p = 2$, alors $p = -i(1+i)^2$, et $1+i$ est irréductible dans $\Z[i]$.
  \item Si $p \equiv 1 \pmod{4}$, alors $p = \pi\bar\pi$ où $\pi$ est irréductible dans
    $\Z[i]$ et $\pi \neq \bar\pi$ (à association près).
  \item Si $p \equiv 3 \pmod{4}$, alors $p$ reste irréductible dans $\Z[i]$.
\end{enumerate}
\end{lemme}

\begin{proof}
Si $p$ n'est pas irréductible dans $\Z[i]$, alors $p = \alpha\beta$ avec
$1 < N(\alpha), N(\beta) < p^2$, donc $N(\alpha) = N(\beta) = p$, ce qui signifie que $p$
est somme de deux carrés.

Pour $p = 2$~: $2 = 1^2 + 1^2$, et $2 = -i(1+i)^2$.

Pour $p \equiv 3 \pmod{4}$~: si $p = a^2 + b^2$, alors $a^2 + b^2 \equiv 3 \pmod{4}$.
Or $a^2 \equiv 0$ ou $1$ modulo $4$, de même pour $b^2$, donc $a^2 + b^2 \equiv 0, 1$
ou $2$ modulo $4$. Contradiction.

Pour $p \equiv 1 \pmod{4}$~: par le théorème de Wilson, $(p-1)! \equiv -1 \pmod{p}$.
Posons $m = (p-1)/2$. Alors
\[
  (p-1)! = 1 \cdot 2 \cdots m \cdot (m+1) \cdots (p-1).
\]
Or $p - k \equiv -k \pmod{p}$, donc $(p-1)! \equiv (-1)^m (m!)^2 \pmod{p}$. Puisque
$p \equiv 1 \pmod{4}$, $m$ est pair, et $(-1)^m = 1$, d'où $(m!)^2 \equiv -1 \pmod{p}$.
Posons $r = m!$. Alors $p \mid r^2 + 1 = (r+i)(r-i)$ dans $\Z[i]$. Si $p$ était
irréductible dans $\Z[i]$, il diviserait $r+i$ ou $r-i$, ce qui est impossible car
$r/p \notin \Z$. Donc $p$ n'est pas irréductible dans $\Z[i]$, et l'argument initial
montre $p = \pi\bar\pi$ avec $N(\pi) = p$.
\end{proof}

\subsection{Théorème de Fermat sur les sommes de deux carrés}

\begin{theoreme}[Fermat, 1640 -- Euler, 1749]
\label{thm:fermat-deux-carres}
\index{théorème!de Fermat (deux carrés)}
Un nombre premier $p$ est somme de deux carrés d'entiers si et seulement si $p = 2$ ou
$p \equiv 1 \pmod{4}$.
\end{theoreme}

\begin{proof}
La preuve découle directement du lemme~\ref{lem:premiers-gaussiens}. Si $p = 2$, alors
$2 = 1^2 + 1^2$. Si $p \equiv 1 \pmod{4}$, alors $p = \pi\bar\pi$ dans $\Z[i]$ avec
$\pi = a + bi$, d'où $p = N(\pi) = a^2 + b^2$. Si $p \equiv 3 \pmod{4}$, alors $p$
est irréductible dans $\Z[i]$, donc $p$ ne peut s'écrire $a^2 + b^2$ avec $a, b \in \Z$.
\end{proof}

\begin{theoreme}[Caractérisation des sommes de deux carrés]
\label{thm:caract-sommes-deux-carres}
\index{somme de deux carrés!caractérisation}
Un entier $n \geq 1$ est somme de deux carrés si et seulement si, dans la factorisation
$n = \prod p_i^{e_i}$, tout facteur premier $p_i \equiv 3 \pmod{4}$ apparaît avec un
exposant $e_i$ pair.
\end{theoreme}

\begin{proof}
\emph{Condition suffisante.} Si tout premier $p \equiv 3 \pmod{4}$ divisant $n$ apparaît
à une puissance paire, alors $n = m^2 \cdot n'$ où $n'$ n'a que des facteurs premiers
$p = 2$ ou $p \equiv 1 \pmod{4}$. Chacun de ces facteurs est somme de deux carrés, et
par l'identité de Brahmagupta (lemme~\ref{lem:brahmagupta}), $n'$ l'est aussi. Donc
$n = m^2 n' = m^2(a^2 + b^2) = (ma)^2 + (mb)^2$.

\emph{Condition nécessaire.} Soit $p \equiv 3 \pmod{4}$ et supposons $p^{2k+1} \| n$ avec
$n = a^2 + b^2$. Dans $\Z[i]$, on a $n = (a+bi)(a-bi)$. Puisque $p$ est irréductible
dans $\Z[i]$ (lemme~\ref{lem:premiers-gaussiens}), $p$ divise $a+bi$ et $a-bi$ le même
nombre de fois. Si $p^r \mid (a+bi)$ et $p^s \mid (a-bi)$ dans $\Z[i]$, alors
$p^{r+s} \mid n$ dans $\Z$, et $r = s$ car $\overline{a+bi} = a-bi$. Donc $p$ divise
$n$ à une puissance paire.
\end{proof}

\begin{exemple}[Application du critère]
\label{ex:application-critere-deux-carres}
\begin{enumerate}[label=(\alph*)]
  \item $n = 50 = 2 \cdot 5^2$. Les facteurs premiers $2$ et $5$ vérifient
    $2 = 1^2 + 1^2$ et $5 \equiv 1 \pmod{4}$. Pas de facteur $\equiv 3 \pmod{4}$, donc
    $50$ est somme de deux carrés~: $50 = 1^2 + 7^2 = 5^2 + 5^2$.
  \item $n = 21 = 3 \cdot 7$. Les deux facteurs premiers sont $\equiv 3 \pmod{4}$ et
    apparaissent à la puissance~$1$ (impaire). Donc $21$ n'est pas somme de deux carrés.
  \item $n = 45 = 3^2 \cdot 5$. Le seul facteur $\equiv 3 \pmod{4}$ est $3$, qui
    apparaît au carré. Donc $45$ est somme de deux carrés~: $45 = 3^2 + 6^2$.
\end{enumerate}
\end{exemple}

\subsection{Dénombrement des représentations}

\begin{definition}[Fonction $r_2(n)$]
\label{def:r2}
\index{r2n@$r_2(n)$}
Pour $n \geq 0$, on note $r_2(n)$ le nombre de couples $(a,b) \in \Z^2$ tels que
$a^2 + b^2 = n$. On compte les représentations \emph{ordonnées} et avec signes.
\end{definition}

\begin{theoreme}[Jacobi]
\label{thm:jacobi-r2}
\index{théorème!de Jacobi ($r_2$)}
Pour tout $n \geq 1$,
\[
  r_2(n) = 4\sum_{d \mid n} \chi(d)
    = 4\bigl(d_1(n) - d_3(n)\bigr),
\]
où $\chi$ est le caractère de Dirichlet non principal modulo~$4$ (défini par $\chi(d) = 0$
si $d$ est pair, $\chi(d) = 1$ si $d \equiv 1 \pmod{4}$, $\chi(d) = -1$ si
$d \equiv 3 \pmod{4}$), et $d_k(n)$ désigne le nombre de diviseurs de $n$ congrus à $k$
modulo $4$.
\end{theoreme}

\begin{exemple}
\label{ex:r2-calculs}
\begin{enumerate}[label=(\alph*)]
  \item $r_2(5) = 4(d_1(5) - d_3(5)) = 4(2 - 0) = 8$. En effet, $5 = 1^2 + 2^2$ donne
    les huit couples $(\pm 1, \pm 2)$ et $(\pm 2, \pm 1)$.
  \item $r_2(25) = 4(d_1(25) - d_3(25)) = 4(3 - 0) = 12$. Les diviseurs de $25$ congrus
    à $1 \pmod{4}$ sont $1, 5, 25$.
\end{enumerate}
\end{exemple}

%% ─── 7.3 Figure TikZ : points entiers sur un cercle ──────────────
\subsection{Visualisation : points entiers sur des cercles}

La figure suivante illustre les points $(a,b) \in \Z^2$ tels que $a^2 + b^2 = n$ pour
différentes valeurs de $n$.

\begin{figure}[ht]
\centering
\begin{tikzpicture}[scale=0.38]
  % Grille
  \draw[gray!20, very thin] (-8,-8) grid (8,8);
  \draw[->] (-8.5,0) -- (8.5,0) node[right] {$a$};
  \draw[->] (0,-8.5) -- (0,8.5) node[above] {$b$};

  % Cercle n=5
  \draw[blue!60, thick] (0,0) circle ({sqrt(5)});
  \foreach \a/\b in {1/2, 2/1, -1/2, -2/1, 1/-2, 2/-1, -1/-2, -2/-1} {
    \fill[blue!80] (\a,\b) circle (4pt);
  }
  \node[blue!80] at (3.2,2.2) {\small $n=5$};

  % Cercle n=25
  \draw[red!60, thick] (0,0) circle (5);
  \foreach \a/\b in {0/5, 5/0, 0/-5, -5/0, 3/4, 4/3, -3/4, -4/3,
                      3/-4, 4/-3, -3/-4, -4/-3} {
    \fill[red!70] (\a,\b) circle (4pt);
  }
  \node[red!70] at (5.5,3.5) {\small $n=25$};

  % Cercle n=50
  \draw[green!50!black, thick] (0,0) circle ({sqrt(50)});
  \foreach \a/\b in {1/7, 7/1, -1/7, -7/1, 1/-7, 7/-1, -1/-7, -7/-1,
                      5/5, -5/5, 5/-5, -5/-5} {
    \fill[green!60!black] (\a,\b) circle (4pt);
  }
  \node[green!60!black] at (7.2,4.5) {\small $n=50$};
\end{tikzpicture}
\caption{Points entiers sur les cercles $a^2 + b^2 = n$ pour $n \in \{5, 25, 50\}$.}
\label{fig:points-cercles}
\end{figure}

%% ─── 7.4 Entiers de Gauss : visualisation ─────────────────────────
\begin{figure}[ht]
\centering
\begin{tikzpicture}[scale=0.7]
  % Grille des entiers de Gauss
  \draw[gray!15, very thin] (-4,-4) grid (4,4);
  \draw[->] (-4.3,0) -- (4.5,0) node[right] {$\Re$};
  \draw[->] (0,-4.3) -- (0,4.5) node[above] {$\Im$};

  % Points du réseau
  \foreach \a in {-4,...,4} {
    \foreach \b in {-4,...,4} {
      \fill[black!30] (\a,\b) circle (1.5pt);
    }
  }

  % Premier gaussien pi = 2+i, norme 5
  \fill[blue!80] (2,1) circle (3.5pt) node[above right] {\small $2+i$};
  \fill[blue!80] (2,-1) circle (3.5pt) node[below right] {\small $2-i$};
  \fill[red!80] (1,2) circle (3.5pt) node[above right] {\small $1+2i$};
  \fill[red!80] (1,-2) circle (3.5pt) node[below right] {\small $1-2i$};

  % Inversibles
  \foreach \a/\b/\lab in {1/0/1, -1/0/{-1}, 0/1/i, 0/-1/{-i}} {
    \fill[orange!80!black] (\a,\b) circle (3pt);
  }
  \node[orange!80!black] at (1.3,-0.4) {\small $1$};
  \node[orange!80!black] at (0.5,1) {\small $i$};
\end{tikzpicture}
\caption{L'anneau des entiers de Gauss $\Z[i]$ : inversibles (orange), premiers
associés à $5 = (2+i)(2-i)$ (bleu et rouge).}
\label{fig:entiers-gauss}
\end{figure}

%% ─── 7.5 Sommes de quatre carrés ──────────────────────────────────
\section{Sommes de quatre carrés}
\label{sec:sommes-quatre-carres}
\index{somme de quatre carrés}
\index{théorème!de Lagrange (quatre carrés)}

\begin{theoreme}[Lagrange, 1770]
\label{thm:lagrange-quatre-carres}
Tout entier naturel est somme de quatre carrés d'entiers~:
\[
  \forall n \in \N, \quad \exists\, a,b,c,d \in \Z : \quad
    n = a^2 + b^2 + c^2 + d^2.
\]
\end{theoreme}

La preuve repose sur deux ingrédients~: une identité multiplicative et un argument
arithmétique.

\begin{lemme}[Identité d'Euler pour les quaternions]
\label{lem:identite-euler-quaternions}
\index{identité!d'Euler (quatre carrés)}
Pour tous $a_1, \ldots, a_4, b_1, \ldots, b_4 \in \Z$,
\[
  (a_1^2 + a_2^2 + a_3^2 + a_4^2)(b_1^2 + b_2^2 + b_3^2 + b_4^2)
    = c_1^2 + c_2^2 + c_3^2 + c_4^2,
\]
où \vspace{-2mm}
\begin{align*}
  c_1 &= a_1 b_1 - a_2 b_2 - a_3 b_3 - a_4 b_4, &
  c_2 &= a_1 b_2 + a_2 b_1 + a_3 b_4 - a_4 b_3, \\
  c_3 &= a_1 b_3 - a_2 b_4 + a_3 b_1 + a_4 b_2, &
  c_4 &= a_1 b_4 + a_2 b_3 - a_3 b_2 + a_4 b_1.
\end{align*}
\end{lemme}

\begin{proof}
On identifie $(a_1, a_2, a_3, a_4)$ au quaternion\index{quaternion}
$\alpha = a_1 + a_2\,\mathbf{i} + a_3\,\mathbf{j} + a_4\,\mathbf{k}$, et de même pour
$\beta$. Alors $N(\alpha) = a_1^2 + a_2^2 + a_3^2 + a_4^2$, et la multiplicativité
de la norme des quaternions donne $N(\alpha\beta) = N(\alpha)N(\beta)$. Le produit
$\alpha\beta$ a pour composantes exactement $c_1, c_2, c_3, c_4$.
\end{proof}

\begin{proof}[Preuve du théorème~\ref{thm:lagrange-quatre-carres} (schéma)]
Grâce au lemme~\ref{lem:identite-euler-quaternions}, il suffit de montrer que tout nombre
premier $p$ est somme de quatre carrés.

\emph{Étape 1.} On montre qu'il existe $m$ avec $1 \leq m < p$ tel que $mp$ est somme de
quatre carrés. En effet, l'équation $x^2 + y^2 + 1 \equiv 0 \pmod{p}$ admet une solution
car l'ensemble $\{x^2 \bmod p : x = 0, 1, \ldots, (p-1)/2\}$ a $(p+1)/2$ éléments, et
$\{-1 - y^2 \bmod p : y = 0, 1, \ldots, (p-1)/2\}$ en a autant. Par le principe des
tiroirs, ces ensembles ont une intersection non vide~: il existe $x_0, y_0$ tels que
$x_0^2 + y_0^2 + 1 \equiv 0 \pmod{p}$, soit $mp = x_0^2 + y_0^2 + 1^2 + 0^2$ avec
$m < p$.

\emph{Étape 2.} On montre par descente infinie que si $mp = a_1^2 + a_2^2 + a_3^2 +
a_4^2$ avec $m > 1$, on peut trouver $m'$ avec $1 \leq m' < m$ tel que $m'p$ est aussi
somme de quatre carrés. On choisit $b_i \equiv a_i \pmod{m}$ avec $\abs{b_i} \leq m/2$,
et on utilise l'identité d'Euler pour réduire le multiplicateur.

Par descente, on atteint $m = 1$, c'est-à-dire $p = a_1^2 + a_2^2 + a_3^2 + a_4^2$.
\end{proof}

\begin{exemple}
\label{ex:sommes-quatre-carres}
\begin{enumerate}[label=(\alph*)]
  \item $7 = 4 + 1 + 1 + 1 = 2^2 + 1^2 + 1^2 + 1^2$.
  \item $15 = 9 + 4 + 1 + 1 = 3^2 + 2^2 + 1^2 + 1^2$.
  \item $23 = 9 + 9 + 4 + 1 = 3^2 + 3^2 + 2^2 + 1^2$.
\end{enumerate}
\end{exemple}

%% ─── 7.6 Problème de Waring ────────────────────────────────────────
\section{Le problème de Waring}
\label{sec:probleme-waring}
\index{problème de Waring}
\index{Waring}

\begin{theoreme}[Hilbert, 1909]
\label{thm:hilbert-waring}
Pour tout entier $k \geq 2$, il existe un entier $g(k)$ tel que tout entier naturel est
somme de $g(k)$ puissances $k$-ièmes d'entiers naturels.
\end{theoreme}

\begin{remarque}
\label{rem:valeurs-gk}
Les premières valeurs connues de $g(k)$ sont~:
\begin{center}
\begin{tabular}{c|cccccc}
\toprule
$k$ & 2 & 3 & 4 & 5 & 6 & 7 \\
\midrule
$g(k)$ & 4 & 9 & 19 & 37 & 73 & 143 \\
\bottomrule
\end{tabular}
\end{center}
Le théorème de Lagrange donne $g(2) = 4$. La valeur $g(3) = 9$ fut établie par
Wieferich\index{Wieferich} et Kempner\index{Kempner}. On conjecture (et c'est démontré
pour $k$ assez grand) que
\[
  g(k) = 2^k + \floor*{\left(\frac{3}{2}\right)^k} - 2.
\]
\end{remarque}

%% ─── 7.7 Formes ternaires et théorème de Legendre ─────────────────
\section{Formes ternaires et théorème de Legendre}
\label{sec:formes-ternaires}
\index{forme quadratique!ternaire}

\begin{theoreme}[Legendre, 1785]
\label{thm:legendre-ternaire}
\index{théorème!de Legendre (ternaire)}
L'équation $ax^2 + by^2 + cz^2 = 0$ avec $a, b, c \in \Z$ non tous de même signe et
sans facteur carré commun admet une solution non triviale $(x,y,z) \in \Z^3 \setminus
\{(0,0,0)\}$ si et seulement si~:
\begin{enumerate}[label=(\roman*)]
  \item $-bc$ est un résidu quadratique modulo $\abs{a}$,
  \item $-ac$ est un résidu quadratique modulo $\abs{b}$,
  \item $-ab$ est un résidu quadratique modulo $\abs{c}$.
\end{enumerate}
\end{theoreme}

\begin{corollaire}[Théorème des trois carrés -- Gauss--Legendre]
\label{cor:trois-carres}
\index{théorème!des trois carrés}
Un entier $n \geq 0$ est somme de trois carrés si et seulement si $n$ n'est pas de la
forme $4^a(8b+7)$ avec $a, b \in \N$.
\end{corollaire}

%% ─── 7.8 Équivalence et réduction de formes ───────────────────────
\section{Équivalence et réduction de formes}
\label{sec:equivalence-reduction}
\index{forme quadratique!équivalence}
\index{forme quadratique!réduite}

\begin{definition}[Équivalence de formes]
\label{def:equivalence-formes}
Deux formes quadratiques binaires $f$ et $g$ sont dites \emph{proprement équivalentes}
s'il existe une matrice $\begin{psmallmatrix} \alpha & \beta \\ \gamma & \delta
\end{psmallmatrix} \in \mathrm{SL}_2(\Z)$ telle que
\[
  g(x,y) = f(\alpha x + \beta y,\; \gamma x + \delta y).
\]
Des formes équivalentes ont le même discriminant et représentent les mêmes entiers.
\end{definition}

\begin{definition}[Forme réduite au sens de Gauss]
\label{def:forme-reduite}
Une forme définie positive $(a,b,c)$ de discriminant $\Delta < 0$ est \emph{réduite} si
\[
  \abs{b} \leq a \leq c,
\]
et de plus $b \geq 0$ lorsque $\abs{b} = a$ ou $a = c$.
\end{definition}

\begin{theoreme}
\label{thm:unicite-forme-reduite}
Toute forme définie positive de discriminant $\Delta < 0$ est proprement équivalente à
une unique forme réduite.
\end{theoreme}

\begin{exemple}
\label{ex:formes-reduites}
Pour $\Delta = -4$, la seule forme réduite est $(1, 0, 1)$, c'est-à-dire $x^2 + y^2$.
Pour $\Delta = -20$, les formes réduites sont $(1, 0, 5)$ et $(2, 2, 3)$. Le
nombre de classes\index{nombre de classes} est $h(-20) = 2$.
\end{exemple}

%% ─── 7.9 Exercices ────────────────────────────────────────────────
\section{Exercices}
\label{sec:exos-ch7}

\begin{exercice}[$\star$]
\label{exo:sommes-deux-carres-1}
Déterminer lesquels des entiers suivants sont sommes de deux carrés, et trouver une
représentation le cas échéant~: $n = 65$, $n = 99$, $n = 130$, $n = 225$, $n = 1000$.
\end{exercice}

\begin{exercice}[$\star$]
\label{exo:r2-calcul}
Calculer $r_2(n)$ pour $n = 1, 2, 4, 5, 10, 13, 25$.
\end{exercice}

\begin{exercice}[$\star$]
\label{exo:quatre-carres}
Écrire chacun des entiers $30$, $47$, $100$, $200$ comme somme de quatre carrés.
\end{exercice}

\begin{exercice}[$\star\star$]
\label{exo:produit-sommes}
Montrer directement (sans utiliser les entiers de Gauss) que si $p$ est un premier avec
$p \equiv 1 \pmod{4}$, alors il existe $m$ avec $1 \leq m < p$ tel que $mp = a^2 + b^2$
pour certains $a, b \in \Z$.
\end{exercice}

\begin{exercice}[$\star\star$]
\label{exo:infini-sommes}
Montrer qu'il existe une infinité de nombres premiers $p$ tels que $p = a^2 + b^2$ avec
$a, b \in \N^*$.
\textit{Indication~: utiliser le théorème de Dirichlet sur les progressions arithmétiques.}
\end{exercice}

\begin{exercice}[$\star\star$]
\label{exo:trois-carres-exclusion}
Montrer qu'un entier de la forme $4^a(8b+7)$ ne peut pas être somme de trois carrés.
\textit{Indication~: raisonner modulo $8$.}
\end{exercice}

\begin{exercice}[$\star\star$]
\label{exo:forme-reduction}
Trouver la forme réduite équivalente à $f(x,y) = 3x^2 + 10xy + 9y^2$.
\end{exercice}

\begin{exercice}[$\star\star\star$]
\label{exo:nombre-classes}
Montrer que le nombre de classes $h(\Delta)$ est fini pour tout discriminant $\Delta < 0$.
\textit{Indication~: montrer que pour une forme réduite $(a,b,c)$, on a
$a \leq \sqrt{\abs{\Delta}/3}$.}
\end{exercice}

\begin{exercice}[$\star\star\star$]
\label{exo:jacobi-r4}
Soient $r_4(n)$ le nombre de représentations de $n$ comme somme de quatre carrés. Montrer
que $r_4(n) = 8\sum_{\substack{d \mid n \\ 4 \nmid d}} d$. En déduire que $r_4(n) \geq 1$
pour tout $n \geq 1$.
\end{exercice}

%% ─── 7.10 Résumé du chapitre ──────────────────────────────────────
\section*{Résumé du chapitre}
\addcontentsline{toc}{section}{Résumé du chapitre}

\begin{enumerate}[label=\textbf{\arabic*.}]
  \item Une \textbf{forme quadratique binaire} $f(x,y) = ax^2 + bxy + cy^2$ a pour
    discriminant $\Delta = b^2 - 4ac$.
  \item L'\textbf{identité de Brahmagupta} montre que le produit de deux sommes de deux
    carrés est une somme de deux carrés.
  \item L'anneau des \textbf{entiers de Gauss} $\Z[i]$ est un anneau euclidien, et la
    factorisation des premiers dans cet anneau détermine quels entiers sont sommes de
    deux carrés.
  \item \textbf{Théorème de Fermat}~: $p$ premier est somme de deux carrés ssi $p = 2$
    ou $p \equiv 1 \pmod{4}$.
  \item Un entier $n$ est somme de deux carrés ssi tout facteur premier
    $p \equiv 3 \pmod{4}$ apparaît avec un exposant pair.
  \item \textbf{Théorème de Lagrange}~: tout entier est somme de quatre carrés.
  \item Le \textbf{problème de Waring}~: tout entier est somme de $g(k)$ puissances
    $k$-ièmes.
  \item Deux formes de même discriminant sont proprement équivalentes si et seulement si
    elles ont la même forme réduite.
\end{enumerate}

\newpage

%%%%%%%%%%%%%%%%%%%%%%%%%%%%%%%%%%%%%%%%%%%%%%%%%%%%%%%%%%%%%%%%%%%%
%  CHAPITRE 8 — Fonctions Arithmétiques
%%%%%%%%%%%%%%%%%%%%%%%%%%%%%%%%%%%%%%%%%%%%%%%%%%%%%%%%%%%%%%%%%%%%
\chapter{Fonctions Arithmétiques}
\label{chap:fonctions-arithmetiques}
\index{fonction arithmétique}

\section*{Introduction}
\addcontentsline{toc}{section}{Introduction}

Les fonctions arithmétiques --- fonctions de $\N^*$ dans $\C$ --- sont omniprésentes en
théorie des nombres. Elles codent des propriétés fondamentales des entiers~: nombre de
diviseurs, somme des diviseurs, structure multiplicative. Ce chapitre développe la théorie
algébrique de ces fonctions, centrée sur le \emph{produit de Dirichlet}, et culmine avec
la \emph{formule d'inversion de Möbius}\index{inversion de Möbius}, l'un des outils les
plus puissants et les plus élégants de la théorie des nombres élémentaire.

%% ─── 8.1 Définitions fondamentales ────────────────────────────────
\section{Définitions et exemples fondamentaux}
\label{sec:fonctions-definitions}

\begin{definition}[Fonction arithmétique]
\label{def:fonction-arithmetique}
\index{fonction arithmétique!définition}
Une \emph{fonction arithmétique} est une application $f : \N^* \to \C$.
\end{definition}

\begin{definition}[Multiplicativité]
\label{def:multiplicativite}
\index{fonction arithmétique!multiplicative}
\index{fonction arithmétique!complètement multiplicative}
Une fonction arithmétique $f$ est~:
\begin{enumerate}[label=(\roman*)]
  \item \emph{multiplicative} si $f(1) = 1$ et $f(mn) = f(m)f(n)$ pour tout
    $m, n \in \N^*$ avec $\pgcd(m,n) = 1$;
  \item \emph{complètement multiplicative} si $f(1) = 1$ et $f(mn) = f(m)f(n)$ pour tout
    $m, n \in \N^*$.
\end{enumerate}
\end{definition}

\begin{remarque}
\label{rem:multiplicative-determination}
Une fonction multiplicative est entièrement déterminée par ses valeurs sur les puissances
de nombres premiers $f(p^k)$. Une fonction complètement multiplicative est déterminée par
ses valeurs sur les nombres premiers $f(p)$.
\end{remarque}

%% ─── 8.2 Indicatrice d'Euler ──────────────────────────────────────
\section{L'indicatrice d'Euler $\varphi(n)$}
\label{sec:euler-phi}
\index{indicatrice d'Euler}
\index{phi@$\varphi(n)$}

\begin{definition}[Indicatrice d'Euler]
\label{def:euler-phi}
Pour $n \geq 1$, l'\emph{indicatrice d'Euler} est
\[
  \varphi(n) = \#\{k \in \{1, 2, \ldots, n\} : \pgcd(k, n) = 1\}.
\]
\end{definition}

\begin{theoreme}[Formule produit pour $\varphi$]
\label{thm:formule-phi}
Pour tout $n \geq 1$,
\[
  \varphi(n) = n \prod_{p \mid n} \left(1 - \frac{1}{p}\right).
\]
En particulier, si $n = p_1^{a_1} \cdots p_r^{a_r}$, alors
\[
  \varphi(n) = \prod_{i=1}^{r} p_i^{a_i - 1}(p_i - 1)
    = \prod_{i=1}^{r} p_i^{a_i}\left(1 - \frac{1}{p_i}\right).
\]
\end{theoreme}

\begin{proof}
On montre d'abord que $\varphi$ est multiplicative. Soient $m, n$ avec
$\pgcd(m,n) = 1$. Par le théorème chinois des restes\index{théorème chinois des restes},
l'application
\[
  \Z/mn\Z \xrightarrow{\;\sim\;} \Z/m\Z \times \Z/n\Z
\]
est un isomorphisme d'anneaux, qui induit un isomorphisme sur les groupes d'inversibles~:
$(\Z/mn\Z)^\times \simeq (\Z/m\Z)^\times \times (\Z/n\Z)^\times$. Donc
$\varphi(mn) = \varphi(m)\varphi(n)$.

Il reste à calculer $\varphi(p^a)$ pour $p$ premier. Les entiers de $\{1, \ldots, p^a\}$
non premiers à $p^a$ sont exactement les multiples de $p$, au nombre de $p^{a-1}$. Donc
$\varphi(p^a) = p^a - p^{a-1} = p^{a-1}(p-1) = p^a(1 - 1/p)$.

La formule générale s'obtient par multiplicativité.
\end{proof}

\begin{theoreme}[Identité de Gauss]
\label{thm:gauss-phi}
\index{identité!de Gauss ($\varphi$)}
Pour tout $n \geq 1$,
\[
  \sum_{d \mid n} \varphi(d) = n.
\]
\end{theoreme}

\begin{proof}
Considérons les fractions $\frac{1}{n}, \frac{2}{n}, \ldots, \frac{n}{n}$. Chaque
fraction $\frac{k}{n}$, réduite sous forme irréductible, s'écrit $\frac{a}{d}$ où
$d \mid n$ et $\pgcd(a, d) = 1$. Pour un diviseur $d$ fixé de $n$, le nombre de
fractions $\frac{a}{d}$ avec $1 \leq a \leq d$ et $\pgcd(a,d) = 1$ est exactement
$\varphi(d)$. Les $n$ fractions se répartissent donc en classes indexées par les diviseurs
de $n$, ce qui donne
\[
  n = \sum_{d \mid n} \varphi(d). \qedhere
\]
\end{proof}

\begin{exemple}
\label{ex:phi-calculs}
\begin{enumerate}[label=(\alph*)]
  \item $\varphi(12) = 12(1 - 1/2)(1 - 1/3) = 12 \cdot 1/2 \cdot 2/3 = 4$.
    En effet, les entiers premiers à $12$ dans $\{1,\ldots,12\}$ sont $1, 5, 7, 11$.
  \item $\varphi(100) = 100(1-1/2)(1-1/5) = 100 \cdot 1/2 \cdot 4/5 = 40$.
  \item Vérification de $\sum_{d \mid 12} \varphi(d)$~:
    $\varphi(1) + \varphi(2) + \varphi(3) + \varphi(4) + \varphi(6) + \varphi(12)
    = 1 + 1 + 2 + 2 + 2 + 4 = 12$. \checkmark
\end{enumerate}
\end{exemple}

%% ─── 8.3 Fonctions de diviseurs ───────────────────────────────────
\section{Fonctions de diviseurs}
\label{sec:fonctions-diviseurs}

\begin{definition}[Fonctions $\tau$, $\sigma$ et $\sigma_k$]
\label{def:tau-sigma}
\index{tau@$\tau(n)$}
\index{sigma@$\sigma(n)$}
\index{sigma@$\sigma_k(n)$}
Pour $n \geq 1$ et $k \in \N$, on pose~:
\begin{align*}
  \tau(n) &= \sum_{d \mid n} 1 &&\text{(nombre de diviseurs de $n$),} \\
  \sigma_k(n) &= \sum_{d \mid n} d^k &&\text{(somme des puissances $k$-ièmes des diviseurs),} \\
  \sigma(n) &= \sigma_1(n) = \sum_{d \mid n} d &&\text{(somme des diviseurs).}
\end{align*}
On a $\tau(n) = \sigma_0(n)$. On note aussi $d(n) = \tau(n)$.
\end{definition}

\begin{proposition}[Formules explicites]
\label{prop:formules-tau-sigma}
Soit $n = p_1^{a_1} \cdots p_r^{a_r}$ la factorisation de $n$. Alors~:
\begin{enumerate}[label=(\roman*)]
  \item $\tau(n) = \displaystyle\prod_{i=1}^{r} (a_i + 1)$.
  \item $\sigma_k(n) = \displaystyle\prod_{i=1}^{r} \frac{p_i^{k(a_i+1)} - 1}{p_i^k - 1}$
    pour $k \geq 1$.
  \item $\sigma(n) = \displaystyle\prod_{i=1}^{r} \frac{p_i^{a_i+1} - 1}{p_i - 1}$.
\end{enumerate}
\end{proposition}

\begin{proof}
Les trois fonctions sont multiplicatives (ce qui se vérifie directement). Il suffit donc
de les calculer sur les puissances de premiers.
\begin{enumerate}[label=(\roman*)]
  \item $\tau(p^a) = \#\{1, p, p^2, \ldots, p^a\} = a + 1$.
  \item $\sigma_k(p^a) = 1 + p^k + p^{2k} + \cdots + p^{ak}
    = \frac{p^{k(a+1)} - 1}{p^k - 1}$.
  \item Cas particulier de (ii) avec $k = 1$. \qedhere
\end{enumerate}
\end{proof}

\begin{definition}[Nombre parfait]
\label{def:nombre-parfait}
\index{nombre parfait}
Un entier $n \geq 2$ est \emph{parfait} si $\sigma(n) = 2n$, c'est-à-dire si la somme
de ses diviseurs propres (distincts de $n$) est égale à $n$.
\end{definition}

\begin{theoreme}[Euler--Euclide]
\label{thm:euler-euclide-parfaits}
\index{nombre parfait!pair}
Un entier pair $n$ est parfait si et seulement si $n = 2^{p-1}(2^p - 1)$ où $2^p - 1$
est un nombre premier (de Mersenne\index{nombre de Mersenne}).
\end{theoreme}

\begin{proof}
\emph{Suffisance.} Si $q = 2^p - 1$ est premier, alors
$\sigma(2^{p-1}q) = \sigma(2^{p-1})\sigma(q) = (2^p - 1)(q + 1) = q \cdot 2^p
= 2 \cdot 2^{p-1}q$.

\emph{Nécessité.} Soit $n = 2^{a}m$ avec $m$ impair et $a \geq 1$. Alors
$\sigma(n) = \sigma(2^a)\sigma(m) = (2^{a+1}-1)\sigma(m) = 2n = 2^{a+1}m$. Donc
$(2^{a+1}-1) \mid 2^{a+1}m$, et comme $\pgcd(2^{a+1}-1, 2^{a+1}) = 1$, on a
$(2^{a+1}-1) \mid m$. Posons $m = (2^{a+1}-1)k$. Alors
$(2^{a+1}-1)\sigma(m) = 2^{a+1}(2^{a+1}-1)k$, d'où $\sigma(m) = 2^{a+1}k = m + k$.
Si $k > 1$, alors $m$ a au moins les trois diviseurs $1, k, m$, donc
$\sigma(m) \geq 1 + k + m > m + k$, contradiction. Donc $k = 1$, $m = 2^{a+1} - 1$
est premier, et $n = 2^a(2^{a+1}-1)$.
\end{proof}

\begin{exemple}
\label{ex:nombres-parfaits}
Les premiers nombres parfaits pairs sont~: $6 = 2^1(2^2-1)$, $28 = 2^2(2^3-1)$,
$496 = 2^4(2^5-1)$, $8128 = 2^6(2^7-1)$.
L'existence de nombres parfaits impairs reste un problème ouvert célèbre.
\end{exemple}

%% ─── 8.4 Fonction de Möbius ───────────────────────────────────────
\section{La fonction de Möbius $\mu(n)$}
\label{sec:mobius}
\index{fonction de Möbius}
\index{mu@$\mu(n)$}

\begin{definition}[Fonction de Möbius]
\label{def:mobius}
La \emph{fonction de Möbius} est définie par
\[
  \mu(n) = \begin{cases}
    1 & \text{si } n = 1, \\
    (-1)^k & \text{si } n = p_1 p_2 \cdots p_k \text{ avec } p_1, \ldots, p_k
      \text{ premiers distincts,} \\
    0 & \text{si } p^2 \mid n \text{ pour un nombre premier } p.
  \end{cases}
\]
\end{definition}

\begin{proposition}
\label{prop:mu-multiplicative}
La fonction $\mu$ est multiplicative.
\end{proposition}

\begin{proof}
Il est clair que $\mu(1) = 1$. Soient $m, n$ avec $\pgcd(m,n) = 1$. Si $m$ ou $n$ a un
facteur carré, $mn$ en a un aussi, et les deux côtés valent~$0$. Sinon,
$m = p_1 \cdots p_r$ et $n = q_1 \cdots q_s$ sont des produits de premiers distincts.
Comme $\pgcd(m,n)=1$, tous les $p_i, q_j$ sont distincts, et $mn$ est un produit de
$r+s$ premiers distincts. Donc $\mu(mn) = (-1)^{r+s} = (-1)^r(-1)^s = \mu(m)\mu(n)$.
\end{proof}

\begin{theoreme}[Propriété fondamentale de $\mu$]
\label{thm:mu-propriete-fondamentale}
\index{fonction de Möbius!propriété fondamentale}
Pour tout $n \geq 1$,
\[
  \sum_{d \mid n} \mu(d) = [n = 1] =
  \begin{cases} 1 & \text{si } n = 1, \\ 0 & \text{si } n > 1. \end{cases}
\]
\end{theoreme}

\begin{proof}
Pour $n = 1$, la somme vaut $\mu(1) = 1$. Pour $n > 1$, écrivons
$n = p_1^{a_1} \cdots p_r^{a_r}$ avec $r \geq 1$. Un diviseur $d$ de $n$ contribue
$\mu(d) \neq 0$ uniquement si $d$ est sans facteur carré, c'est-à-dire si $d$ est un
produit de premiers distincts choisis parmi $p_1, \ldots, p_r$. Donc
\[
  \sum_{d \mid n} \mu(d) = \sum_{S \subseteq \{p_1, \ldots, p_r\}} (-1)^{\abs{S}}
    = \sum_{k=0}^{r} \binom{r}{k} (-1)^k = (1 - 1)^r = 0.
\]
On a utilisé le binôme de Newton. Ceci achève la preuve.
\end{proof}

\begin{exemple}
\label{ex:mu-calculs}
\begin{enumerate}[label=(\alph*)]
  \item $\mu(1)=1$, $\mu(2)=-1$, $\mu(3)=-1$, $\mu(4)=0$, $\mu(5)=-1$, $\mu(6)=1$.
  \item $\mu(30) = \mu(2 \cdot 3 \cdot 5) = (-1)^3 = -1$.
  \item $\mu(12) = \mu(2^2 \cdot 3) = 0$ car $4 \mid 12$.
  \item $\sum_{d \mid 6} \mu(d) = \mu(1) + \mu(2) + \mu(3) + \mu(6)
    = 1 + (-1) + (-1) + 1 = 0$. \checkmark
\end{enumerate}
\end{exemple}

%% ─── 8.5 Fonction de Liouville ─────────────────────────────────────
\section{La fonction de Liouville $\lambda(n)$}
\label{sec:liouville}
\index{fonction de Liouville}
\index{lambda@$\lambda(n)$ (Liouville)}

\begin{definition}[Fonction de Liouville]
\label{def:liouville}
La \emph{fonction de Liouville} est définie par $\lambda(n) = (-1)^{\Omega(n)}$, où
$\Omega(n)$\index{Omega@$\Omega(n)$} est le nombre total de facteurs premiers de $n$
comptés avec multiplicité.
\end{definition}

\begin{proposition}
\label{prop:liouville-proprietes}
\begin{enumerate}[label=(\roman*)]
  \item La fonction $\lambda$ est complètement multiplicative.
  \item Pour tout $n \geq 1$,
    $\displaystyle\sum_{d \mid n} \lambda(d) = \begin{cases}
      1 & \text{si $n$ est un carré parfait,} \\
      0 & \text{sinon.}
    \end{cases}$
\end{enumerate}
\end{proposition}

\begin{proof}
\begin{enumerate}[label=(\roman*)]
  \item $\Omega(mn) = \Omega(m) + \Omega(n)$, donc $\lambda(mn) = (-1)^{\Omega(m)+\Omega(n)}
    = \lambda(m)\lambda(n)$.
  \item La somme est multiplicative en $n$. Pour $n = p^a$~:
    $\sum_{k=0}^{a} \lambda(p^k) = \sum_{k=0}^{a} (-1)^k = \frac{1-(-1)^{a+1}}{2}$,
    qui vaut $1$ si $a$ est pair et $0$ si $a$ est impair.
    Le résultat suit par multiplicativité. \qedhere
\end{enumerate}
\end{proof}

%% ─── 8.6 Fonction de von Mangoldt ──────────────────────────────────
\section{La fonction de von Mangoldt $\Lambda(n)$}
\label{sec:von-mangoldt}
\index{fonction de von Mangoldt}
\index{Lambda@$\Lambda(n)$ (von Mangoldt)}

\begin{definition}[Fonction de von Mangoldt]
\label{def:von-mangoldt}
La \emph{fonction de von Mangoldt} est définie par
\[
  \Lambda(n) = \begin{cases}
    \ln p & \text{si } n = p^k \text{ pour un premier } p \text{ et un entier } k \geq 1, \\
    0 & \text{sinon.}
  \end{cases}
\]
\end{definition}

\begin{proposition}
\label{prop:von-mangoldt-identite}
Pour tout $n \geq 1$, on a
\[
  \sum_{d \mid n} \Lambda(d) = \ln n.
\]
\end{proposition}

\begin{proof}
Soit $n = p_1^{a_1} \cdots p_r^{a_r}$. Les diviseurs $d$ de $n$ pour lesquels
$\Lambda(d) \neq 0$ sont de la forme $p_i^j$ avec $1 \leq j \leq a_i$. Donc
\[
  \sum_{d \mid n} \Lambda(d) = \sum_{i=1}^{r} \sum_{j=1}^{a_i} \ln p_i
    = \sum_{i=1}^{r} a_i \ln p_i = \ln\!\left(\prod_{i=1}^r p_i^{a_i}\right) = \ln n.
    \qedhere
\]
\end{proof}

\begin{remarque}
\label{rem:von-mangoldt-importance}
La fonction de von Mangoldt joue un rôle central en théorie analytique des nombres. La
\emph{fonction de Tchébychev}\index{fonction de Tchébychev}
$\psi(x) = \sum_{n \leq x} \Lambda(n)$ est liée à la distribution des nombres premiers~:
le théorème des nombres premiers est équivalent à $\psi(x) \sim x$ quand $x \to +\infty$.
\end{remarque}

%% ─── 8.7 Produit de Dirichlet ─────────────────────────────────────
\section{Le produit de Dirichlet}
\label{sec:dirichlet-convolution}
\index{produit de Dirichlet}
\index{convolution de Dirichlet|see{produit de Dirichlet}}

\begin{definition}[Produit de Dirichlet]
\label{def:dirichlet-convolution}
Le \emph{produit de Dirichlet} (ou \emph{convolution de Dirichlet}) de deux fonctions
arithmétiques $f$ et $g$ est la fonction arithmétique $f * g$ définie par
\[
  (f * g)(n) = \sum_{d \mid n} f(d)\, g\!\left(\frac{n}{d}\right)
    = \sum_{\substack{ab = n \\ a, b \geq 1}} f(a)\, g(b).
\]
\end{definition}

\begin{notation}
\label{not:fonctions-speciales}
\index{fonction identité (Dirichlet)}
\index{fonction zeta arithmétique}
On introduit les fonctions arithmétiques suivantes~:
\begin{align*}
  \varepsilon(n) &= [n = 1] = \begin{cases} 1 & \text{si } n = 1, \\ 0 & \text{sinon,}
    \end{cases}
    &&\text{(identité de Dirichlet),} \\
  \mathbf{1}(n) &= 1 \quad \text{pour tout } n,
    &&\text{(fonction constante~$1$),} \\
  \mathrm{id}(n) &= n,
    &&\text{(fonction identité),} \\
  \mathrm{id}_k(n) &= n^k.
\end{align*}
\end{notation}

\begin{theoreme}[Structure d'anneau]
\label{thm:anneau-dirichlet}
\index{produit de Dirichlet!anneau}
L'ensemble des fonctions arithmétiques, muni de l'addition ponctuelle et du produit de
Dirichlet~$*$, forme un anneau commutatif unitaire dont~:
\begin{enumerate}[label=(\roman*)]
  \item le zéro est la fonction nulle,
  \item l'unité est $\varepsilon$,
  \item le produit $*$ est associatif et commutatif.
\end{enumerate}
\end{theoreme}

\begin{proof}
La commutativité est immédiate par le changement $d \mapsto n/d$. Pour l'associativité~:
\begin{align*}
  ((f*g)*h)(n)
    &= \sum_{d \mid n} (f*g)(d)\, h\!\left(\frac{n}{d}\right)
     = \sum_{d \mid n} \sum_{e \mid d} f(e)\, g\!\left(\frac{d}{e}\right)
       h\!\left(\frac{n}{d}\right) \\
    &= \sum_{\substack{abc = n}} f(a)\, g(b)\, h(c)
     = (f*(g*h))(n).
\end{align*}
La propriété $f * \varepsilon = f$ est claire~: $(f * \varepsilon)(n) = \sum_{d \mid n}
f(d)\,\varepsilon(n/d) = f(n)$ car $\varepsilon(n/d) = 0$ sauf pour $d = n$.
\end{proof}

Les résultats des sections précédentes se réexpriment élégamment~:

\begin{proposition}[Identités de convolution]
\label{prop:identites-convolution}
\index{produit de Dirichlet!identités}
\begin{enumerate}[label=(\roman*)]
  \item $\mathbf{1} * \mathbf{1} = \tau$ \quad
    {\small (car $\sum_{d \mid n} 1 = \tau(n)$).}
  \item $\mathrm{id} * \mathbf{1} = \sigma$ \quad
    {\small (car $\sum_{d \mid n} d = \sigma(n)$).}
  \item $\varphi * \mathbf{1} = \mathrm{id}$ \quad
    {\small (identité de Gauss, théorème~\ref{thm:gauss-phi}).}
  \item $\mu * \mathbf{1} = \varepsilon$ \quad
    {\small (propriété fondamentale de $\mu$, théorème~\ref{thm:mu-propriete-fondamentale}).}
  \item $\Lambda * \mathbf{1} = \ln$ \quad
    {\small (proposition~\ref{prop:von-mangoldt-identite}).}
\end{enumerate}
\end{proposition}

\begin{theoreme}[Inversibilité]
\label{thm:inversibilite-dirichlet}
\index{produit de Dirichlet!inversibilité}
Une fonction arithmétique $f$ admet un inverse pour le produit de Dirichlet si et
seulement si $f(1) \neq 0$. L'inverse $f^{-1}$ est défini récursivement par
\[
  f^{-1}(1) = \frac{1}{f(1)}, \qquad
  f^{-1}(n) = -\frac{1}{f(1)} \sum_{\substack{d \mid n \\ d < n}}
    f\!\left(\frac{n}{d}\right) f^{-1}(d) \quad \text{pour } n > 1.
\]
\end{theoreme}

\begin{proof}
La condition $f * f^{-1} = \varepsilon$ donne, pour $n = 1$~:
$f(1)f^{-1}(1) = 1$, d'où la nécessité de $f(1) \neq 0$ et la valeur de $f^{-1}(1)$.
Pour $n > 1$~: $(f*f^{-1})(n) = 0$ donne $\sum_{d \mid n} f(n/d)\,f^{-1}(d) = 0$, soit
$f(1)f^{-1}(n) + \sum_{\substack{d \mid n \\ d < n}} f(n/d)\,f^{-1}(d) = 0$,
d'où la formule récursive.
\end{proof}

\begin{corollaire}
\label{cor:mu-inverse-1}
La fonction de Möbius $\mu$ est l'inverse de Dirichlet de $\mathbf{1}$~:
$\mu * \mathbf{1} = \varepsilon$, c'est-à-dire $\mu = \mathbf{1}^{-1}$.
\end{corollaire}

\begin{proposition}[Produit de fonctions multiplicatives]
\label{prop:produit-multiplicatives}
Si $f$ et $g$ sont multiplicatives, alors $f * g$ est multiplicative.
\end{proposition}

\begin{proof}
Soient $m, n$ avec $\pgcd(m,n) = 1$. Tout diviseur $d$ de $mn$ s'écrit de manière unique
$d = d_1 d_2$ avec $d_1 \mid m$ et $d_2 \mid n$ (et $\pgcd(d_1, d_2) = 1$). Donc
\begin{align*}
  (f*g)(mn) &= \sum_{d \mid mn} f(d)\, g\!\left(\frac{mn}{d}\right)
    = \sum_{d_1 \mid m} \sum_{d_2 \mid n} f(d_1 d_2)\, g\!\left(\frac{m}{d_1}
      \cdot \frac{n}{d_2}\right) \\
    &= \sum_{d_1 \mid m} \sum_{d_2 \mid n} f(d_1)f(d_2)\, g\!\left(\frac{m}{d_1}\right)
      g\!\left(\frac{n}{d_2}\right) \\
    &= \left(\sum_{d_1 \mid m} f(d_1)\, g\!\left(\frac{m}{d_1}\right)\right)
       \left(\sum_{d_2 \mid n} f(d_2)\, g\!\left(\frac{n}{d_2}\right)\right) \\
    &= (f*g)(m)\,(f*g)(n). \qedhere
\end{align*}
\end{proof}

%% ─── 8.8 Visualisation du treillis des diviseurs ──────────────────
\subsection{Visualisation : treillis des diviseurs}

\begin{figure}[ht]
\centering
\begin{tikzpicture}[
  every node/.style={circle, draw, inner sep=2pt, minimum size=20pt,
    font=\small},
  level/.style={sibling distance=30mm/#1},
  edge from parent/.style={draw, thick, -}
]
  % Treillis des diviseurs de 12
  \node (12) at (0,4) {$12$};
  \node (6) at (-2,2.7) {$6$};
  \node (4) at (0,2.7) {$4$};
  \node (3) at (-2,1.3) {$3$};
  \node (2) at (0,1.3) {$2$};
  \node (1) at (-1,0) {$1$};

  \draw[thick] (1) -- (2) -- (4) -- (12);
  \draw[thick] (1) -- (3) -- (6) -- (12);
  \draw[thick] (2) -- (6);
  \draw[thick] (3) -- (12);

  % Annotations
  \node[draw=none, right=5mm] at (12.east) {$\mu(12)=0$};
  \node[draw=none, right=5mm] at (6.east) {$\mu(6)=1$};
  \node[draw=none, left=5mm] at (4.west) {$\mu(4)=0$};
  \node[draw=none, right=5mm] at (3.east) {$\mu(3)=-1$};
  \node[draw=none, left=5mm] at (2.west) {$\mu(2)=-1$};
  \node[draw=none, left=5mm] at (1.west) {$\mu(1)=1$};
\end{tikzpicture}
\caption{Treillis des diviseurs de $12$ avec les valeurs de $\mu$.
On vérifie $\sum_{d \mid 12} \mu(d) = 1 + (-1) + (-1) + 0 + 1 + 0 = 0$.}
\label{fig:treillis-diviseurs-12}
\end{figure}

%% ─── 8.9 Visualisation du produit de Dirichlet ────────────────────
\begin{figure}[ht]
\centering
\begin{tikzpicture}[>=Stealth, scale=0.85]
  % Titre
  \node[font=\bfseries] at (4,5.5) {Produit de Dirichlet $(f*g)(12)$};

  % Paires (d, 12/d)
  \foreach \d/\q/\y in {1/12/4.5, 2/6/3.5, 3/4/2.5, 4/3/1.5, 6/2/0.5, 12/1/-0.5} {
    % Boîtes
    \node[draw, fill=blue!10, minimum width=12mm, minimum height=7mm] (fd\d) at (1.5,\y)
      {$f(\d)$};
    \node[font=\Large] at (3.2,\y) {$\times$};
    \node[draw, fill=red!10, minimum width=12mm, minimum height=7mm] (gq\d) at (4.8,\y)
      {$g(\q)$};
    \node[font=\Large] at (6.3,\y) {$=$};
    \node[draw, fill=green!10, minimum width=20mm, minimum height=7mm] at (8,\y)
      {$f(\d)\,g(\q)$};
  }

  % Accolade et somme
  \draw[decorate, decoration={brace, amplitude=5pt}] (9.5,-0.9) -- (9.5,4.9);
  \node[font=\large] at (11.2,2) {$\displaystyle\sum$};
\end{tikzpicture}
\caption{Schéma du calcul de $(f*g)(12) = \sum_{d \mid 12} f(d)\,g(12/d)$.}
\label{fig:dirichlet-convolution}
\end{figure}

%% ─── 8.10 Inversion de Möbius ─────────────────────────────────────
\section{La formule d'inversion de Möbius}
\label{sec:inversion-mobius}
\index{inversion de Möbius}

\begin{theoreme}[Inversion de Möbius]
\label{thm:inversion-mobius}
Soient $f$ et $g$ deux fonctions arithmétiques. On a l'équivalence~:
\[
  g(n) = \sum_{d \mid n} f(d) \quad \text{pour tout } n \geq 1
  \qquad \Longleftrightarrow \qquad
  f(n) = \sum_{d \mid n} \mu(d)\, g\!\left(\frac{n}{d}\right)
    \quad \text{pour tout } n \geq 1.
\]
Autrement dit~: $g = f * \mathbf{1} \;\Leftrightarrow\; f = g * \mu = \mu * g$.
\end{theoreme}

\begin{proof}
\emph{Sens direct ($\Rightarrow$).} Supposons $g = f * \mathbf{1}$. Alors
\begin{align*}
  (\mu * g)(n) &= \sum_{d \mid n} \mu(d)\, g\!\left(\frac{n}{d}\right)
    = \sum_{d \mid n} \mu(d) \sum_{e \mid (n/d)} f(e) \\
    &= \sum_{e \mid n} f(e) \sum_{\substack{d \mid n \\ e \mid (n/d)}} \mu(d)
     = \sum_{e \mid n} f(e) \sum_{d \mid (n/e)} \mu(d).
\end{align*}
Par le théorème~\ref{thm:mu-propriete-fondamentale}, $\sum_{d \mid (n/e)} \mu(d) =
[n/e = 1] = [e = n]$. Donc $(\mu * g)(n) = f(n)$.

\emph{Sens réciproque ($\Leftarrow$).} Supposons $f = \mu * g$. Alors
\begin{align*}
  (f * \mathbf{1})(n) &= \sum_{d \mid n} f(d) = \sum_{d \mid n} (\mu * g)(d)
    = \sum_{d \mid n} \sum_{e \mid d} \mu(e)\, g\!\left(\frac{d}{e}\right) \\
    &= \sum_{e \mid n} \mu(e) \sum_{\substack{d \mid n \\ e \mid d}}
      g\!\left(\frac{d}{e}\right)
    = \sum_{e \mid n} \mu(e) \sum_{k \mid (n/e)} g(k).
\end{align*}
Posant $h(m) = \sum_{k \mid m} g(k) = (g * \mathbf{1})(m)$, cela vaut
$\sum_{e \mid n} \mu(e)\, h(n/e) = (\mu * h)(n) = (\mu * g * \mathbf{1})(n)$.
Or $\mu * \mathbf{1} = \varepsilon$, donc par associativité,
$\mu * g * \mathbf{1} = g * (\mu * \mathbf{1}) = g * \varepsilon = g$.
Donc $(f * \mathbf{1})(n) = g(n)$.
\end{proof}

\subsection{Applications de l'inversion de Möbius}

\begin{corollaire}[Formule de $\varphi$ par inversion]
\label{cor:phi-par-inversion}
On a $\varphi = \mu * \mathrm{id}$, c'est-à-dire
\[
  \varphi(n) = \sum_{d \mid n} \mu(d) \cdot \frac{n}{d}
    = n \sum_{d \mid n} \frac{\mu(d)}{d}.
\]
\end{corollaire}

\begin{proof}
L'identité de Gauss (théorème~\ref{thm:gauss-phi}) s'écrit
$\mathrm{id} = \varphi * \mathbf{1}$. Par inversion de Möbius,
$\varphi = \mu * \mathrm{id}$.
\end{proof}

\begin{exemple}
\label{ex:phi-par-inversion}
Calculons $\varphi(12)$ par la formule d'inversion~:
\begin{align*}
  \varphi(12) &= \sum_{d \mid 12} \mu(d) \cdot \frac{12}{d} \\
    &= \mu(1) \cdot 12 + \mu(2) \cdot 6 + \mu(3) \cdot 4 + \mu(4) \cdot 3
      + \mu(6) \cdot 2 + \mu(12) \cdot 1 \\
    &= 1 \cdot 12 + (-1) \cdot 6 + (-1) \cdot 4 + 0 \cdot 3 + 1 \cdot 2 + 0 \cdot 1
     = 12 - 6 - 4 + 2 = 4.
\end{align*}
\end{exemple}

\begin{corollaire}[Formule de $\Lambda$ par inversion]
\label{cor:Lambda-par-inversion}
On a $\Lambda = \mu * (-\ln)$, c'est-à-dire, en écrivant $L(n) = \ln n$~:
$L = \Lambda * \mathbf{1}$ implique $\Lambda = \mu * L$, soit
\[
  \Lambda(n) = -\sum_{d \mid n} \mu(d) \ln d = \sum_{d \mid n} \mu(d) \ln\frac{n}{d}.
\]
\end{corollaire}

\subsection{Forme sommatoire de l'inversion de Möbius}

\begin{theoreme}[Inversion de Möbius -- forme sommatoire]
\label{thm:inversion-sommatoire}
\index{inversion de Möbius!forme sommatoire}
Soient $f, g : \N^* \to \C$. On a l'équivalence~:
\[
  g(x) = \sum_{n \leq x} f\!\left(\floor*{\frac{x}{n}}\right)
  \qquad \Longleftrightarrow \qquad
  f(x) = \sum_{n \leq x} \mu(n)\, g\!\left(\floor*{\frac{x}{n}}\right).
\]
Plus généralement, si $G(x) = \sum_{n \leq x} f(n)$ et $H(x) = \sum_{n \leq x} g(n)$
avec $g = f * \mathbf{1}$, alors
\[
  F(x) = \sum_{n \leq x} f(n) = \sum_{n \leq x} \mu(n)\, H\!\left(\frac{x}{n}\right).
\]
\end{theoreme}

%% ─── 8.11 Séries de Dirichlet ─────────────────────────────────────
\section{Introduction aux séries de Dirichlet}
\label{sec:series-dirichlet}
\index{série de Dirichlet}

\begin{definition}[Série de Dirichlet]
\label{def:serie-dirichlet}
À toute fonction arithmétique $f$, on associe la \emph{série de Dirichlet}
\[
  F(s) = \sum_{n=1}^{\infty} \frac{f(n)}{n^s}, \quad s \in \C.
\]
\end{definition}

\begin{theoreme}[Produit d'Euler]
\label{thm:produit-euler}
\index{produit d'Euler}
Si $f$ est multiplicative et la série $F(s)$ converge absolument, alors
\[
  F(s) = \prod_{p \text{ premier}} \left(\sum_{k=0}^{\infty}
    \frac{f(p^k)}{p^{ks}}\right).
\]
Si $f$ est complètement multiplicative, cela se simplifie en
\[
  F(s) = \prod_{p \text{ premier}} \frac{1}{1 - f(p)\, p^{-s}}.
\]
\end{theoreme}

\begin{exemple}[Séries de Dirichlet classiques]
\label{ex:series-dirichlet-classiques}
\begin{enumerate}[label=(\alph*)]
  \item La \emph{fonction zêta de Riemann}\index{fonction zêta de Riemann}~:
    $\displaystyle\zeta(s) = \sum_{n=1}^{\infty} \frac{1}{n^s}
    = \prod_p \frac{1}{1 - p^{-s}}$ pour $\Re(s) > 1$.
  \item $\displaystyle\sum_{n=1}^{\infty} \frac{\mu(n)}{n^s} = \frac{1}{\zeta(s)}$
    pour $\Re(s) > 1$.
  \item $\displaystyle\sum_{n=1}^{\infty} \frac{\varphi(n)}{n^s}
    = \frac{\zeta(s-1)}{\zeta(s)}$ pour $\Re(s) > 2$.
  \item $\displaystyle\sum_{n=1}^{\infty} \frac{\tau(n)}{n^s} = \zeta(s)^2$
    pour $\Re(s) > 1$.
  \item $\displaystyle\sum_{n=1}^{\infty} \frac{\sigma(n)}{n^s}
    = \zeta(s)\,\zeta(s-1)$ pour $\Re(s) > 2$.
\end{enumerate}
\end{exemple}

\begin{remarque}
\label{rem:convolution-series}
Le produit de Dirichlet $f * g$ correspond au \emph{produit} des séries de Dirichlet~: si
$F(s) = \sum f(n)/n^s$ et $G(s) = \sum g(n)/n^s$, alors
$F(s) \cdot G(s) = \sum (f*g)(n)/n^s$, dans la zone de convergence absolue commune. Cela
justifie par exemple $\mu * \mathbf{1} = \varepsilon \Leftrightarrow
1/\zeta(s) \cdot \zeta(s) = 1$.
\end{remarque}

%% ─── 8.12 Tableau récapitulatif ───────────────────────────────────
\section{Tableau récapitulatif}
\label{sec:tableau-fonctions}

\begin{center}
\renewcommand{\arraystretch}{1.4}
\begin{tabular}{llll}
\toprule
\textbf{Fonction} & \textbf{Définition} & \textbf{Mult.} & \textbf{Série de Dirichlet} \\
\midrule
$\varepsilon(n) = [n=1]$ & identité de $*$ & oui & $1$ \\
$\mathbf{1}(n) = 1$ & constante & comp.\ mult.\ & $\zeta(s)$ \\
$\mathrm{id}(n) = n$ & identité & comp.\ mult.\ & $\zeta(s-1)$ \\
$\mu(n)$ & Möbius & mult. & $1/\zeta(s)$ \\
$\varphi(n)$ & Euler & mult. & $\zeta(s-1)/\zeta(s)$ \\
$\tau(n) = \sigma_0(n)$ & nb.\ de diviseurs & mult. & $\zeta(s)^2$ \\
$\sigma(n) = \sigma_1(n)$ & somme des diviseurs & mult. & $\zeta(s)\zeta(s-1)$ \\
$\lambda(n)$ & Liouville & comp.\ mult.\ & $\zeta(2s)/\zeta(s)$ \\
$\Lambda(n)$ & von Mangoldt & non & $-\zeta'(s)/\zeta(s)$ \\
\bottomrule
\end{tabular}
\end{center}

%% ─── 8.13 Exemples de calcul ──────────────────────────────────────
\section{Exemples de calcul}
\label{sec:exemples-calcul-ch8}

\begin{exemple}[Calcul complet pour $n = 60$]
\label{ex:calcul-60}
Soit $n = 60 = 2^2 \cdot 3 \cdot 5$.
\begin{enumerate}[label=(\alph*)]
  \item $\varphi(60) = 60(1 - 1/2)(1-1/3)(1-1/5) = 60 \cdot 1/2 \cdot 2/3 \cdot 4/5 = 16$.
  \item $\tau(60) = (2+1)(1+1)(1+1) = 12$.
  \item $\sigma(60) = \frac{2^3-1}{2-1} \cdot \frac{3^2-1}{3-1} \cdot \frac{5^2-1}{5-1}
    = 7 \cdot 4 \cdot 6 = 168$.
  \item $\mu(60) = 0$ car $4 \mid 60$.
  \item $\lambda(60) = (-1)^{2+1+1} = (-1)^4 = 1$.
  \item $\Lambda(60) = 0$ car $60$ n'est pas une puissance de premier.
\end{enumerate}
Vérification~: $\sum_{d \mid 60} \varphi(d) = \varphi(1) + \varphi(2) + \varphi(3) +
\varphi(4) + \varphi(5) + \varphi(6) + \varphi(10) + \varphi(12) + \varphi(15) +
\varphi(20) + \varphi(30) + \varphi(60) = 1 + 1 + 2 + 2 + 4 + 2 + 4 + 4 + 8 + 8 +
8 + 16 = 60$. \checkmark
\end{exemple}

\begin{exemple}[Identités via convolution]
\label{ex:identites-convolution}
\begin{enumerate}[label=(\alph*)]
  \item De $\varphi * \mathbf{1} = \mathrm{id}$ et $\mu * \mathbf{1} = \varepsilon$, on
    déduit~:
    $\varphi = \mathrm{id} * \mu$, puis $\sigma = \mathrm{id} * \mathbf{1}$, d'où
    $\sigma * \mu = \mathrm{id} * \mathbf{1} * \mu = \mathrm{id} * \varepsilon
    = \mathrm{id}$.
    Conclusion~: $\sum_{d \mid n} \mu(d)\,\sigma(n/d) = n$.
  \item De $\tau = \mathbf{1} * \mathbf{1}$, on obtient
    $\tau * \mu = \mathbf{1} * \mathbf{1} * \mu = \mathbf{1}$, soit
    $\sum_{d \mid n} \mu(d)\,\tau(n/d) = 1$ pour tout $n \geq 1$.
\end{enumerate}
\end{exemple}

\begin{exemple}[Application~: polynômes cyclotomiques]
\label{ex:cyclotomiques}
\index{polynôme cyclotomique}
Le $n$-ième polynôme cyclotomique s'exprime par inversion de Möbius~:
\[
  \Phi_n(x) = \prod_{d \mid n} (x^d - 1)^{\mu(n/d)}.
\]
En effet, $x^n - 1 = \prod_{d \mid n} \Phi_d(x)$, et l'inversion de Möbius
multiplicative donne la formule ci-dessus.
\end{exemple}

%% ─── 8.14 Exercices ───────────────────────────────────────────────
\section{Exercices}
\label{sec:exos-ch8}

\begin{exercice}[$\star$]
\label{exo:calcul-fonctions}
Calculer $\varphi(n)$, $\tau(n)$, $\sigma(n)$, $\mu(n)$ et $\lambda(n)$ pour
$n \in \{36, 72, 100, 105, 210, 360\}$.
\end{exercice}

\begin{exercice}[$\star$]
\label{exo:verification-gauss}
Vérifier l'identité $\sum_{d \mid n} \varphi(d) = n$ pour $n = 24$ et $n = 30$.
\end{exercice}

\begin{exercice}[$\star$]
\label{exo:convolution-calcul}
Calculer le produit de Dirichlet $(\mu * \mathrm{id})(n)$ pour $n = 1, 2, \ldots, 10$
et vérifier que l'on obtient $\varphi(n)$.
\end{exercice}

\begin{exercice}[$\star\star$]
\label{exo:phi-pair}
Montrer que $\varphi(n)$ est pair pour tout $n \geq 3$.
\end{exercice}

\begin{exercice}[$\star\star$]
\label{exo:sigma-multiplicative}
Démontrer directement (sans utiliser les séries de Dirichlet) que $\sigma$ est
multiplicative.
\end{exercice}

\begin{exercice}[$\star\star$]
\label{exo:identite-piltz}
Montrer que $\sum_{d \mid n} \mu(d)^2 = 2^{\omega(n)}$ où $\omega(n)$ est le nombre de
facteurs premiers distincts de $n$.
\end{exercice}

\begin{exercice}[$\star\star$]
\label{exo:mu-carre-formule}
Montrer que $\sum_{d \mid n} \abs{\mu(d)} \cdot \varphi(n/d) = n \prod_{p \mid n}
\frac{2}{p+1} \cdot \frac{p+1}{1} \cdot \frac{1}{1}$.
Plus précisément, montrer que $\mu^2 * \varphi = \mathrm{id} \cdot (\text{une fonction
multiplicative})$, et identifier cette fonction.
\textit{Indication~: évaluer les deux côtés sur les puissances de premiers.}
\end{exercice}

\begin{exercice}[$\star\star$]
\label{exo:inversion-application}
En utilisant l'identité $\ln n = \sum_{d \mid n} \Lambda(d)$, exprimer la fonction de
von Mangoldt $\Lambda(n)$ comme somme impliquant $\mu$ et $\ln$.
\end{exercice}

\begin{exercice}[$\star\star$]
\label{exo:nombre-sans-carre}
Soit $Q(x) = \#\{n \leq x : n \text{ est sans facteur carré}\}$. Montrer que
$Q(x) = \sum_{d \leq \sqrt{x}} \mu(d) \floor{x/d^2}$.
\textit{Indication~: utiliser $\mu^2(n) = \sum_{d^2 \mid n} \mu(d)$.}
\end{exercice}

\begin{exercice}[$\star\star\star$]
\label{exo:mertens}
Soit $M(x) = \sum_{n \leq x} \mu(n)$ la fonction de Mertens\index{fonction de Mertens}.
\begin{enumerate}[label=(\alph*)]
  \item Montrer que $\sum_{n \leq x} M(\floor{x/n}) = 1$ pour tout $x \geq 1$.
  \item En admettant le théorème des nombres premiers sous la forme $\psi(x) \sim x$,
    montrer que $M(x) = o(x)$.
\end{enumerate}
\end{exercice}

\begin{exercice}[$\star\star\star$]
\label{exo:euler-product-zeta}
\begin{enumerate}[label=(\alph*)]
  \item Montrer que pour $\Re(s) > 1$, le produit d'Euler
    $\prod_p (1-p^{-s})^{-1}$ converge et vaut $\zeta(s)$.
  \item En déduire que $\sum_{n=1}^{\infty} \mu(n)/n^s = 1/\zeta(s)$ pour $\Re(s) > 1$.
  \item Montrer que $\zeta(s)$ a un pôle simple en $s = 1$ avec résidu~$1$, et interpréter
    en termes de densité des nombres premiers.
\end{enumerate}
\end{exercice}

\begin{exercice}[$\star\star\star$]
\label{exo:ramanujan}
La \emph{somme de Ramanujan}\index{somme de Ramanujan} est
$c_q(n) = \sum_{\substack{a=1 \\ \pgcd(a,q)=1}}^{q} e^{2\pi i an/q}$.
Montrer que $c_q(n) = \sum_{d \mid \pgcd(q,n)} d\,\mu(q/d)$.
\end{exercice}

%% ─── 8.15 Résumé du chapitre ──────────────────────────────────────
\section*{Résumé du chapitre}
\addcontentsline{toc}{section}{Résumé du chapitre}

\begin{enumerate}[label=\textbf{\arabic*.}]
  \item Une \textbf{fonction arithmétique} est une application $f : \N^* \to \C$. Elle
    est \textbf{multiplicative} si $f(mn) = f(m)f(n)$ pour $\pgcd(m,n) = 1$.
  \item L'\textbf{indicatrice d'Euler} satisfait $\varphi(n) = n\prod_{p \mid n}(1-1/p)$
    et $\sum_{d \mid n} \varphi(d) = n$.
  \item La \textbf{fonction de Möbius} satisfait
    $\sum_{d \mid n} \mu(d) = [n=1]$ et est l'inverse de Dirichlet de $\mathbf{1}$.
  \item Le \textbf{produit de Dirichlet} $(f*g)(n) = \sum_{d \mid n} f(d)\,g(n/d)$ munit
    les fonctions arithmétiques d'une structure d'anneau commutatif unitaire, d'identité
    $\varepsilon$.
  \item La \textbf{formule d'inversion de Möbius} stipule que $g = f * \mathbf{1}
    \Leftrightarrow f = g * \mu$.
  \item Les \textbf{séries de Dirichlet} $F(s) = \sum f(n)/n^s$ transforment la
    convolution de Dirichlet en produit ordinaire, et admettent un \textbf{produit
    d'Euler} lorsque $f$ est multiplicative.
  \item La fonction $\zeta$ de Riemann, les fonctions $\tau$, $\sigma$, $\varphi$, $\mu$,
    $\lambda$, $\Lambda$ forment un réseau d'identités reliées par le produit de Dirichlet.
\end{enumerate}

% === FIN DU CHUNK ch7-8 ===
% === DÉBUT DU CHUNK ch9-10+end ===

%%%%%%%%%%%%%%%%%%%%%%%%%%%%%%%%%%%%%%%%%%%%%%%%%%%%%%%%%%%%%%%%%%%%
%%  CHAPITRE 9 — Introduction aux Nombres p-adiques
%%%%%%%%%%%%%%%%%%%%%%%%%%%%%%%%%%%%%%%%%%%%%%%%%%%%%%%%%%%%%%%%%%%%
\chapter{Introduction aux Nombres \texorpdfstring{$p$}{p}-adiques}
\label{chap:p-adiques}
\index{nombres p-adiques@nombres $p$-adiques|(}

\begin{center}
\itshape
<<~Les nombres $p$-adiques ne sont pas un simple artifice technique ;\\
ils révèlent une géométrie arithmétique invisible dans $\Q$.~>>\\[4pt]
--- Kurt Hensel (1861--1941)
\end{center}

\bigskip

\section{Introduction historique}
\label{sec:p-adic-intro}

En 1897, Kurt Hensel\index{Hensel, Kurt} introduisit les nombres $p$-adiques
dans le but d'appliquer les méthodes de la théorie des séries de puissances
à la théorie des nombres. L'idée fondamentale est la suivante : de même que les
nombres réels $\R$ s'obtiennent en complétant $\Q$ par rapport à la valeur absolue
usuelle $|\cdot|$, on peut compléter $\Q$ par rapport à d'autres distances, indexées
par les nombres premiers $p$. Chaque complétion produit un corps
$\Q_p$\index{Qp@$\Q_p$} radicalement différent de $\R$, mais tout aussi
fondamental pour la compréhension de l'arithmétique.

Cette construction offre un point de vue \emph{local} sur les propriétés des
nombres : résoudre une équation dans $\Q_p$ revient à analyser sa solubilité
<<~modulo les puissances de $p$~>>. Le \emph{principe local-global}\index{principe
local-global} affirme que, pour certaines classes d'équations, la solubilité dans
$\Q$ équivaut à la solubilité simultanée dans $\R$ et dans tous les $\Q_p$.

\section{Valuation \texorpdfstring{$p$}{p}-adique}
\label{sec:valuation-p-adique}
\index{valuation p-adique@valuation $p$-adique|(}

\begin{definition}[Valuation $p$-adique sur $\Z$]
\label{def:valuation-p-adique-Z}
Soit $p$ un nombre premier. Pour tout entier $n \in \Z \setminus \{0\}$, la
\emph{valuation $p$-adique}\index{valuation p-adique@valuation $p$-adique} de $n$,
notée $v_p(n)$, est le plus grand entier $k \geq 0$ tel que $p^k \mid n$.
Autrement dit,
\[
  v_p(n) = \max\bigl\{ k \in \N : p^k \mid n \bigr\}.
\]
Par convention, on pose $v_p(0) = +\infty$.
\end{definition}

\begin{exemple}[Calculs de valuations]
\label{ex:calculs-valuations}
\begin{enumerate}[label=(\roman*)]
  \item $v_2(12) = v_2(2^2 \cdot 3) = 2$, \quad $v_3(12) = 1$, \quad $v_5(12) = 0$.
  \item $v_3(81) = v_3(3^4) = 4$.
  \item $v_7(2401) = v_7(7^4) = 4$.
  \item $v_2(1000) = v_2(2^3 \cdot 125) = 3$, \quad $v_5(1000) = 3$.
\end{enumerate}
\end{exemple}

\begin{definition}[Valuation $p$-adique sur $\Q$]
\label{def:valuation-p-adique-Q}
On étend $v_p$ à $\Q^*$ en posant, pour $r = a/b$ avec $a,b \in \Z$, $b \neq 0$ :
\[
  v_p\!\left(\frac{a}{b}\right) = v_p(a) - v_p(b).
\]
Cette définition ne dépend pas du représentant choisi pour la fraction.
\end{definition}

\begin{proposition}[Propriétés de la valuation $p$-adique]
\label{prop:proprietes-valuation}
Pour tout $x, y \in \Q^*$, on a :
\begin{enumerate}[label=(\roman*)]
  \item $v_p(xy) = v_p(x) + v_p(y)$ \quad (morphisme multiplicatif) ;
  \item $v_p(x + y) \geq \min\bigl(v_p(x), v_p(y)\bigr)$, avec égalité si
    $v_p(x) \neq v_p(y)$.
\end{enumerate}
\end{proposition}

\begin{proof}
(i) Soient $x = p^a \cdot u$ et $y = p^b \cdot w$ avec $\pgcd(u, p) = \pgcd(w, p) = 1$.
Alors $xy = p^{a+b} \cdot uw$, et $\pgcd(uw, p) = 1$, d'où $v_p(xy) = a + b = v_p(x) + v_p(y)$.

(ii) Écrivons $x = p^a \cdot u/v$ et $y = p^b \cdot r/s$ sous forme réduite,
avec $a = v_p(x)$, $b = v_p(y)$. Supposons $a \leq b$. Alors
\[
  x + y = p^a \left( \frac{u}{v} + p^{b-a} \cdot \frac{r}{s} \right),
\]
d'où $v_p(x+y) \geq a = \min(v_p(x), v_p(y))$. Si $a < b$, le terme entre
parenthèses a valuation $0$, d'où l'égalité.
\end{proof}

\begin{remarque}
\label{rem:decomposition-fondamentale}
La décomposition en facteurs premiers de tout $n \in \Z \setminus \{0\}$ se lit
directement sur les valuations :
\[
  n = \pm \prod_{p \text{ premier}} p^{v_p(n)},
\]
où le produit est fini car $v_p(n) = 0$ pour presque tout $p$.
\end{remarque}

\index{valuation p-adique@valuation $p$-adique|)}

\section{Valeur absolue \texorpdfstring{$p$}{p}-adique}
\label{sec:valeur-absolue-p-adique}
\index{valeur absolue p-adique@valeur absolue $p$-adique|(}

\begin{definition}[Valeur absolue $p$-adique]
\label{def:valeur-absolue-p-adique}
Pour $x \in \Q$, on définit la \emph{valeur absolue $p$-adique} par
\[
  \abs{x}_p =
  \begin{cases}
    p^{-v_p(x)} & \text{si } x \neq 0, \\
    0 & \text{si } x = 0.
  \end{cases}
\]
\end{definition}

\begin{exemple}
\label{ex:valeur-absolue-p-adique}
\begin{enumerate}[label=(\roman*)]
  \item $\abs{12}_3 = 3^{-v_3(12)} = 3^{-1} = 1/3$.
  \item $\abs{75}_5 = 5^{-v_5(75)} = 5^{-2} = 1/25$.
  \item $\abs{1/7}_7 = 7^{-v_7(1/7)} = 7^{-(-1)} = 7$.
  \item Pour $\pgcd(n, p) = 1$ : $\abs{n}_p = 1$.
\end{enumerate}
\end{exemple}

\begin{theoreme}[Propriétés de $\abs{\cdot}_p$]
\label{thm:proprietes-valeur-absolue}
L'application $\abs{\cdot}_p : \Q \to \R_{\geq 0}$ vérifie :
\begin{enumerate}[label=(\roman*)]
  \item \textbf{Positivité :} $\abs{x}_p = 0 \iff x = 0$ ;
  \item \textbf{Multiplicativité :} $\abs{xy}_p = \abs{x}_p \cdot \abs{y}_p$ ;
  \item \textbf{Inégalité ultramétrique\index{inégalité ultramétrique} :}
    $\abs{x + y}_p \leq \max\bigl(\abs{x}_p, \abs{y}_p\bigr)$.
\end{enumerate}
\end{theoreme}

\begin{proof}
Les points (i) et (ii) découlent directement des propriétés de $v_p$.

Pour (iii), si $x = 0$ ou $y = 0$, c'est immédiat. Sinon :
\begin{align*}
  \abs{x + y}_p &= p^{-v_p(x+y)} \\
  &\leq p^{-\min(v_p(x),\, v_p(y))} \quad \text{(par la Proposition~\ref{prop:proprietes-valuation})} \\
  &= \max\bigl(p^{-v_p(x)},\, p^{-v_p(y)}\bigr) \\
  &= \max\bigl(\abs{x}_p,\, \abs{y}_p\bigr). \qedhere
\end{align*}
\end{proof}

\begin{remarque}[Conséquence de l'ultramétricité]
\label{rem:ultrametrique-consequences}
L'inégalité ultramétrique est \emph{plus forte} que l'inégalité triangulaire usuelle :
\[
  \abs{x+y}_p \leq \max(\abs{x}_p, \abs{y}_p) \leq \abs{x}_p + \abs{y}_p.
\]
Cela entraîne des propriétés géométriques surprenantes : dans un espace
ultramétrique, tout triangle est isocèle, et toute boule ouverte est aussi fermée.
\end{remarque}

\begin{proposition}[Formule du produit]
\label{prop:formule-produit}
\index{formule du produit}
Pour tout $x \in \Q^*$,
\[
  \abs{x}_\infty \cdot \prod_{p \text{ premier}} \abs{x}_p = 1,
\]
où $\abs{\cdot}_\infty$ désigne la valeur absolue usuelle sur $\Q$.
\end{proposition}

\begin{proof}
Il suffit de vérifier pour $x = \pm p_1^{a_1} \cdots p_k^{a_k}$. On a
$\abs{x}_\infty = p_1^{a_1} \cdots p_k^{a_k}$ et $\abs{x}_{p_i} = p_i^{-a_i}$
pour chaque $i$, tandis que $\abs{x}_q = 1$ pour tout premier $q \notin \{p_1, \ldots, p_k\}$.
Le produit vaut donc $p_1^{a_1} \cdots p_k^{a_k} \cdot p_1^{-a_1} \cdots p_k^{-a_k} = 1$.
\end{proof}

\index{valeur absolue p-adique@valeur absolue $p$-adique|)}

\section{Métrique \texorpdfstring{$p$}{p}-adique et convergence}
\label{sec:metrique-p-adique}
\index{métrique p-adique@métrique $p$-adique}

\begin{definition}[Distance $p$-adique]
\label{def:distance-p-adique}
La \emph{distance $p$-adique} sur $\Q$ est définie par
\[
  d_p(x, y) = \abs{x - y}_p \quad \text{pour tout } x, y \in \Q.
\]
Le couple $(\Q, d_p)$ est un espace métrique\index{espace métrique} (ultramétrique).
\end{definition}

\begin{exemple}[Convergence $p$-adique]
\label{ex:convergence-p-adique}
Dans $(\Q, d_5)$, la suite $(a_n)$ définie par $a_n = \sum_{k=0}^{n} 5^k$ converge,
car $\abs{a_{n+1} - a_n}_5 = \abs{5^{n+1}}_5 = 5^{-(n+1)} \to 0$. En fait,
\[
  \sum_{k=0}^{\infty} 5^k = \frac{1}{1 - 5} = -\frac{1}{4}
\]
dans $\Q_5$. La série géométrique converge pour $\abs{5}_5 = 1/5 < 1$.
\end{exemple}

\begin{remarque}
\label{rem:tout-triangle-isocele}
Dans un espace ultramétrique, si $\abs{x}_p \neq \abs{y}_p$, alors
$\abs{x + y}_p = \max(\abs{x}_p, \abs{y}_p)$. Géométriquement : dans tout triangle
$p$-adique, les deux plus grands côtés sont égaux.
\end{remarque}

\section{Entiers \texorpdfstring{$p$}{p}-adiques et corps \texorpdfstring{$\Q_p$}{Qp}}
\label{sec:entiers-p-adiques}
\index{entiers p-adiques@entiers $p$-adiques|(}

\begin{definition}[Corps $\Q_p$ des nombres $p$-adiques]
\label{def:Qp}
Le corps $\Q_p$\index{Qp@$\Q_p$} est le complété de $\Q$ par rapport à la
distance $d_p$. C'est un corps valué complet contenant $\Q$ comme sous-corps
dense.
\end{definition}

\begin{definition}[Anneau $\Z_p$ des entiers $p$-adiques]
\label{def:Zp}
L'anneau des \emph{entiers $p$-adiques}\index{Zp@$\Z_p$} est
\[
  \Z_p = \bigl\{ x \in \Q_p : \abs{x}_p \leq 1 \bigr\}
  = \bigl\{ x \in \Q_p : v_p(x) \geq 0 \bigr\}.
\]
C'est un anneau local\index{anneau local} d'idéal maximal $p\Z_p$.
\end{definition}

\begin{proposition}[Représentation en série]
\label{prop:serie-p-adique}
\index{série p-adique@série $p$-adique}
Tout élément $x \in \Z_p$ s'écrit de manière unique sous la forme
\[
  x = \sum_{i=0}^{\infty} a_i \, p^i, \quad a_i \in \{0, 1, \ldots, p-1\}.
\]
Plus généralement, tout $x \in \Q_p^*$ s'écrit
\[
  x = \sum_{i=v_p(x)}^{\infty} a_i \, p^i, \quad a_i \in \{0, 1, \ldots, p-1\},
  \quad a_{v_p(x)} \neq 0.
\]
\end{proposition}

\begin{exemple}[Calculs dans $\Z_5$]
\label{ex:calculs-Z5}
Représentons $-1$ dans $\Z_5$. On cherche $x = \sum a_i \cdot 5^i$ tel que $x + 1 = 0$.
On a $-1 = (p-1) + (p-1)p + (p-1)p^2 + \cdots$, soit
\[
  -1 = 4 + 4 \cdot 5 + 4 \cdot 5^2 + 4 \cdot 5^3 + \cdots = \sum_{i=0}^{\infty} 4 \cdot 5^i.
\]
Vérification : $1 + (4 + 4 \cdot 5 + 4 \cdot 5^2 + \cdots) = 5 + 4 \cdot 5 + 4 \cdot 5^2 + \cdots
= 5(1 + 4 + 4 \cdot 5 + \cdots)$. En itérant, tous les chiffres s'annulent,
donnant $0$.
\end{exemple}

\begin{exemple}[Représentation de $1/3$ dans $\Z_5$]
\label{ex:un-tiers-Z5}
On cherche $x \in \Z_5$ tel que $3x = 1$. On résout par développement successif :
$3 a_0 \equiv 1 \pmod{5}$ donne $a_0 = 2$. Puis $3(2 + a_1 \cdot 5 + \cdots) = 1$,
soit $6 + 3a_1 \cdot 5 + \cdots = 1$, d'où $(6-1)/5 = 1$ et $3a_1 \equiv -1 \pmod{5}$,
soit $a_1 = 3$. En continuant :
\[
  \frac{1}{3} = 2 + 3 \cdot 5 + 1 \cdot 5^2 + 3 \cdot 5^3 + 1 \cdot 5^4 + \cdots
\]
Les chiffres alternent périodiquement entre $3$ et $1$ (après le premier terme).
\end{exemple}

\begin{proposition}[Structure de $\Z_p$]
\label{prop:structure-Zp}
\begin{enumerate}[label=(\roman*)]
  \item $\Z_p$ est un anneau intègre, local, principal.
  \item Les idéaux de $\Z_p$ sont exactement les $p^n \Z_p$ pour $n \geq 0$, et
    l'idéal nul.
  \item $\Z_p / p^n \Z_p \cong \Z / p^n \Z$ pour tout $n \geq 1$.
  \item $\Z_p^{\times} = \{ x \in \Z_p : \abs{x}_p = 1 \} = \Z_p \setminus p\Z_p$.
\end{enumerate}
\end{proposition}

\index{entiers p-adiques@entiers $p$-adiques|)}

% --- Illustration TikZ : arbre p-adique ---
\begin{figure}[ht]
\centering
\begin{tikzpicture}[
  level distance=1.4cm,
  sibling distance=3.5cm,
  edge from parent/.style={draw, -Stealth, thick},
  every node/.style={draw, rounded corners, fill=blue!8, minimum width=1.2cm,
    minimum height=0.6cm, font=\small},
  level 1/.style={sibling distance=3.5cm},
  level 2/.style={sibling distance=1.2cm},
]
  \node {$\Z_3$}
    child { node {$0 + 3\Z_3$}
      child { node {$0$} }
      child { node {$3$} }
      child { node {$6$} }
    }
    child { node {$1 + 3\Z_3$}
      child { node {$1$} }
      child { node {$4$} }
      child { node {$7$} }
    }
    child { node {$2 + 3\Z_3$}
      child { node {$2$} }
      child { node {$5$} }
      child { node {$8$} }
    };
\end{tikzpicture}
\caption{Arbre des classes de $\Z_3$ modulo $3$ et $9$ :
chaque élément de $\Z_3$ est déterminé par ses chiffres successifs en base $3$.}
\label{fig:arbre-3-adique}
\end{figure}

% --- Illustration TikZ : disques p-adiques ---
\begin{figure}[ht]
\centering
\begin{tikzpicture}
  % Disque principal
  \fill[blue!8] (0,0) circle (3);
  \draw[thick, blue!60!black] (0,0) circle (3);
  \node[font=\small] at (0, 3.3) {$\Z_p = \{x : \abs{x}_p \leq 1\}$};

  % Sous-disques
  \fill[red!10] (0,0) circle (2);
  \draw[thick, red!60!black] (0,0) circle (2);
  \node[font=\small] at (0, 2.3) {$p\Z_p$};

  \fill[green!10] (0,0) circle (1.2);
  \draw[thick, green!60!black] (0,0) circle (1.2);
  \node[font=\small] at (0, 0) {$p^2\Z_p$};
\end{tikzpicture}
\caption{Structure emboîtée des disques $p$-adiques : $\cdots \subset p^2\Z_p \subset p\Z_p \subset \Z_p$.}
\label{fig:disques-p-adiques}
\end{figure}

\section{Lemme de Hensel}
\label{sec:hensel}
\index{lemme de Hensel|(}

Le lemme de Hensel\index{Hensel, lemme de} est l'analogue $p$-adique du théorème
des fonctions implicites : il permet de <<~relever~>> une solution approchée en
une solution exacte.

\begin{theoreme}[Lemme de Hensel]
\label{thm:hensel}
Soit $f(x) \in \Z_p[x]$ un polynôme à coefficients dans $\Z_p$. Supposons qu'il
existe $a_0 \in \Z_p$ tel que
\[
  v_p\bigl(f(a_0)\bigr) > 2\, v_p\bigl(f'(a_0)\bigr).
\]
Alors il existe un unique $a \in \Z_p$ tel que $f(a) = 0$ et
$\abs{a - a_0}_p \leq \abs{f(a_0)/f'(a_0)^2}_p^{1/2}$.
\end{theoreme}

\begin{proof}[Preuve (cas simple : $f'(a_0) \not\equiv 0 \pmod{p}$)]
On suppose $v_p(f'(a_0)) = 0$ (c'est-à-dire $f'(a_0) \in \Z_p^{\times}$) et
$v_p(f(a_0)) \geq 1$. On construit une suite $(a_n)$ par la méthode de Newton
$p$-adique :
\[
  a_{n+1} = a_n - \frac{f(a_n)}{f'(a_n)}.
\]

\textbf{Étape 1 : Bonne définition.}
Puisque $f'(a_0) \in \Z_p^{\times}$, la division est bien définie dans $\Q_p$.

\textbf{Étape 2 : Convergence quadratique.}
Par le développement de Taylor dans $\Z_p[x]$ :
\[
  f(a_{n+1}) = f(a_n) + f'(a_n)(a_{n+1} - a_n) + c(a_{n+1} - a_n)^2
\]
pour un certain $c \in \Z_p$. Or $a_{n+1} - a_n = -f(a_n)/f'(a_n)$, d'où
\[
  f(a_{n+1}) = f(a_n) - f(a_n) + c \cdot \frac{f(a_n)^2}{f'(a_n)^2}
  = c \cdot \frac{f(a_n)^2}{f'(a_n)^2}.
\]
Ainsi $v_p(f(a_{n+1})) \geq 2\, v_p(f(a_n))$, ce qui donne la convergence
quadratique $v_p(f(a_n)) \geq 2^n$.

\textbf{Étape 3 : $f'(a_n) \in \Z_p^{\times}$ pour tout $n$.}
On montre par récurrence que $a_{n+1} \equiv a_n \pmod{p}$, d'où
$f'(a_{n+1}) \equiv f'(a_n) \pmod{p}$ et $v_p(f'(a_n)) = 0$ pour tout $n$.

\textbf{Étape 4 : Convergence et unicité.}
La suite $(a_n)$ est de Cauchy dans $\Z_p$ (car $\abs{a_{n+1}-a_n}_p \to 0$),
donc elle converge vers un $a \in \Z_p$ vérifiant $f(a) = 0$. L'unicité dans le
disque considéré découle de l'estimation.
\end{proof}

\begin{corollaire}[Version modulo $p$]
\label{cor:hensel-mod-p}
Si $f(x) \in \Z[x]$ et $a_0 \in \Z$ vérifient $f(a_0) \equiv 0 \pmod{p}$ et
$f'(a_0) \not\equiv 0 \pmod{p}$, alors il existe $a \in \Z_p$ tel que $f(a) = 0$
et $a \equiv a_0 \pmod{p}$.
\end{corollaire}

\index{lemme de Hensel|)}

\section{Applications du lemme de Hensel}
\label{sec:applications-hensel}

\subsection{Résolution de $x^2 = -1$ dans $\Q_5$}
\label{subsec:x2-moins1-Q5}
\index{racine carrée p-adique@racine carrée $p$-adique}

\begin{exemple}[{$x^2 + 1 = 0$ dans $\Q_5$}]
\label{ex:x2-plus-1-Q5}
Cherchons $x \in \Z_5$ tel que $x^2 = -1$.

\textbf{Étape 1 :} Solution modulo $5$. On vérifie $2^2 = 4 \equiv -1 \pmod{5}$.
Prenons $a_0 = 2$.

\textbf{Étape 2 :} Vérification des hypothèses de Hensel. Soit $f(x) = x^2 + 1$.
On a $f(2) = 5 \equiv 0 \pmod{5}$ et $f'(2) = 4 \not\equiv 0 \pmod{5}$.

\textbf{Étape 3 :} Application du lemme de Hensel (Corollaire~\ref{cor:hensel-mod-p}).
Il existe $a \in \Z_5$ tel que $a^2 = -1$ et $a \equiv 2 \pmod{5}$.

\textbf{Relèvement explicite :}
On cherche $a_1$ tel que $a_1 \equiv 2 \pmod{5}$ et $a_1^2 \equiv -1 \pmod{25}$.
La méthode de Newton donne
\[
  a_1 = a_0 - \frac{f(a_0)}{f'(a_0)} = 2 - \frac{5}{4}.
\]
Or $4^{-1} \equiv 4 \pmod{5}$, donc $5/4 \equiv 5 \cdot 4 = 20 \pmod{25}$, et
$a_1 = 2 - 20 = -18 \equiv 7 \pmod{25}$. Vérifions : $7^2 = 49 \equiv -1 \pmod{25}$.

On obtient les premiers chiffres : $i = 2 + 1 \cdot 5 + \cdots$ dans $\Z_5$,
où $i^2 = -1$.
\end{exemple}

\begin{remarque}
\label{rem:x2-moins1-Q3}
En revanche, $x^2 = -1$ n'a pas de solution dans $\Q_3$, car $-1$ n'est pas un
carré modulo $3$ (les carrés modulo $3$ sont $\{0, 1\}$).
\end{remarque}

\subsection{Principe local-global}
\label{subsec:local-global}
\index{Hasse--Minkowski, théorème de}

\begin{theoreme}[Hasse--Minkowski]
\label{thm:hasse-minkowski}
Soit $Q(x_1, \ldots, x_n) = \sum_{i,j} a_{ij} x_i x_j$ une forme quadratique
à coefficients dans $\Q$. Alors $Q$ représente $0$ dans $\Q$ (c'est-à-dire
il existe $(x_1, \ldots, x_n) \in \Q^n \setminus \{0\}$ tel que $Q(x_1, \ldots, x_n) = 0$)
si et seulement si $Q$ représente $0$ dans $\R$ et dans $\Q_p$ pour tout premier $p$.
\end{theoreme}

\begin{remarque}
\label{rem:local-global-limitation}
Le théorème de Hasse--Minkowski ne se généralise pas aux formes de degré supérieur.
Par exemple, Selmer\index{Selmer, Ernst} a montré que l'équation $3x^3 + 4y^3 + 5z^3 = 0$
a des solutions non triviales dans $\R$ et dans tous les $\Q_p$, mais pas dans $\Q$.
\end{remarque}

\section{Théorème d'Ostrowski}
\label{sec:ostrowski}
\index{Ostrowski, théorème d'}

\begin{theoreme}[Ostrowski]
\label{thm:ostrowski}
Toute valeur absolue non triviale sur $\Q$ est équivalente soit à la valeur
absolue usuelle $\abs{\cdot}_\infty$, soit à une valeur absolue $p$-adique
$\abs{\cdot}_p$ pour un certain nombre premier $p$.
\end{theoreme}

\begin{remarque}
\label{rem:ostrowski-places}
Ce théorème signifie que les <<~places~>>\index{place} de $\Q$ --- c'est-à-dire
les classes d'équivalence de valeurs absolues non triviales --- sont naturellement
indexées par $\{2, 3, 5, 7, 11, \ldots\} \cup \{\infty\}$. Chaque complétion
correspondante donne soit $\R$ (pour $\infty$), soit $\Q_p$ (pour le premier $p$).
\end{remarque}

\section{Exemples détaillés}
\label{sec:p-adic-exemples}

\begin{exemple}[Calculs dans $\Z_7$]
\label{ex:calculs-Z7}
Représentons $1/2$ dans $\Z_7$. On cherche $x = \sum a_i \cdot 7^i$ tel que $2x = 1$.
\begin{itemize}
  \item $2a_0 \equiv 1 \pmod{7}$ : $a_0 = 4$ (car $2 \cdot 4 = 8 \equiv 1$).
  \item $2(4 + a_1 \cdot 7) \equiv 1 \pmod{49}$ : $8 + 14a_1 \equiv 1 \pmod{49}$,
    soit $14a_1 \equiv -7 \pmod{49}$, $2a_1 \equiv -1 \equiv 6 \pmod{7}$, $a_1 = 3$.
  \item En continuant : $a_2 = 3$, $a_3 = 3$, \ldots
\end{itemize}
On obtient $1/2 = 4 + 3 \cdot 7 + 3 \cdot 7^2 + 3 \cdot 7^3 + \cdots = 4 + 3 \cdot \frac{7}{1 - 7}$
dans $\Z_7$, ce qui est cohérent car $4 + 3 \cdot 7 / (1-7) = 4 - 7/2 = 1/2$.
\end{exemple}

\begin{exemple}[{Résolution de $x^2 = 2$ dans $\Q_7$}]
\label{ex:x2-egal-2-Q7}
On a $3^2 = 9 \equiv 2 \pmod{7}$. De plus, $f(x) = x^2 - 2$, $f'(3) = 6 \not\equiv 0
\pmod{7}$. Par le lemme de Hensel, $\sqrt{2} \in \Z_7$ et $\sqrt{2} \equiv 3 \pmod{7}$.

Relèvement : $a_1 = 3 - (9-2)/(2 \cdot 3) = 3 - 7/6$. Or $6^{-1} \equiv 6 \pmod{7}$,
donc $7/6 \equiv 7 \cdot 6 = 42 \equiv 0 \pmod{7}$, soit $a_1 = 3 + 1 \cdot 7 = 10$.
Vérifions : $10^2 = 100 = 2 + 2 \cdot 49$, donc $v_7(10^2 - 2) = v_7(98) = 2 \geq 2$.
\end{exemple}

\section{Résumé du chapitre}
\label{sec:resume-ch9}

\begin{center}
\begin{tabular}{lp{9cm}}
\toprule
\textbf{Concept} & \textbf{Description} \\
\midrule
Valuation $v_p(n)$ & Plus grande puissance de $p$ divisant $n$ \\
$\abs{x}_p = p^{-v_p(x)}$ & Valeur absolue $p$-adique \\
Inégalité ultramétrique & $\abs{x+y}_p \leq \max(\abs{x}_p, \abs{y}_p)$ \\
$\Z_p$ & Entiers $p$-adiques : séries $\sum a_i p^i$, $a_i \in \{0,\ldots,p-1\}$ \\
$\Q_p$ & Corps des nombres $p$-adiques, complétion de $(\Q, d_p)$ \\
Lemme de Hensel & Relèvement de racines modulo $p$ en racines $p$-adiques \\
Ostrowski & Classification de toutes les valeurs absolues sur $\Q$ \\
\bottomrule
\end{tabular}
\end{center}

\section{Exercices}
\label{sec:exercices-ch9}

\begin{exercice}[$\star$ --- Calculs de valuations]
\label{exo:calculs-valuations}
Calculer les valuations suivantes :
\begin{enumerate}[label=(\alph*)]
  \item $v_2(360)$, $v_3(360)$, $v_5(360)$.
  \item $v_7(50!/7)$.
  \item $v_p(n!)$ en fonction de $n$ et $p$ (formule de Legendre\index{formule de Legendre}).
\end{enumerate}
\end{exercice}

\begin{exercice}[$\star$ --- Valeur absolue $p$-adique]
\label{exo:valeur-absolue-p-adique}
\begin{enumerate}[label=(\alph*)]
  \item Calculer $\abs{63}_3$, $\abs{140}_7$, $\abs{-250}_5$.
  \item Montrer que $\abs{n}_p = 1$ pour tout $n \in \Z$ avec $\pgcd(n,p)=1$.
  \item Montrer que pour tout $n \in \Z \setminus \{0\}$, il n'existe qu'un nombre
    fini de premiers $p$ tels que $\abs{n}_p \neq 1$.
\end{enumerate}
\end{exercice}

\begin{exercice}[$\star\star$ --- Propriétés ultramétriques]
\label{exo:proprietes-ultrametriques}
Soit $(K, \abs{\cdot})$ un corps muni d'une valeur absolue ultramétrique.
\begin{enumerate}[label=(\alph*)]
  \item Montrer que si $\abs{x} \neq \abs{y}$, alors $\abs{x+y} = \max(\abs{x}, \abs{y})$.
  \item En déduire que dans un espace ultramétrique, tout point d'une boule ouverte
    en est le centre.
  \item Montrer que toute boule ouverte est aussi un ensemble fermé.
\end{enumerate}
\end{exercice}

\begin{exercice}[$\star\star$ --- Développements $p$-adiques]
\label{exo:developpements-p-adiques}
\begin{enumerate}[label=(\alph*)]
  \item Trouver le développement $3$-adique de $-1$, $1/2$ et $2/5$.
  \item Montrer que $x \in \Z_p$ est rationnel si et seulement si son développement
    $p$-adique est ultimement périodique.
\end{enumerate}
\end{exercice}

\begin{exercice}[$\star\star\star$ --- Applications de Hensel]
\label{exo:applications-hensel}
\begin{enumerate}[label=(\alph*)]
  \item Montrer que $\sqrt{-1} \in \Q_p$ si et seulement si $p \equiv 1 \pmod{4}$ ou $p = 2$.
    \textit{Indication :} pour $p=2$, vérifier qu'il n'y a pas de racine modulo $8$.
  \item Soit $p$ un premier impair et $a \in \Z$ avec $\pgcd(a,p) = 1$. Montrer que
    $a$ est un carré dans $\Z_p$ si et seulement si $\legendre{a}{p} = 1$.
  \item Montrer que $\sqrt{6} \in \Q_5$ mais $\sqrt{6} \notin \Q_7$.
\end{enumerate}
\end{exercice}

\begin{exercice}[$\star\star\star$ --- Formule de Legendre]
\label{exo:formule-legendre}
\index{formule de Legendre}
Démontrer la formule de Legendre : pour tout premier $p$ et tout entier $n \geq 1$,
\[
  v_p(n!) = \sum_{k=1}^{\infty} \floor{\frac{n}{p^k}}
  = \frac{n - s_p(n)}{p - 1},
\]
où $s_p(n)$ désigne la somme des chiffres de $n$ en base $p$.
\end{exercice}

\index{nombres p-adiques@nombres $p$-adiques|)}

%%%%%%%%%%%%%%%%%%%%%%%%%%%%%%%%%%%%%%%%%%%%%%%%%%%%%%%%%%%%%%%%%%%%
%%  CHAPITRE 10 — Aperçu sur la Théorie Analytique des Nombres
%%%%%%%%%%%%%%%%%%%%%%%%%%%%%%%%%%%%%%%%%%%%%%%%%%%%%%%%%%%%%%%%%%%%
\chapter{Aperçu sur la Théorie Analytique des Nombres}
\label{chap:analytique}
\index{théorie analytique des nombres|(}

\begin{center}
\itshape
<<~Les nombres premiers sont les atomes de l'arithmétique,\\
et la fonction zêta est le spectre qui révèle leur structure.~>>\\[4pt]
--- Bernhard Riemann (1826--1866)
\end{center}

\bigskip

\section{Introduction historique}
\label{sec:analytique-intro}

La théorie analytique des nombres est née de l'idée audacieuse d'utiliser les
outils de l'analyse --- séries, intégrales, fonctions de variable complexe ---
pour étudier les propriétés des nombres entiers et, en particulier, la distribution
des nombres premiers.

\begin{itemize}
  \item \textbf{Euler\index{Euler, Leonhard} (1737)} établit le produit eulérien
    $\sum n^{-s} = \prod_p (1 - p^{-s})^{-1}$ et en déduit une nouvelle preuve
    de l'infinité des premiers.
  \item \textbf{Dirichlet\index{Dirichlet, Johann Peter Gustav Lejeune} (1837)}
    introduit les $L$-fonctions et démontre que toute progression arithmétique
    $a, a+q, a+2q, \ldots$ avec $\pgcd(a,q)=1$ contient une infinité de premiers.
  \item \textbf{Riemann\index{Riemann, Bernhard} (1859)} publie son mémoire
    fondateur <<~Über die Anzahl der Primzahlen unter einer gegebenen Grösse~>>,
    introduisant la prolongation analytique de $\zeta$ et formulant sa célèbre
    hypothèse.
  \item \textbf{Hadamard\index{Hadamard, Jacques}} et
    \textbf{de la Vallée-Poussin\index{de la Vallée-Poussin, Charles-Jean}}
    (1896) démontrent indépendamment le théorème des nombres premiers.
\end{itemize}

\section{La fonction zêta de Riemann}
\label{sec:zeta}
\index{fonction zêta de Riemann|(}

\begin{definition}[Fonction zêta de Riemann]
\label{def:zeta}
La \emph{fonction zêta de Riemann}\index{zêta, fonction} est définie pour
$s \in \C$ avec $\operatorname{Re}(s) > 1$ par la série de Dirichlet
\[
  \zeta(s) = \sum_{n=1}^{\infty} \frac{1}{n^s}.
\]
\end{definition}

\begin{proposition}[Convergence]
\label{prop:convergence-zeta}
La série $\sum_{n=1}^{\infty} n^{-s}$ converge absolument pour
$\operatorname{Re}(s) > 1$ et y définit une fonction holomorphe.
\end{proposition}

\begin{proof}
Pour $s = \sigma + it$ avec $\sigma > 1$, on a $\abs{n^{-s}} = n^{-\sigma}$.
La série $\sum n^{-\sigma}$ est une série de Riemann convergente pour $\sigma > 1$.
La convergence est uniforme sur tout demi-plan $\{\operatorname{Re}(s) \geq \sigma_0\}$
avec $\sigma_0 > 1$, ce qui assure l'holomorphie par le théorème de Weierstrass.
\end{proof}

\subsection{Produit eulérien}
\label{subsec:euler-product}
\index{produit eulérien|(}

\begin{theoreme}[Produit eulérien d'Euler]
\label{thm:euler-product}
Pour tout $s \in \C$ avec $\operatorname{Re}(s) > 1$,
\[
  \zeta(s) = \prod_{p \text{ premier}} \frac{1}{1 - p^{-s}}.
\]
\end{theoreme}

\begin{proof}
Fixons $s$ avec $\sigma = \operatorname{Re}(s) > 1$. Pour tout $N \geq 1$, notons
$P_N = \{p_1, \ldots, p_k\}$ l'ensemble des premiers $\leq N$. Développons le
produit fini :
\[
  \prod_{j=1}^{k} \frac{1}{1 - p_j^{-s}}
  = \prod_{j=1}^{k} \sum_{e_j=0}^{\infty} p_j^{-e_j s}
  = \sum_{\substack{n \geq 1 \\ p \mid n \Rightarrow p \leq N}} \frac{1}{n^s},
\]
où la dernière égalité résulte du théorème fondamental de l'arithmétique :
chaque entier dont tous les facteurs premiers sont $\leq N$ apparaît exactement
une fois dans le développement.

Or, la somme du membre droit contient au moins tous les termes $1/n^s$ pour
$n \leq N$ (car si $n \leq N$, tous ses facteurs premiers sont $\leq N$). Ainsi :
\[
  \left| \zeta(s) - \prod_{j=1}^{k} \frac{1}{1 - p_j^{-s}} \right|
  \leq \sum_{n > N} \frac{1}{n^\sigma}.
\]
Quand $N \to \infty$, le reste de la série convergente $\sum n^{-\sigma}$ tend
vers $0$, d'où la convergence du produit infini vers $\zeta(s)$.
\end{proof}

\begin{corollaire}[Infinité des nombres premiers]
\label{cor:euler-infinité-premiers}
Il existe une infinité de nombres premiers.
\end{corollaire}

\begin{proof}
Si l'ensemble des premiers était fini, le produit eulérien serait un produit fini,
donc fini pour $s = 1$. Or $\zeta(1) = \sum 1/n = +\infty$, contradiction.
\end{proof}

\begin{remarque}
\label{rem:euler-quantitatif}
Plus quantitativement, Euler a montré que $\sum_{p} 1/p = +\infty$, ce qui
donne une mesure de la <<~densité~>> des premiers : ils sont non seulement
infinis, mais suffisamment nombreux pour que la somme de leurs inverses diverge.
\end{remarque}

\index{produit eulérien|)}

\subsection{Valeurs spéciales de $\zeta$}
\label{subsec:valeurs-speciales-zeta}

\begin{theoreme}[{$\zeta(2) = \pi^2/6$} (Euler, 1735)]
\label{thm:zeta-2}
\index{problème de Bâle}
On a
\[
  \zeta(2) = \sum_{n=1}^{\infty} \frac{1}{n^2} = \frac{\pi^2}{6}.
\]
\end{theoreme}

\begin{proof}[Esquisse de preuve (méthode d'Euler)]
On considère la fonction $\sin(\pi x)/(\pi x) = \prod_{n=1}^{\infty}(1 - x^2/n^2)$.
En développant le produit et en identifiant le coefficient de $x^2$, on obtient
\[
  -\frac{\pi^2}{6} \cdot \frac{1}{\pi^2} = - \sum_{n=1}^{\infty} \frac{1}{n^2},
\]
en comparant avec le développement en série $\sin(\pi x)/(\pi x) = 1 - \pi^2 x^2/6 + \cdots$,
d'où le résultat. La justification rigoureuse du produit infini utilise le
théorème de factorisation de Weierstrass.
\end{proof}

\begin{remarque}
\label{rem:zeta-pairs}
Plus généralement, Euler a calculé $\zeta(2k)$ pour tout $k \geq 1$ en termes
des nombres de Bernoulli\index{nombres de Bernoulli} :
\[
  \zeta(2k) = \frac{(-1)^{k+1} (2\pi)^{2k} B_{2k}}{2 \cdot (2k)!}.
\]
En revanche, les valeurs $\zeta(2k+1)$ (et en particulier $\zeta(3)$, la
\emph{constante d'Apéry}\index{constante d'Apéry}) restent mystérieuses :
$\zeta(3)$ est irrationnel (Apéry, 1978), mais on ignore si $\zeta(5)$ l'est.
\end{remarque}

\index{fonction zêta de Riemann|)}

\section{\texorpdfstring{$L$}{L}-fonctions de Dirichlet}
\label{sec:L-fonctions}
\index{L-fonction de Dirichlet@$L$-fonction de Dirichlet|(}

\begin{definition}[Caractère de Dirichlet]
\label{def:caractere-dirichlet}
\index{caractère de Dirichlet}
Un \emph{caractère de Dirichlet} modulo $q$ est un morphisme de groupes
$\chi : (\Z/q\Z)^{\times} \to \C^{\times}$, étendu à $\Z$ par périodicité et
en posant $\chi(n) = 0$ si $\pgcd(n,q) > 1$.
\end{definition}

\begin{definition}[$L$-fonction de Dirichlet]
\label{def:L-fonction}
Pour un caractère $\chi$ modulo $q$, la \emph{$L$-fonction de Dirichlet} associée est
\[
  L(s, \chi) = \sum_{n=1}^{\infty} \frac{\chi(n)}{n^s}, \quad \operatorname{Re}(s) > 1.
\]
Elle admet un produit eulérien :
\[
  L(s, \chi) = \prod_{p \text{ premier}} \frac{1}{1 - \chi(p)\, p^{-s}}.
\]
\end{definition}

\begin{remarque}
\label{rem:L-fonction-zeta}
Pour le caractère trivial $\chi_0$ modulo $1$ (c'est-à-dire $\chi_0(n) = 1$
pour tout $n$), on retrouve $L(s, \chi_0) = \zeta(s)$.
\end{remarque}

\index{L-fonction de Dirichlet@$L$-fonction de Dirichlet|)}

\section{Théorème de Dirichlet sur les progressions arithmétiques}
\label{sec:dirichlet-progressions}
\index{Dirichlet, théorème de|(}

\begin{theoreme}[Dirichlet, 1837]
\label{thm:dirichlet-progressions}
Soient $a, q \in \Z$ avec $q \geq 1$ et $\pgcd(a, q) = 1$. Alors la progression
arithmétique
\[
  a, \; a + q, \; a + 2q, \; a + 3q, \; \ldots
\]
contient une infinité de nombres premiers. Plus précisément,
\[
  \sum_{\substack{p \text{ premier} \\ p \equiv a \!\!\pmod{q}}} \frac{1}{p} = +\infty.
\]
\end{theoreme}

\begin{remarque}
\label{rem:dirichlet-preuve}
La preuve de Dirichlet repose sur deux ingrédients principaux :
\begin{enumerate}[label=(\roman*)]
  \item L'utilisation des caractères de Dirichlet pour <<~isoler~>> la
    progression arithmétique via la relation d'orthogonalité :
    \[
      \frac{1}{\varphi(q)} \sum_{\chi \bmod q} \chi(n) \overline{\chi(a)}
      = \begin{cases} 1 & \text{si } n \equiv a \pmod{q}, \\ 0 & \text{sinon.} \end{cases}
    \]
  \item La non-annulation $L(1, \chi) \neq 0$ pour tout caractère $\chi \neq \chi_0$
    modulo $q$, qui est le c\oe ur technique de la preuve.
\end{enumerate}
\end{remarque}

\index{Dirichlet, théorème de|)}

\section{Le théorème des nombres premiers}
\label{sec:TNP}
\index{théorème des nombres premiers|(}

\begin{notation}
\label{not:pi-x}
On note $\pi(x) = \#\{p \leq x : p \text{ premier}\}$\index{pi(x)@$\pi(x)$}
la fonction de comptage des nombres premiers.
\end{notation}

\begin{theoreme}[Théorème des nombres premiers (Hadamard, de la Vallée-Poussin, 1896)]
\label{thm:TNP}
\[
  \pi(x) \sim \frac{x}{\ln x} \quad \text{quand } x \to \infty,
\]
c'est-à-dire $\displaystyle\lim_{x \to \infty} \frac{\pi(x)}{x / \ln x} = 1$.
\end{theoreme}

\begin{remarque}
\label{rem:TNP-historique}
Ce théorème, conjecturé indépendamment par Legendre\index{Legendre, Adrien-Marie}
(vers 1798) et Gauss\index{Gauss, Carl Friedrich} (vers 1793), fut démontré
indépendamment en 1896 par Hadamard et de la Vallée-Poussin. Les deux preuves
utilisent de manière essentielle la non-annulation de $\zeta(s)$ sur la droite
$\operatorname{Re}(s) = 1$.
\end{remarque}

\begin{definition}[Logarithme intégral]
\label{def:Li}
\index{logarithme intégral}
Le \emph{logarithme intégral} est défini par
\[
  \operatorname{Li}(x) = \int_2^{x} \frac{dt}{\ln t}.
\]
C'est une meilleure approximation de $\pi(x)$ que $x/\ln x$ :
$\pi(x) \sim \operatorname{Li}(x)$.
\end{definition}

% --- Figure TikZ : comparaison π(x), x/ln(x), Li(x) ---
\begin{figure}[ht]
\centering
\begin{tikzpicture}[scale=0.85]
  \begin{scope}
    \draw[->] (0,0) -- (10.5,0) node[right] {$x$};
    \draw[->] (0,0) -- (0,7) node[above] {$y$};

    % Axes labels
    \foreach \x/\lab in {2/20, 4/40, 6/60, 8/80, 10/100} {
      \draw (\x, 0.1) -- (\x, -0.1) node[below, font=\footnotesize] {$\lab$};
    }
    \foreach \y/\lab in {1/5, 2/10, 3/15, 4/20, 5/25} {
      \draw (0.1, \y) -- (-0.1, \y) node[left, font=\footnotesize] {$\lab$};
    }

    % π(x) - step function (scaled: x -> x/10, y -> y/5)
    % Values: π(10)=4, π(20)=8, π(30)=10, π(40)=12, π(50)=15,
    %         π(60)=17, π(70)=19, π(80)=22, π(90)=24, π(100)=25
    \draw[thick, blue] (0.2,0) -- (0.2,0.2) -- (0.3,0.2) -- (0.3,0.4)
      -- (0.5,0.4) -- (0.5,0.6) -- (0.7,0.6) -- (0.7,0.8)
      -- (1.1,0.8) -- (1.1,1) -- (1.3,1) -- (1.3,1.2)
      -- (1.7,1.2) -- (1.7,1.4) -- (1.9,1.4) -- (1.9,1.6);
    \draw[thick, blue] (1.9,1.6) -- (2.3,1.6) -- (2.3,1.8) -- (2.9,1.8)
      -- (2.9,2) -- (3.1,2) -- (3.1,2.2) -- (3.7,2.2)
      -- (3.7,2.4) -- (4.1,2.4) -- (4.1,2.6) -- (4.3,2.6)
      -- (4.3,2.8) -- (4.7,2.8) -- (4.7,3);
    \draw[thick, blue] (4.7,3) -- (5.3,3) -- (5.3,3.2) -- (5.9,3.2)
      -- (5.9,3.4) -- (6.1,3.4) -- (6.1,3.6) -- (6.7,3.6)
      -- (6.7,3.8) -- (7.1,3.8) -- (7.1,4) -- (7.3,4)
      -- (7.3,4.2) -- (7.9,4.2) -- (7.9,4.4)
      -- (8.3,4.4) -- (8.3,4.6) -- (8.9,4.6) -- (8.9,4.8)
      -- (9.7,4.8) -- (9.7,5) -- (10,5);

    % x/ln(x) curve (scaled)
    \draw[thick, red, dashed, domain=0.5:10, samples=40]
      plot (\x, {\x / (ln(\x*10) ) * (10/5)});

    % Li(x) approximation curve (scaled)
    \draw[thick, green!60!black, densely dotted, domain=0.5:10, samples=40]
      plot (\x, {\x / (ln(\x*10) - 1.08) * (10/5)});

    % Legend
    \draw[thick, blue] (6,6.5) -- (7,6.5) node[right, font=\small] {$\pi(x)$};
    \draw[thick, red, dashed] (6,6) -- (7,6) node[right, font=\small] {$x/\ln x$};
    \draw[thick, green!60!black, densely dotted] (6,5.5) -- (7,5.5)
      node[right, font=\small] {$\operatorname{Li}(x)$};
  \end{scope}
\end{tikzpicture}
\caption{Comparaison de $\pi(x)$ avec ses approximations $x/\ln x$ et
$\operatorname{Li}(x)$ pour $x \leq 100$.}
\label{fig:pi-x-comparaison}
\end{figure}

\section{Fonctions de Tchebychev}
\label{sec:tchebychev}
\index{fonctions de Tchebychev|(}

\begin{definition}[Fonctions $\theta$ et $\psi$ de Tchebychev]
\label{def:tchebychev}
On définit :
\[
  \theta(x) = \sum_{\substack{p \leq x \\ p \text{ premier}}} \ln p,
  \qquad
  \psi(x) = \sum_{\substack{p^k \leq x \\ p \text{ premier}, \, k \geq 1}} \ln p
  = \sum_{n \leq x} \Lambda(n),
\]
où $\Lambda$ est la \emph{fonction de von Mangoldt}\index{fonction de von Mangoldt} :
\[
  \Lambda(n) = \begin{cases} \ln p & \text{si } n = p^k \text{ pour un premier } p
    \text{ et } k \geq 1, \\ 0 & \text{sinon.} \end{cases}
\]
\end{definition}

\begin{proposition}[Relation avec le TNP]
\label{prop:tchebychev-TNP}
Les assertions suivantes sont équivalentes :
\begin{enumerate}[label=(\roman*)]
  \item $\pi(x) \sim x/\ln x$ \quad (théorème des nombres premiers) ;
  \item $\theta(x) \sim x$ ;
  \item $\psi(x) \sim x$.
\end{enumerate}
\end{proposition}

\begin{proof}[Esquisse]
L'équivalence (ii)$\iff$(iii) résulte de l'estimation $\psi(x) = \theta(x) + O(\sqrt{x})$,
car les termes avec $k \geq 2$ dans $\psi$ contribuent au plus $\sum_{p \leq \sqrt{x}} \ln p
\leq \sqrt{x} \ln x$.

Pour (i)$\iff$(ii), on utilise la sommation d'Abel\index{sommation d'Abel} :
$\theta(x) = \pi(x) \ln x - \int_2^x \pi(t)/t \, dt$, et on montre que si l'un
des deux équivalents vaut, l'autre aussi.
\end{proof}

\begin{remarque}
\label{rem:tchebychev-encadrement}
Avant la preuve du TNP, Tchebychev\index{Tchebychev, Pafnouti} (1852) avait établi
l'encadrement
\[
  0{,}921 \cdot \frac{x}{\ln x} < \pi(x) < 1{,}106 \cdot \frac{x}{\ln x}
\]
pour $x$ assez grand, prouvant ainsi que \emph{si} $\pi(x)/(x/\ln x)$ admet une
limite, celle-ci vaut nécessairement $1$.
\end{remarque}

\index{fonctions de Tchebychev|)}

\section{L'hypothèse de Riemann}
\label{sec:RH}
\index{hypothèse de Riemann|(}

\begin{theoreme}[Prolongement analytique de $\zeta$]
\label{thm:prolongement-zeta}
La fonction $\zeta(s)$ admet un prolongement méromorphe à tout le plan complexe
$\C$, avec un unique pôle simple en $s = 1$ de résidu $1$. Elle satisfait
l'équation fonctionnelle
\[
  \pi^{-s/2}\, \Gamma\!\left(\frac{s}{2}\right) \zeta(s)
  = \pi^{-(1-s)/2}\, \Gamma\!\left(\frac{1-s}{2}\right) \zeta(1-s).
\]
\end{theoreme}

\begin{definition}[Bande critique et zéros]
\label{def:bande-critique}
\index{bande critique}
\begin{itemize}
  \item Les \emph{zéros triviaux}\index{zéros triviaux} de $\zeta$ sont les
    entiers négatifs pairs : $s = -2, -4, -6, \ldots$
  \item Les \emph{zéros non triviaux}\index{zéros non triviaux} se trouvent tous
    dans la \emph{bande critique} $0 < \operatorname{Re}(s) < 1$.
  \item La \emph{droite critique}\index{droite critique} est $\operatorname{Re}(s) = 1/2$.
\end{itemize}
\end{definition}

\begin{theoreme}[Hypothèse de Riemann (1859, non démontrée)]
\label{thm:RH}
\index{hypothèse de Riemann}
Tous les zéros non triviaux de la fonction $\zeta$ de Riemann se trouvent sur la
droite critique $\operatorname{Re}(s) = 1/2$.
\end{theoreme}

\begin{remarque}[Importance et statut]
\label{rem:RH-statut}
L'hypothèse de Riemann est considérée comme l'un des problèmes ouverts les plus
importants des mathématiques. Elle fait partie des sept problèmes du millénaire
du \emph{Clay Mathematics Institute}\index{Clay Mathematics Institute},
assortis d'un prix d'un million de dollars.

Parmi ses conséquences :
\begin{itemize}
  \item Un terme d'erreur optimal dans le TNP :
    $\pi(x) = \operatorname{Li}(x) + O(\sqrt{x} \ln x)$.
  \item Des estimations fines sur les écarts entre nombres premiers consécutifs.
  \item Des résultats sur la distribution des nombres premiers dans les
    progressions arithmétiques.
\end{itemize}

À ce jour (2026), l'hypothèse de Riemann reste non démontrée, bien que des
vérifications numériques aient confirmé qu'elle est vraie pour les premiers
$10^{13}$ zéros non triviaux.
\end{remarque}

% --- Figure TikZ : bande critique et zéros de ζ ---
\begin{figure}[ht]
\centering
\begin{tikzpicture}[scale=0.75]
  % Axes
  \draw[->] (-4,0) -- (4.5,0) node[right] {$\operatorname{Re}(s)$};
  \draw[->] (0,-4.5) -- (0,5) node[above] {$\operatorname{Im}(s)$};

  % Critical strip
  \fill[yellow!15] (0,-4.5) rectangle (2,5);
  \node[font=\footnotesize, yellow!50!black] at (1,4.5) {bande critique};

  % Critical line
  \draw[thick, red, dashed] (1,-4.5) -- (1,5);
  \node[font=\footnotesize, red!70!black, anchor=west] at (1.15,4.2) {$\operatorname{Re}(s) = \tfrac{1}{2}$};

  % Pole at s=1
  \node[font=\small] at (2,0) {$\times$};
  \node[font=\footnotesize, anchor=south west] at (2.1,0.1) {pôle $s=1$};

  % Trivial zeros
  \foreach \y in {-2,-4,-6} {
    \fill[blue] (\y/2,0) circle (2.5pt);
  }
  \node[font=\footnotesize, blue, anchor=north] at (-1,-0.2) {$-2$};
  \node[font=\footnotesize, blue, anchor=north] at (-2,-0.2) {$-4$};
  \node[font=\footnotesize, blue, anchor=north] at (-3,-0.2) {$-6$};

  % Non-trivial zeros (on critical line)
  \foreach \y in {-3.73, -2.51, -1.41, 1.41, 2.51, 3.73} {
    \fill[red] (1,\y) circle (3pt);
  }
  \node[font=\footnotesize, anchor=west] at (1.3,1.41) {$\tfrac{1}{2} + 14{,}13\ldots\, i$};
  \node[font=\footnotesize, anchor=west] at (1.3,2.51) {$\tfrac{1}{2} + 21{,}02\ldots\, i$};

  % Scale labels
  \foreach \x in {-3,-2,-1,1,2} {
    \draw (\x,0.08) -- (\x,-0.08);
  }
  \node[font=\footnotesize, anchor=north] at (0,-0.15) {$0$};
  \node[font=\footnotesize, anchor=north] at (2,-0.15) {$1$};
\end{tikzpicture}
\caption{La bande critique $0 < \operatorname{Re}(s) < 1$, la droite critique
$\operatorname{Re}(s) = 1/2$ (tirets rouges), les zéros triviaux (points bleus)
et les premiers zéros non triviaux (points rouges).}
\label{fig:zeros-zeta}
\end{figure}

\index{hypothèse de Riemann|)}

\section{Problèmes ouverts célèbres}
\label{sec:problemes-ouverts}
\index{problèmes ouverts|(}

Nous terminons ce cours par un bref panorama de quelques-unes des conjectures
les plus célèbres de la théorie des nombres, toujours non résolues.

\subsection{Conjecture de Goldbach}
\label{subsec:goldbach}
\index{conjecture de Goldbach}

\begin{theoreme}[Conjecture de Goldbach (1742, non démontrée)]
\label{conj:goldbach}
Tout entier pair $n \geq 4$ est somme de deux nombres premiers.
\end{theoreme}

\begin{remarque}
\label{rem:goldbach-progres}
Parmi les résultats partiels :
\begin{itemize}
  \item Vinogradov\index{Vinogradov, Ivan} (1937) a montré que tout entier
    \emph{impair} suffisamment grand est somme de trois premiers
    (<<~conjecture faible de Goldbach~>>).
  \item Helfgott\index{Helfgott, Harald} (2013) a rendu le résultat de Vinogradov
    effectif pour tout entier impair $> 5$.
  \item Chen\index{Chen, Jingrun} (1966) a montré que tout entier pair assez
    grand s'écrit $p + q$ où $p$ est premier et $q$ est premier ou semi-premier.
  \item La conjecture a été vérifiée numériquement jusqu'à $4 \times 10^{18}$.
\end{itemize}
\end{remarque}

\subsection{Conjecture des nombres premiers jumeaux}
\label{subsec:premiers-jumeaux}
\index{nombres premiers jumeaux}

\begin{theoreme}[Conjecture des premiers jumeaux (non démontrée)]
\label{conj:jumeaux}
Il existe une infinité de nombres premiers $p$ tels que $p + 2$ soit aussi premier.
\end{theoreme}

\begin{remarque}
\label{rem:jumeaux-progres}
\begin{itemize}
  \item Brun\index{Brun, Viggo} (1919) a montré que $\sum_{p,\,p+2\text{ premiers}} 1/p$
    converge (la \emph{constante de Brun}\index{constante de Brun} $\approx 1{,}902$),
    même si la conjecture est vraie.
  \item Zhang\index{Zhang, Yitang} (2013) a prouvé un résultat spectaculaire :
    il existe une infinité de paires de premiers dont l'écart est $\leq 70\,000\,000$.
  \item Le projet Polymath\index{Polymath} a réduit cette borne à $246$ (2014).
    La conjecture des premiers jumeaux correspond à la borne $2$.
\end{itemize}
\end{remarque}

\index{problèmes ouverts|)}

\section{Résumé du chapitre}
\label{sec:resume-ch10}

\begin{center}
\begin{tabular}{lp{9cm}}
\toprule
\textbf{Concept} & \textbf{Description} \\
\midrule
$\zeta(s) = \sum n^{-s}$ & Fonction zêta de Riemann, convergente pour $\operatorname{Re}(s) > 1$ \\
Produit eulérien & $\zeta(s) = \prod_p (1-p^{-s})^{-1}$ : lien zêta--premiers \\
$L(s,\chi)$ & $L$-fonctions de Dirichlet, généralisations de $\zeta$ \\
Dirichlet & Infinité de premiers dans toute progression $a + nq$, $\pgcd(a,q)=1$ \\
TNP & $\pi(x) \sim x/\ln x$ (Hadamard, de la Vallée-Poussin, 1896) \\
$\theta(x)$, $\psi(x)$ & Fonctions de Tchebychev, TNP $\iff \psi(x) \sim x$ \\
Hypothèse de Riemann & Zéros non triviaux sur $\operatorname{Re}(s) = 1/2$ (ouverte) \\
\bottomrule
\end{tabular}
\end{center}

\section{Exercices}
\label{sec:exercices-ch10}

\begin{exercice}[$\star$ --- Produit eulérien]
\label{exo:produit-eulerien}
\begin{enumerate}[label=(\alph*)]
  \item En utilisant le produit eulérien, montrer que $\zeta(s) \neq 0$ pour
    $\operatorname{Re}(s) > 1$.
  \item Calculer $\prod_{p} (1 - p^{-2})^{-1}$ et retrouver ainsi $\zeta(2) = \pi^2/6$.
  \item En déduire que $\sum_{n=1}^{\infty} \mu(n)/n^2 = 6/\pi^2$, où $\mu$ est
    la fonction de Möbius\index{fonction de Möbius}.
\end{enumerate}
\end{exercice}

\begin{exercice}[$\star$ --- Fonction de comptage]
\label{exo:fonction-comptage}
\begin{enumerate}[label=(\alph*)]
  \item Calculer $\pi(100)$, $\pi(200)$ et $\pi(1000)$.
  \item Comparer avec $x/\ln x$ et $\operatorname{Li}(x)$ pour ces valeurs.
  \item Pour quels $x \leq 1000$ l'écart $|\pi(x) - x/\ln x|$ est-il maximal ?
\end{enumerate}
\end{exercice}

\begin{exercice}[$\star\star$ --- Identités de la fonction zêta]
\label{exo:identites-zeta}
\begin{enumerate}[label=(\alph*)]
  \item Montrer que $\zeta(s)^2 = \sum_{n=1}^{\infty} d(n)/n^s$ pour $\operatorname{Re}(s) > 1$,
    où $d(n)$ est le nombre de diviseurs de $n$.
  \item Montrer que $\zeta(s)/\zeta(2s) = \sum_{n=1}^{\infty} \abs{\mu(n)}/n^s$.
  \item En déduire que la densité des entiers sans facteur carré est $6/\pi^2$.
\end{enumerate}
\end{exercice}

\begin{exercice}[$\star\star$ --- Fonctions de Tchebychev]
\label{exo:tchebychev}
\begin{enumerate}[label=(\alph*)]
  \item Calculer $\theta(30)$ et $\psi(30)$.
  \item Montrer que $\psi(x) = \theta(x) + \theta(x^{1/2}) + \theta(x^{1/3}) + \cdots$
  \item En déduire que $\psi(x) - \theta(x) = O(\sqrt{x} \ln x)$.
\end{enumerate}
\end{exercice}

\begin{exercice}[$\star\star\star$ --- Divergence de $\sum 1/p$]
\label{exo:divergence-somme-premiers}
\index{somme des inverses des premiers}
En utilisant le produit eulérien et le fait que $\zeta(s) \to +\infty$ quand
$s \to 1^+$, montrer que
\[
  \sum_{p \text{ premier}} \frac{1}{p} = +\infty.
\]
\textit{Indication :} prendre le logarithme du produit eulérien et estimer
$-\ln(1-p^{-s})$.
\end{exercice}

\begin{exercice}[$\star\star\star$ --- Formule de Mertens]
\label{exo:mertens}
\index{formule de Mertens}
Démontrer que
\[
  \sum_{\substack{p \leq x \\ p \text{ premier}}} \frac{1}{p}
  = \ln \ln x + M + O\!\left(\frac{1}{\ln x}\right),
\]
où $M \approx 0{,}2615$ est la constante de Meissel--Mertens\index{constante de Meissel--Mertens}.
\end{exercice}

\begin{exercice}[$\star\star\star$ --- Hypothèse de Riemann et terme d'erreur]
\label{exo:RH-erreur}
Admettre l'hypothèse de Riemann.
\begin{enumerate}[label=(\alph*)]
  \item Montrer que $\psi(x) = x + O(\sqrt{x}\, (\ln x)^2)$.
  \item En déduire que $\pi(x) = \operatorname{Li}(x) + O(\sqrt{x} \ln x)$.
  \item Comparer avec le meilleur terme d'erreur connu sans l'HR :
    $\pi(x) = \operatorname{Li}(x) + O(x \exp(-c\sqrt{\ln x}))$.
\end{enumerate}
\end{exercice}

\index{théorie analytique des nombres|)}

%%%%%%%%%%%%%%%%%%%%%%%%%%%%%%%%%%%%%%%%%%%%%%%%%%%%%%%%%%%%%%%%%%%%
%%  BIBLIOGRAPHIE
%%%%%%%%%%%%%%%%%%%%%%%%%%%%%%%%%%%%%%%%%%%%%%%%%%%%%%%%%%%%%%%%%%%%

\begin{thebibliography}{99}

\bibitem{HardyWright}
\textsc{Hardy}, G.H. et \textsc{Wright}, E.M.,
\emph{An Introduction to the Theory of Numbers},
6\ieme{} édition,
Oxford University Press, 2008.

\bibitem{Serre}
\textsc{Serre}, J.-P.,
\emph{Cours d'arithmétique},
3\ieme{} édition,
Presses Universitaires de France, 1988.

\bibitem{IrelandRosen}
\textsc{Ireland}, K. et \textsc{Rosen}, M.,
\emph{A Classical Introduction to Modern Number Theory},
2\ieme{} édition,
Graduate Texts in Mathematics~84,
Springer, 1990.

\bibitem{Apostol}
\textsc{Apostol}, T.M.,
\emph{Introduction to Analytic Number Theory},
Undergraduate Texts in Mathematics,
Springer, 1976.

\bibitem{NivenZuckermanMontgomery}
\textsc{Niven}, I., \textsc{Zuckerman}, H.S. et \textsc{Montgomery}, H.L.,
\emph{An Introduction to the Theory of Numbers},
5\ieme{} édition,
Wiley, 1991.

\bibitem{Nathanson}
\textsc{Nathanson}, M.B.,
\emph{Elementary Methods in Number Theory},
Graduate Texts in Mathematics~195,
Springer, 2000.

\bibitem{Koblitz}
\textsc{Koblitz}, N.,
\emph{$p$-adic Numbers, $p$-adic Analysis, and Zeta-Functions},
2\ieme{} édition,
Graduate Texts in Mathematics~58,
Springer, 1984.

\bibitem{Davenport}
\textsc{Davenport}, H.,
\emph{Multiplicative Number Theory},
3\ieme{} édition, révisée par H.L. Montgomery,
Graduate Texts in Mathematics~74,
Springer, 2000.

\end{thebibliography}

%%%%%%%%%%%%%%%%%%%%%%%%%%%%%%%%%%%%%%%%%%%%%%%%%%%%%%%%%%%%%%%%%%%%
\printindex
\end{document}
